
% !LW recipe=2xpdfLaTeX
% !LW recipe=pdfLaTeX ➞ bibTeX ➞ pdfLaTeX

\documentclass[a4paper]{article}

% Shared packages, common definitions and external TikZ-figure machinery

%\begingroup For image generation 
  \usepackage[strict]{changepage}

  \usepackage{shellesc}

  \usepackage{etoolbox}
  \newbool{forcetikzrecompile}

  \usepackage{tikz}
  \usetikzlibrary{
    positioning,
    shapes.geometric,
    plotmarks,
    arrows.meta,
    calc,
    decorations.pathreplacing,
    backgrounds,
    patterns.meta,
    fpu, % «floating point unit»
    % external,
  } 

  \usepackage{pgfplots}
  \pgfplotsset{compat=newest}
  \usepgfplotslibrary{
    patchplots,
    groupplots,
    colorbrewer
  }
  \pgfdeclarelayer{background}
  \pgfdeclarelayer{foreground}
  \pgfsetlayers{background,main,foreground}
%\endgroup

%\begingroup For the paper
  \usepackage{microtype}

  \usepackage{amsmath}
  \usepackage{amsthm}
  \usepackage{accents}

  \usepackage{zref-savepos}
  \newcounter{figposcount}

  \usepackage{xstring}

  %\usepackage[numbers]{natbib}
  %\usepackage[authoryear]{natbib}
  % \usepackage[authoryear,longnamesfirst]{natbib}
  \usepackage[numbers,square]{natbib}

  \usepackage{
    graphicx,
    caption,
    subcaption,
    stfloats,
    float
  }

  \usepackage[%
    ruled,%
    vlined,%
    linesnumbered%
  ]{algorithm2e}

  \usepackage{ragged2e}

  \usepackage{enumitem}
  \setlist[enumerate]{%
    label=\arabic*.,
    topsep=0.5em,%
    parsep=0em,%
    itemsep=0.2em,%
    % leftmargin=2em,%
    rightmargin=1.5em,%
  }
  \setlist[itemize]{%
    label=\(\diamond\),
    topsep=0.5em,%
    parsep=0em,%
    itemsep=0.2em,%
    % leftmargin=2em,%
    rightmargin=1.5em,%
  }

  \usepackage{hyperref}
  \usepackage{xurl}

  \usepackage[english]{cleveref}

  \usepackage{siunitx}
  \sisetup{%
    exponent-mode=threshold,%
    exponent-thresholds=-3:3,%
    round-mode=places,% figures
    round-precision=3,%
    round-pad=false,%
    separate-uncertainty=true,%
    print-unity-mantissa=false,%
    tight-spacing=true,%
  }
%\endgroup


\newif\ifcasscclass
\makeatletter
\@ifclassloaded{cas-sc}{\casscclasstrue}{\casscclassfalse}
\makeatother

\newif\ifpageisodd % È usato (nelle macro delle infopagini) dalla meccanica

\ifcasscclass
%\begingroup For image generation
  %\begingroup Tipografia
    \newcommand\textsize[1]{%
      \pgfmathparse{#1*1.2}%
      \fontsize{#1pt}{\pgfmathresult pt}\selectfont%
    }

    \DeclareMathOperator\FRC{\mathcal R}

    % (il rilevamento della classe è in cima al file)

    \providecommand\Ifthispageodd[2]{\checkoddpage\ifoddpage#1\else#2\fi}
  %\endgroup
\fi

  %\begingroup TikZ
    %\begingroup Didascalie
      \newcommand\newcaption[2]{%
        \captionsetup{width=#2} % aboveskip=-.2cm,
        \caption{#1}%
      }
      
      \newcommand\newsubcaption[2][ ]{%
        \refstepcounter{subfigure}%
        (\thesubfigure)#1#2%
      }
    %\endgroup
    
\ifcasscclass
    %\begingroup RedefineTikZbounds
      % La differenza è che colla prima bisogna catturare anche le etichette perché si crea un pdf isolato mentre colla non è importante perché viene tutto compilato sul momento

      %\begingroup Versione valida con esternalizzazione [ma anche senza]
        \newcommand\redefineTikZbounds[2]{
          \coordinate (RealSW) at (current bounding box.south west);
          \coordinate (RealNE) at (current bounding box.north east);
          \coordinate (Center) at ($(#1.west)!0.5!(#2.east)$);
    
          \pgfresetboundingbox 

          \useasboundingbox 
            let \p1=(RealSW),\p2=(RealNE),\p3=(Center),
            \n1={max(\x3-\x1,\x2-\x3)} in
            (\x3-\n1,\y1) rectangle (\x3+\n1,\y2+.25em);

          % \path[use as bounding box] 
          %   let \p1=(RealSW),\p2=(RealNE),\p3=(Center),
          %   \n1={max(\x3-\x1,\x2-\x3)} in
          %   (\x3-\n1,\y1) rectangle (\x3+\n1,\y2+.25em);
          % Lo «.25em» discende dalla traslazione imposta all'etichetta della scala verticale mediante «yScaleLabelShift»
        } % Scatola delimitante simmetrica
      %\endgroup

      %\begingroup Versione valida senza esternalizzazione
        % \newcommand\redefineTikZbounds[2]{
        %   % Salva le coordinate della cima e del fondo della figura corrente (cosí da mantenere i titoli e le didascalie)
        %   \coordinate (Cima) at (current bounding box.north);
        %   \coordinate (Fondo) at (current bounding box.south);
    
        %   \pgfresetboundingbox % Resetta la scatola dei limiti
    
        %   \useasboundingbox % Definisce la nuova scatola come un rettangolo
        %     (#1.west |- Cima)  % con ascissa sinistra data da «Col1 c1r1.west», 
        %     rectangle      % ascissa destra data da «Col2 c1r1.east»
        %     (#2.east |- Fondo); % ordinate date dalla cima e fondo originali
        %     % L'operatore «|-» in «(A |- B)» prende come ascissa quella del punto A mentre come coordinata quella del punto B
        % }
      %\endgroup
    %\endgroup

    %\begingroup TikZBox
      % Parzialmente ispirato da https://gist.github.com/Caramdir/899107

      \makeatletter
      \newsavebox\tikzBufferBox
      \xdef\offsetX{0pt}
      \xdef\offsetY{0pt}
      
      % \tikzbox[#1]{#2}{#3}[#4]{#5}
      % #1 -> Ancora rispetto cui posizionare l'intera scatola (di base è «south west», ossia l'angolo in basso a sinistra)
      % #2 -> Nodo dove posizionare l'origine (per la natura di come LaTeX riceve gli argomenti è buon pratica passare questo circondato da parentesi graffe)
      % #3 -> Nome della scatola limitante definita dall'ambiente «scope»
      % #4 -> Argomenti opzionali da passare all'ambiente «scope»; di base sono vuoti
      % #5 -> Contenuto della scatola circondato dall'ambiente «scope»
      % Es. \tikzbox{{(0,0)}}{figuraTikZ}{blocco TikZ}

      % Macro di apertura che controlla se v'è l'argomento facoltativo dell'ancora
      \def\tikzbox{\pgfutil@ifnextchar[{\tikzbox@anchor}{\tikzbox@anchor[south west]}}
      % La funzione «\pgfutil@ifnextchar» serve per controllare se il prossimo carattere è l'apertura di una paretensi graffa «[»
      
      % Macro intermedia per cogliere gli argomenti obbligatori e controllare quelli successivi facoltativi
      \def\tikzbox@anchor[#1]#2#3{\pgfutil@ifnextchar[{\tikzbox@opt[#1]#2#3}{\tikzbox@opt[#1]#2#3[]}}

      % Macro degli argomenti obbligatori
      \long\def\tikzbox@opt[#1]#2#3[#4]#5{%
        \node[ % Misurazione della scatola in un nodo invisibile
          overlay,%
          opacity=0,%
          inner sep=0pt,%
          outer sep=0pt%
        ] {%
          \sbox{\tikzBufferBox}{%
            \begin{tikzpicture}%
              \begin{scope}[local bounding box=measuringTikZBox]%
                #5%
              \end{scope}%
              \pgfpointanchor{measuringTikZBox}{#1}%
              \pgfgetlastxy{\tempX}{\tempY}% 
              \xdef\offsetX{\tempX}%
              \xdef\offsetY{\tempY}%
            \end{tikzpicture}%
          }%
        };%
        %
        % Disegno reale traslato
        \begin{scope}[%
          local bounding box=#3,%
          shift={#2},%
          xshift=-\offsetX,%
          yshift=-\offsetY,%
          #4%
        ]%
          #5%
        \end{scope}%
      }%
      \makeatother
    %\endgroup

    %\begingroup Chiavi-Valori
      %\begingroup Testo
        % La dimensione del testo di base è 11pt
        \newcommand\titleSize{11}
        \newcommand\subCaptionSize{10}
        \newcommand\ParameterSize{9}
        \newcommand\labelSize{10}
        \newcommand\tickSize{9}
        \newcommand\legendSize{9}
      %\endgroup
    
      %\begingroup Infopagine
        %\begingroup Chiavi condivise
        \newcounter{lpc} % «loopcounter»
        \newcommand\infopageSharedKeys{
          %\begingroup Testo
            \renewcommand\titleSize{10}
            \renewcommand\subCaptionSize8
            \renewcommand\ParameterSize7
            \renewcommand\labelSize9
            \renewcommand\tickSize6
            \renewcommand\legendSize7
          %\endgroup
    
          %\begingroup Dimensioni
            \pgfmathsetmacro\plotWidth7
            \pgfmathsetmacro\plotHeight{.618*\plotWidth}
            \pgfmathsetmacro\halfPlotHeight{\plotHeight/2}
            %
            \pgfmathparse{0.825*\plotWidth}
            \pgfmathsetmacro\plotWidth\pgfmathresult
            \pgfmathsetmacro\halfWidth{\plotWidth/2}
            \pgfmathsetmacro\plotLastWidth{1.25*\plotWidth}
    
            \pgfmathsetmacro\xPlotSep{1}
            \pgfmathsetmacro\yPlotSep{1.25}
            \pgfmathsetmacro\yGroupSep{2.75}
          %\endgroup
    
          %\begingroup Legende
            \newcommand\placeLegend[5]{%
              /pgfplots/plotLegend,
              at={([shift={(##4,##5)}]axis description cs:##1,##2)},
              anchor=##3,
            }
            %
            \pgfplotsset{%
              plotLegend/.style={%
                font=\textsize{\legendSize},
                /pgfplots/legend cell align=left,
                % align=left,
                % draw=white!15!black,
                draw=none, % Bordo invisibile
                fill=white,
                fill opacity=0.5,
                text opacity=1,
                % text width=2cm,
              },
              plotLegendNW/.style={legend style={\placeLegend{0}{1}{north west}{0em}{0em}}},
              plotLegendNE/.style={legend style={\placeLegend{1}{1}{north east}{0em}{0em}}},
              plotLegendSW/.style={legend style={\placeLegend{0}{0}{south west}{0em}{0em}}},
              plotLegendSE/.style={legend style={\placeLegend{1}{0}{south east}{0em}{0em}}},
              %
              % plotLegendW/.style={legend style={plotLegend,at={(0,.5)},anchor=west}},
              %
              % plotLegendNW/.style={legend style={plotLegend,legend pos=north west}},
              % plotLegendNE/.style={legend style={plotLegend,legend pos=north east}},
              % plotLegendSW/.style={legend style={plotLegend,legend pos=south west}},
              % plotLegendSE/.style={legend style={plotLegend,legend pos=south east}},
              %
              % plotLegendNWcloser/.style={legend style={plotLegend,at={(0,1)},anchor=north west}},
              % plotLegendNEcloser/.style={legend style={plotLegend,at={(1,1)},anchor=north east}},
              % plotLegendSWcloser/.style={legend style={plotLegend,at={(0,0)},anchor=south west}},
              % plotLegendSEcloser/.style={legend style={plotLegend,at={(0,1)},anchor=south east}},
            }
            %
            \pgfmathsetmacro\lineLegendLength{1}
            \pgfmathsetmacro\areaLegendHeigth{.25}
            \pgfplotsset{
              legend image code/.code={
                \draw [mark repeat=2,mark phase=2,##1]
                  plot coordinates {
                    (0em,0em)
                    (.5*\lineLegendLength em,0em)
                    (\lineLegendLength em,0em)
                  };
              },
              /pgfplots/area legend/.style={
                legend image code/.code={
                  \draw [##1,fill,draw=none] 
                    (0em,-\areaLegendHeigth em) 
                    rectangle 
                    (\lineLegendLength em,\areaLegendHeigth em);
                },
              },
            }
            %
            \newcommand\ExactHistogramLegend{Exact histogram}
            \newcommand\ApproxHistogramLegend{Approx. histogram}
            \newcommand\RealHistogramLegend{Real histogram}
            %
            \newcommand\ExactBLNFitLegend{Exact BLN fit}
            \newcommand\ApproxBLNFitLegend{Approx BLN fit}
            \newcommand\RealBLNFitLegend{Real BLN fit}
            %
            \newcommand\ExactEmpiricCCDFLegend{Exact empiric CCDF}
            \newcommand\ApproxEmpiricCCDFLegend{Approx. empiric CCDF}
            \newcommand\RealEmpiricCCDFLegend{Real empiric CCDF}
            %
            \newcommand\ExactParetoFitLegend{Exact Pareto fit}
            \newcommand\ApproxParetoFitLegend{Approx. Pareto fit}
            \newcommand\RealParetoFitLegend{Real Pareto fit}
            %
            \newcommand\ExactScatterLegend{Exact scatter}
            \newcommand\ApproxScatterLegend{Approx. scatter}
            \newcommand\RealScatterLegend{Real scatter}
            %
            \newcommand\ClasskILegend{Class k=283}
            \newcommand\ClasskIILegend{Class k=83}
            \newcommand\ClasskIIILegend{Class k=43}
            \newcommand\ClasskIVLegend{Class k=10}
          %\endgroup
    
          %\begingroup Etichette vuote
            \pgfplotsset{
              yLabelEmpty/.style={ylabel={}},
              yTickLabelsEmpty/.style={yticklabels={}},
              yTickScaleLabelEmpty/.style={ytick scale label code/.code={}},
              xTickScaleLabelEmpty/.style={xtick scale label code/.code={}},
              emptyLeged/.style={
                legend style={draw=none,fill=none,nodes={draw=none}},
                legend image post style={opacity=0},
              },
            }
          %\endgroup
    
          %\begingroup Estetica
            \newcommand\smallminus{\scalebox{0.5}[1]{$-$}\hspace{-.11em}}
            \pgfplotsset{
              /pgfplots/log number format basis/.code 2 args={%
                \pgfmathsetmacro\exponent{##2}%
                \pgfmathparse{sign(\exponent)}%
                \ifnum\pgfmathresult<0%
                  \pgfmathparse{abs(\exponent)}%
                  $##1^{\smash{\smallminus\pgfmathprintnumber{\pgfmathresult}}}$%
                \else%
                  $##1^{\smash{\pgfmathprintnumber{\exponent}}}$%
                \fi%
              },
              ytick scale label code/.code={%
                \pgfmathsetmacro\exponent{##1}%
                \pgfmathparse{sign(\exponent)}%
                \ifnum\pgfmathresult<0
                  \pgfmathparse{abs(\exponent)}%
                  $\cdot 10^{\smash{\smallminus\pgfmathprintnumber{\pgfmathresult}}}$%
                \else
                  $\cdot 10^{\smash{\pgfmathprintnumber{\exponent}}}$%
                \fi
              },
            }
          %\endgroup
    
          %\begingroup Colori
            \definecolor{blueish}{rgb}{0.00000,0.44700,0.74100}%
            \definecolor{reddish}{rgb}{0.83529,0.36863,0.00000}%
          %\endgroup
    
          %\begingroup Commenti
            %\begingroup Impostazioni
            \tikzset{
              title/.style={
                align=center,
                anchor=south,
                font=\bfseries\textsize{\titleSize},
              },
              fittedParameters/.style={
                align=center,
                font=\textsize{\ParameterSize},
              },
              varyingParameter/.style={
                align=center,
                % rotate=90,
                font=\textsize{\subCaptionSize},
              },
              subCaption/.style={
                % text width=\plotWidth cm,
                align=center,
                font=\textsize{\subCaptionSize},
              },
            }
            \newcommand\codetocolor[1]{\ifcase##1blueish\or reddish\else black\fi}
            %\endgroup
    
            %\begingroup Parametri Sardegna
            \newcommand\betaReal{\num{1.273}}
            \newcommand\zetaReal{\num{0.929}}
            \newcommand\mIReal{\num{3.227}}
            \newcommand\sIReal{\num{0.408}}
            \newcommand\mIIReal{\num{4.099}}
            \newcommand\sIIReal{\num{0.556}}
            %\endgroup
    
            %\begingroup Parametri Pareto
            \newcommand\codetoParetovalue[1]{
              \ifcase##1\betaExact\or\betaAprx\else\betaReal\fi
            }
            \newcommand\writeParetoParameters[7][center]{
              \node[fittedParameters,anchor=##1] at ([shift={(##5,##6)}]##4) {%
                \(\begin{aligned}
                  \color{\codetocolor{##2}}\beta
                  &\color{\codetocolor{##2}}=\codetoParetovalue{##2}\\[##7]
                  \color{\codetocolor{##3}}\beta
                  &\color{\codetocolor{##3}}=\codetoParetovalue{##3}
                \end{aligned}\)
              };
            }
            %\endgroup

            %\begingroup % Vecchia versione in cui si scrivevano esplicitamente tutte le 9 combinazioni possibili
            % \newcommand\writeParameters[6]{%
            %   \node[parameter] at ([shift={(##4,##5)}]##1) {%
            %     \ifcase##2% Caso 0
            %       \ifcase##3% Caso 0
            %         \(\begin{aligned}
            %           \color{blueish}\beta&\color{blueish}=\betaExact\\[##6]
            %           \color{blueish}\beta&\color{blueish}=\betaExact
            %         \end{aligned}\)
            %       \or% Caso 1
            %         \(\begin{aligned}
            %           \color{blueish}\beta&\color{blueish}=\betaExact\\[##6]
            %           \color{reddish}\beta&\color{reddish}=\betaAprx% Caso 1
            %         \end{aligned}\)
            %       \else% Caso 2
            %         \(\begin{aligned}
            %           \color{blueish}\beta&\color{blueish}=\betaExact\\[##6]
            %           \color{black}\beta&\color{black}=\betaReal% Caso 2
            %         \end{aligned}\)
            %       \fi
            %     \or% Caso 1
            %       \ifcase##3% Caso 0
            %         \(\begin{aligned}
            %           \color{reddish}\beta&\color{reddish}=\betaAprx\\[##6]
            %           \color{blueish}\beta&\color{blueish}=\betaExact
            %         \end{aligned}\)
            %       \or% Caso 1
            %         \(\begin{aligned}
            %           \color{reddish}\beta&\color{reddish}=\betaAprx\\[##6]
            %           \color{reddish}\beta&\color{reddish}=\betaAprx
            %         \end{aligned}\)
            %       \else% Caso 2
            %         \(\begin{aligned}
            %           \color{reddish}\beta&\color{reddish}=\betaAprx\\[##6]
            %           \color{black}\beta&\color{black}=\betaReal
            %         \end{aligned}\)
            %       \fi
            %     \else% Caso 2
            %       \ifcase##3% Caso 0
            %         \(\begin{aligned}
            %           \color{black}\beta&\color{black}=\betaReal\\[##6]
            %           \color{blueish}\beta&\color{blueish}=\betaExact
            %         \end{aligned}\)
            %       \or% Caso 1
            %         \(\begin{aligned}
            %           \color{black}\beta&\color{black}=\betaReal\\[##6]
            %           \color{reddish}\beta&\color{reddish}=\betaAprx
            %         \end{aligned}\)
            %       \else% Caso 2
            %         \(\begin{aligned}
            %           \color{black}\beta&\color{black}=\betaReal\\[##6]
            %           \color{black}\beta&\color{black}=\betaReal
            %         \end{aligned}\)
            %       \fi
            %     \fi
            %   };
            % }
            %\endgroup
    
            %\begingroup Parametri bilognormale
            \newcommand\codetoparameter[1]{
              \ifcase##1\zeta  % Caso 0
              \or\mu_1     % Caso 1
              \or\sigma_1   % Caso 2
              \or\mu_2     % Caso 3
              \else\sigma_2\fi % Caso 4
            }
            \newcommand\codetoBiLognormalvalue[2]{
              \ifcase##1 % Caso 0
                \ifcase##2\zetaExact % Caso 0
                \or\mIExact    % Caso 1
                \or\sIExact    % Caso 2
                \or\mIIExact    % Caso 3
                \else\sIIExact\fi % Caso 4
              \or % Caso 1
                \ifcase##2\zetaAprx % Caso 0
                \or\mIAprx     % Caso 1
                \or\sIAprx     % Caso 2
                \or\mIIAprx    % Caso 3
                \else\sIIAprx\fi  % Caso 4
              \else % Caso 2
                \ifcase##2\zetaReal % Caso 0
                \or\mIReal     % Caso 1
                \or\sIReal     % Caso 2
                \or\mIIReal    % Caso 3
                \else\sIIReal\fi  % Caso 4
              \fi
            }
            \newcommand\matrixline[2]{\codetoparameter{##2}&=\codetoBiLognormalvalue{##1}{##2}}
            \def\lineloop##1##2{%
              \matrixline{##1}{\number\value{lpc}}\stepcounter{lpc}%
              \ifnum\number\value{lpc}<5\\[##2]\lineloop{##1}{##2}\fi%
            } % https://tex.stackexchange.com/questions/491358/tikz-foreach-to-work-with-the-aligned-environment
            \newcommand\writeBiLognormalParameters[6][center]{
              \node[fittedParameters,anchor=##1] at ([shift={(##4,##5)}]##3) {%
                \setcounter{lpc}{0}
                \(\color{\codetocolor{##2}}
                \begin{aligned}
                  \lineloop{##2}{##6}
                \end{aligned}\)
              };
            }
            %\endgroup
    
            %\begingroup Altre etichette
            \newcommand\writeSizeEvolutionDescriptions[1][5]{%
              \foreach \col/\clr/\txt in {%
                1/blueish/Esatto,
                2/reddish/Approssimato
              }{
                \node[
                  fittedParameters,
                  anchor=south,
                  yshift=.1em,
                ] at (Row##1 c\col r1.south) {\color{\clr}\txt};
              }
            }
    
            \newcommand\writeVaryingParameter[5]{%
              \foreach \r in {1,...,5}{
                \node[
                  varyingParameter,
                  anchor=##3,
                ] at ([shift={(##4,##5)}]Row\r\space c##1r1.##2)
                {\(\varyingParameterSymbol=\varyingParameterList\r\)};
              }
            }
            %\endgroup
          %\endgroup
    
          %\begingroup Sottodidascalie
            \newcommand\writeSubCaption[5]{%
              \node[subCaption,anchor=##3]
              at ([shift={(##4,##5)}]##1.##2)
              {\newsubcaption[]{}};
            }
            \newcommand\writeSubCaptions[5][3]{%
              \foreach \r in {1,...,5}{
                \foreach \c in {1,...,##1}{
                  \writeSubCaption{Row\r\space c\c r1}{##2}{##3}{##4}{##5}
                }
              }
            }
          %\endgroup
        }
        %\endgroup
    
        %\begingroup Chiavi del primo tipo d'infopagina
        \newcommand\infopageKeysI{
          \infopageSharedKeys
    
          %\begingroup Gruppi
            \ifpageisodd
              \pgfplotsset{/pgfplots/group/yDescPos/.style={y descriptions at=edge right}}
            \else
              \pgfplotsset{/pgfplots/group/yDescPos/.style={y descriptions at=edge left}}
            \fi
            \pgfplotsset{
              /pgfplots/group/plotGroup/.style={
                horizontal sep=\xPlotSep em,
                vertical sep=\yPlotSep em,
                yDescPos,
                group size=3 by 1,
              },
              /pgfplots/group/plotLastGroup/.style={
                plotGroup,
                x descriptions at=edge bottom,
              },
            }
          %\endgroup
    
          %\begingroup Stili
            %\begingroup Selezione «recto»/«verso»
            \ifpageisodd
              \pgfplotsset{%
                yLineAxis/.style={axis y line*=right},
                /tikz/yLabelShift/.style={%
                  yshift=-1.25em,
                  at={(axis description cs:1,0.5)},
                  anchor=north,
                  % rotate=180,
                },
                /tikz/yScaleLabelShift/.style={%
                  yshift=.25em,
                  at={(1,1)},
                  anchor=south east,
                },
                /tikz/yTickLabelShift/.style={xshift=-0.15em},
              }
            \else
              \pgfplotsset{%
                yLineAxis/.style={axis y line*=left},
                /tikz/yLabelShift/.style={%
                  yshift=1.25em,
                  at={(axis description cs:0,0.5)},
                  anchor=south,
                  % rotate=180,
                },
                /tikz/yScaleLabelShift/.style={%
                  yshift=.25em,
                  at={(0,1)},
                  anchor=south west,
                },
                /tikz/yTickLabelShift/.style={xshift=+0.15em},
              }
            \fi
            %\endgroup
            %
            %\begingroup Grafici
            \pgfplotsset{
              plotRow/.style={
                scale only axis,
                width=\plotWidth cm,height=\plotHeight cm,
                xlabel style={font=\color{white!15!black}\textsize{\labelSize}},
                ylabel style={font=\color{white!15!black}\textsize{\labelSize},yLabelShift},
                xticklabel style={font=\textsize{\tickSize}},
                yticklabel style={font=\textsize{\tickSize},yTickLabelShift},
                y tick scale label style={yScaleLabelShift},
                % axis x line*=bottom,yLineAxis,
                enlarge x limits=0.1,enlarge y limits=0.1,
                xmajorgrids,ymajorgrids,
                xminorticks=true,yminorticks=true,
                scale ticks above exponent=2,
                axis background/.style={fill=white},
                grid style={dashed}
              },
              % xlabel123/.style={xlabel={\(h\)}},
              xlabel3/.style={xlabel={\(h\)}},
              ylabel23/.style={ylabel={\(\FRC(h)\)}},
              row1Keys/.style={
                % ylabel={\(\underhat{\bar l}_N(h)\)},
                ylabel={\(\bar l_N(h)\)},
                % enlarge x limits={lower,value=.05},
                enlarge x limits=.05,
                enlarge y limits={upper,value=.05},
                xmode=linear,
                log origin=infty,
              },
              row2KeysWithLabels/.style={
                ylabel23,
                xmode=linear,ymode=log
              },
              row2Keys/.style={
                ylabel23,
                xmode=linear,ymode=log,
                xticklabels={},
              },
              row3Keys/.style={
                xlabel3,ylabel23,
                xmode=linear,ymode=log
              },
              row4Keys/.style={
                xlabel={\(\xi\)},ylabel={\(s(\xi)\)},
                xmode=log,ymode=log
              },
              row5KeysLin/.style={
                width=\plotLastWidth cm,
                xlabel={\(t\)},ylabel={\(s(t)\)},
                enlarge x limits=0,
                enlarge y limits=0.05,
                x tick scale label style={
                  yshift=-1em,
                  at={(.95,0)},
                  anchor=north,
                },
              },
              row5Keys/.style={
                row5KeysLin,
                ymode=log,
              },
            }
            %\endgroup
          %\endgroup
        }
        %\endgroup
    
        %\begingroup Chiavi del secondo tipo d'infopagina
        \newcommand\infopageKeysII{
          \infopageSharedKeys
    
          %\begingroup Gruppi
            \ifpageisodd
              \pgfplotsset{/pgfplots/group/yDescPos/.style={y descriptions at=edge right}}
            \else
              \pgfplotsset{/pgfplots/group/yDescPos/.style={y descriptions at=edge left}}
            \fi
            \pgfplotsset{
              /pgfplots/group/plotGroup/.style={
                horizontal sep=\xPlotSep em,
                vertical sep=\yPlotSep em,
                y descriptions at=all,
                group size=3 by 1,
              },
              /pgfplots/group/plotLastGroup/.style={
                plotGroup,
                x descriptions at=edge bottom,
              },
            }
          %\endgroup
    
          %\begingroup Stili
            %\begingroup Selezione «recto»/«verso»
            \ifpageisodd
              \pgfplotsset{
                yLabel/.style={
                  ylabel style={
                    font=\color{white!15!black}\textsize{\labelSize},
                    yshift=-1.25em,
                    at={(axis description cs:1,0.5)},
                    anchor=north,
                    % rotate=180,
                  },
                },
                yLabelInv/.style={yLabelEmpty},
                yScaleLabelShift/.style={
                  y tick scale label style={
                    yshift=.25em,
                    % xshift=1em,
                    at={(1,1)},
                    anchor=south east,
                  },
                },
                yScaleLabelShiftInv/.style={
                  y tick scale label style={
                    yshift=.25em,
                    % xshift=-1em,
                    at={(0,1)},
                    anchor=south west,
                  },
                },
                % yLineAxis/.style={axis y line*=right},
              }
            \else
              \pgfplotsset{
                yLabel/.style={yLabelEmpty},
                yLabelInv/.style={
                  ylabel style={
                    font=\color{white!15!black}\textsize{\labelSize},
                    yshift=1.25em,
                    at={(axis description cs:0,0.5)},
                    anchor=south,
                    % rotate=180,
                  },
                },
                yScaleLabelShift/.style={
                  y tick scale label style={
                    yshift=.25em,
                    % xshift=1em,
                    at={(1,1)},
                    anchor=south east,
                  },
                },
                yScaleLabelShiftInv/.style={
                  y tick scale label style={
                    yshift=.25em,
                    % xshift=-1em,
                    at={(0,1)},
                    anchor=south west,
                  },
                },
                % yLineAxis/.style={axis y line*=left},
              }
            \fi
            %\endgroup
            %
            %\begingroup Grafici
            \pgfplotsset{
              plotRow/.style={
                scale only axis,
                width=\plotWidth cm,height=\plotHeight cm,
                xlabel style={font=\color{white!15!black}\textsize{\labelSize}},
                % ylabel style={font=\color{white!15!black}\textsize{\labelSize},yLabelShift},
                xticklabel style={font=\textsize{\tickSize}},
                % axis x line*=bottom,yLineAxis,
                enlarge x limits=0.1,enlarge y limits=0.1,
                xmajorgrids,ymajorgrids,
                xminorticks=true,yminorticks=true,
                scale ticks above exponent=2,
                axis background/.style={fill=white},
                grid style={dashed},
              },
              plotLeftCol/.style={
                yLabelInv,
                yticklabel style={font=\textsize{\tickSize},xshift=0.15em},
                yticklabel pos=left,
                yScaleLabelShiftInv,
              },
              plotMiddleCol/.style={
                yLabelEmpty,
                yTickLabelsEmpty,
                yTickScaleLabelEmpty,
              },
              plotRightCol/.style={
                yLabel,
                yticklabel style={font=\textsize{\tickSize},xshift=-0.15em},
                yticklabel pos=right,
                yScaleLabelShift,
              },
              xlabel3/.style={xlabel={\(h\)}},
              ylabel23/.style={ylabel={\(\FRC(h)\)}},
              row1Keys/.style={
                ylabel={\(\bar l_N(h)\)},
                enlarge x limits=.05,
                enlarge y limits={upper,value=.05},
                xmode=linear,
                log origin=infty,
              },
              row2Keys/.style={
                ylabel23,
                xmode=linear,ymode=log,
                % xticklabels={},
              },
              row3Keys/.style={
                xlabel3,ylabel23,
                xmode=linear,ymode=log,
              },
              row4KeysLin/.style={
                xlabel={\(\xi\)},ylabel={\(s(\xi)\)},
                xmode=log,
              },
              row4Keys/.style={row4KeysLin,ymode=log},
              row5KeysLin/.style={
                width=\plotLastWidth cm,
                xlabel={\(t\)},ylabel={\(s(t)\)},
                enlarge x limits=0,
                enlarge y limits=0.05,
                x tick scale label style={
                  yshift=-1em,
                  at={(.95,0)},
                  anchor=north,
                },
              },
              row5Keys/.style={
                row5KeysLin,
                ymode=log,
              },
            }
            %\endgroup
          %\endgroup
        }
        %\endgroup
      %\endgroup
    %\endgroup

    %\begingroup Diagrammi dei nodi
    \newcommand\nodeDiagramSettings{
      %\begingroup Dimensioni
        \pgfmathsetmacro\xPlotSep{10}
        \pgfmathsetmacro\yPlotSep{4}
        \pgfmathsetmacro\plotScale{.8}

        \pgfmathsetmacro\titleSize{11}
        \pgfmathsetmacro\commentSize{10}

        \pgfmathsetmacro\nodeSize{20pt} % Diametro del nodo
        \pgfmathsetmacro\starSize{2.25*\nodeSize}
        \pgfmathsetmacro\numberSize{.53*\nodeSize}
      %\endgroup

      %\begingroup Colori
        % \definecolor{cRed}{HTML}{D68490}
        % \definecolor{cPurple}{HTML}{7E6EAC}
        % \definecolor{cTeal}{HTML}{87BDB1}
        % \definecolor{cOrange}{HTML}{F2D0A4}
        % \definecolor{cEdgeLight}{HTML}{EEE8D5} 
        % \definecolor{cEdgeDark}{HTML}{998888}
        % \definecolor{cStar}{HTML}{FDF6E3}
        %
        \definecolor{color1}{HTML}{800000} % Deep Red
        \definecolor{color2}{HTML}{4682B4} % Steel Blue
        \definecolor{color3}{HTML}{6B8E23} % Olive Green
        \definecolor{color4}{HTML}{708090} % Slate Gray
        \definecolor{color5}{HTML}{BC8F8F} % Rosy Brown
        \definecolor{color6}{HTML}{008B8B} % Dark Cyan
        \definecolor{color7}{HTML}{B8860B} % Dark Goldenrod
      %\endgroup

      %\begingroup Testo
        \tikzset{
          titleStyle/.style={
            font=\textsize{\titleSize},
            align=center,
            anchor=south,
          },
          commentStyle/.style={font=\textsize{\commentSize}},
          arrowCommentStyle/.style={
            midway,
            commentStyle,
            left,
            sloped,
            below
          },
          numberStyle/.style={font=\textsize{\numberSize}}
        }
      %\endgroup

      %\begingroup Frecce
        \tikzset{
          transformingArrow/.style={
            % =Latex,
            ->,
            >=Stealth,
            dashed,
            line width=.8pt,
            black!90,
            shorten >=5pt,
            shorten <=10pt
          },
          indirectArrow/.style={
            draw=black!95,
            line width=.5pt,
          },
          nodeArrowExact/.style={
            indirectArrow,
            % ->,
            % -{Stealth[scale=1.2]},
            -{Triangle[scale=.75]},
            % draw=cEdgeLight,
            % draw=red!50,
            % shorten >=2pt,
            % shorten <=2pt
          },
          nodeArrowApprox/.style={
            indirectArrow,
            -,
            dashed,
          },
          transparentNodeArrow/.style={
            opacity=.4,
          },
          arrowAxis/.style={
            ->,
            line width=.3pt,
            >=Latex,
          },
        }
      %\endgroup

      %\begingroup Nodi
        \tikzset{
          nodeStyle/.style={
            circle,
            minimum size=\nodeSize,
            % draw=black!80,
            draw=black,
            line width=.5pt,
            inner sep=0pt,
          },
          starStyle/.style={
            star,
            star points=5,
            star point ratio=2,
            % fill=cStar,
            fill=black!50,
            fill opacity=.25,
            minimum size=\starSize,
            draw=none
          },
        }
      %\endgroup
    }
    %\endgroup

    %\begingroup GNU-Plot
      % \IfFileExists{\sourceauxpath.aux}{
      %   \externaldocument{\sourceauxpath}
      % }{\relax}
      
      \pgfkeys{/pgf/plot/gnuplot call={cd /d "\sourceoutdir" && gnuplot}}
    %\endgroup
  %\endgroup
\fi
%\endgroup

%\begingroup For the paper
  \def\tsc#1{\csdef{#1}{\textsc{\lowercase{#1}}\xspace}}
  \tsc{WGM}
  \tsc{QE}
  
  \newcommand\etal{et al.}
  \newcommand\KTMAS{\hyperlink{acrKineticTheoryMultiAgentSystems}{KTMAS}}
  
  % \newcommand\MathSpacing[3]{
  %   \thinmuskip = #1mu\relax
  %   \medmuskip = #2mu\relax
  %   \thickmuskip = #3mu\relax
  % }
  % \MathSpacing123
  
  \newcommand\RPlus{\mathbb R_+}
  \newcommand\NPlus{\mathbb N_+}
  \DeclareMathAlphabet\pazocal{OMS}{zplm}{m}{n}
  
  \DeclareMathOperator\Lognormal{Lognormal}
  \DeclareMathOperator\BiLognormal{BiLognormal}
  \DeclareMathOperator\Bernoulli{Bernoulli}
  \DeclareMathOperator\Exponential{Exp}
  \DeclareMathOperator\Student{Student}
  \DeclareMathOperator\Variance{var}
  \DeclareMathOperator\IC{IC}
  \DeclareMathOperator\Histogram{\mathcal B}
  \DeclareMathOperator\CCDF{\mathcal R}
    
  \newcommand\ThreeQuarters{{3\hspace{-0.15em}/\hspace{-0.15em}4}}
  \newcommand\OneQuarter{{1\hspace{-0.15em}/\hspace{-0.15em}4}}
  
  \newcommand\de[1][]{\mathrm{d}#1}
  \newcommand\derS[3][d]{\frac{\mathrm{#1}#2}{\mathrm{#1}#3}}
  \newcommand\derP[2]{\frac{\partial#1}{\partial#2}}
  \newcommand\underhat[1]{\underaccent{\hat}{#1}}
  \renewcommand\undertilde[1]{\underaccent{\tilde}{#1}}
  
  \makeatletter
  \newcommand\eqtag{\@ifstar{\eqtag@star}{\eqtag@nostar}}
  
  \newcounter{eqtag}
  \newcommand\eqtag@nostar[2]{%
    \begin{equation*}%
      \refstepcounter{eqtag}%
      % \protected@edef\@currentlabeltype{eqtag}%
      \protected@edef\@currentlabel{#1}%
      #2%
      \tag*{#1}%
    \end{equation*}%
  }
  \newcommand\eqtag@star[2]{%
    \begin{equation*}%
      \refstepcounter{eqtag}%
      % \protected@edef\@currentlabeltype{eqtag}%
      \protected@edef\@currentlabel{#1}%
      \let\label\ltx@label%
      #2%
      % \notag%
    \end{equation*}%
  }
  \makeatother
  
  \newcommand\crefsubfigInt[4][-]{Figures~\ref{#2}#3#1\ref{#2}#4}
  \newcommand\crefsubfigSep[4]{Figures~\ref{#1}#2\foreach \l in#3{,~\ref{#1}\l}~and~\ref{#1}#4}
  
  \crefname{section}{§}{§§}
  \crefrangeformat{section}{§§~#3#1#4--#5#2#6}
  \crefmultiformat{section}{§§~#2#1#3}{ and~#2#1#3}{, #2#1#3}{ and~#2#1#3}
  \crefrangemultiformat{section}{§§~#3#1#4--#5#2#6}{ and~#3#1#4--#5#2#6}{, #3#1#4--#5#2#6}{ and~#3#1#4--#5#2#6}
  
  \crefname{equation}{}{}
  \crefrangeformat{equation}{(#3#1#4--#5#2#6)}
  \crefmultiformat{equation}{(#2#1#3}{,~#2#1#3)}{,~#2#1#3}{,~#2#1#3)}
  \crefrangemultiformat{equation}{(#3#1#4--#5#2#6}{,~#3#1#4--#5#2#6)}{,~#3#1#4--#5#2#6}{,~#3#1#4--#5#2#6)}
  
  \crefname{figure}{Figure}{Figures}
  \crefrangeformat{figure}{Figures~#3#1#4--#5#2#6}
  \crefmultiformat{figure}{Figures~#2#1#3}{ and~#2#1#3}{, #2#1#3}{ and~#2#1#3}
  \crefrangemultiformat{figure}{Figures~#3#1#4--#5#2#6}{ and~#3#1#4--#5#2#6}{, #3#1#4--#5#2#6}{ and~#3#1#4--#5#2#6}
  
  \crefname{table}{Table}{Tables}
  \crefrangeformat{table}{Tables~#3#1#4--#5#2#6}
  \crefmultiformat{table}{Tables~#2#1#3}{ and~#2#1#3}{, #2#1#3}{ and~#2#1#3}
  \crefrangemultiformat{table}{Tables~#3#1#4--#5#2#6}{ and~#3#1#4--#5#2#6}{, #3#1#4--#5#2#6}{ and~#3#1#4--#5#2#6}
  
  \crefname{algorithm}{Algorithm}{Algorithms}
  \crefrangeformat{algorithm}{Algorithms~#3#1#4--#5#2#6}
  \crefmultiformat{algorithm}{Algorithms~#2#1#3}{ and~#2#1#3}{, #2#1#3}{ and~#2#1#3}
  \crefrangemultiformat{algorithm}{Algorithms~#3#1#4--#5#2#6}{ and~#3#1#4--#5#2#6}{, #3#1#4--#5#2#6}{ and~#3#1#4--#5#2#6}
  
  %\newtheorem{theorem}{Theorem}
  %\newtheorem{lemma}[theorem]{Lemma}
  \ifcasscclass
    % «cas-sc» fornisce \newdefinition
    \newdefinition{dfn}{Definition}
    \newdefinition{rmk}{Remark}
    \newdefinition{hyp}{Hypotesis}
  \else
    % Con «article» (Articolo) si usano gli ambienti del kernel, con
    % contatore condiviso e «Assumption» come titolo di «hyp»
    \newtheorem{dfn}{Definition}
    \newtheorem{rmk}[dfn]{Remark}
    \newtheorem{hyp}[dfn]{Assumption}
  \fi
  %\newproof{pf}{Proof}
  %\newproof{pot}{Proof of Theorem \ref{thm}}
  
  \crefname{dfn}{Definition}{Definitions}
  \crefrangeformat{dfn}{Definitions~#3#1#4--#5#2#6}
  \crefmultiformat{dfn}{Definitions~#2#1#3}{ and~#2#1#3}{, #2#1#3}{ and~#2#1#3}
  \crefrangemultiformat{dfn}{Definitions~#3#1#4--#5#2#6}{ and~#3#1#4--#5#2#6}{, #3#1#4--#5#2#6}{ and~#3#1#4--#5#2#6}
  \crefname{hyp}{Hypothesis}{Hypotheses}
  
  \crefrangeformat{hyp}{Hypotheses~#3#1#4--#5#2#6}
  \crefmultiformat{hyp}{Hypotheses~#2#1#3}{ and~#2#1#3}{, #2#1#3}{ and~#2#1#3}
  \crefrangemultiformat{hyp}{Hypotheses~#3#1#4--#5#2#6}{ and~#3#1#4--#5#2#6}{, #3#1#4--#5#2#6}{ and~#3#1#4--#5#2#6}
  
  \crefname{rmk}{Remark}{Remarks}
  \crefrangeformat{rmk}{Remarks~#3#1#4--#5#2#6}
  \crefmultiformat{rmk}{Remarks~#2#1#3}{ and~#2#1#3}{, #2#1#3}{ and~#2#1#3}
  \crefrangemultiformat{rmk}{Remarks~#3#1#4--#5#2#6}{ and~#3#1#4--#5#2#6}{, #3#1#4--#5#2#6}{ and~#3#1#4--#5#2#6}
  
  % \newcommand\tikzfigure[5][0em]{%
  %   \begin{figure}[%
  %     pos=#2,%
  %     width=#4,%
  %     align=\centering%
  %   ]%
  %     \centering%
  %     \vspace*{#1}%
  %     \IfSubStr{#2}{p}{%
  %       \leavevmode\stepcounter{figposcount}%
  %       \zsaveposx{figpos-\thefigposcount}%
  %       \makebox[0pt][l]{%
  %         \hspace*{-\dimexpr\zposx{figpos-\thefigposcount}sp\relax}%
  %         \hspace*{\dimexpr0.5\paperwidth\relax}%
  %         \makebox[0pt][c]{%
  %           \scalebox{0.86}{\includegraphics{figs/#5.pdf}}%
  %         }%
  %       }
  %     }{%
  %       \includegraphics{figs/#5.pdf}%
  %     }%
  %     \caption{%
  %       \hyphenpenalty=0%
  %       \exhyphenpenalty=0%
  %       #3%
  %     }%
  %     \protect\label{#5}%
  %   \end{figure}%
  % }
%\endgroup

%\begingroup Versione per la classe «cas-sc»
  % 
  % %\begingroup For image generation
  %   %\begingroup Tipografia
  %     \newcommand\textsize[1]{%
  %       \pgfmathparse{#1*1.2}%
  %       \fontsize{#1pt}{\pgfmathresult pt}\selectfont%
  %     }
  % 
  %     \DeclareMathOperator\FRC{\mathcal R}
  % 
  %     \newif\ifcasscclass
  %     \makeatletter
  %     \@ifclassloaded{cas-sc}{\casscclasstrue}{\casscclassfalse}
  %     \makeatother
  % 
  %     \providecommand\Ifthispageodd[2]{\checkoddpage\ifoddpage#1\else#2\fi}
  %   %\endgroup
  % 
  %   %\begingroup TikZ
  %     %\begingroup Didascalie
  %       \newcommand\newcaption[2]{%
  %         \captionsetup{width=#2} % aboveskip=-.2cm,
  %         \caption{#1}%
  %       }
  %       
  %       \newcommand\newsubcaption[2][ ]{%
  %         \refstepcounter{subfigure}%
  %         (\thesubfigure)#1#2%
  %       }
  %     %\endgroup
  %     
  %     %\begingroup RedefineTikZbounds
  %       % La differenza è che colla prima bisogna catturare anche le etichette perché si crea un pdf isolato mentre colla non è importante perché viene tutto compilato sul momento
  % 
  %       %\begingroup Versione valida con esternalizzazione [ma anche senza]
  %         \newcommand\redefineTikZbounds[2]{
  %           \coordinate (RealSW) at (current bounding box.south west);
  %           \coordinate (RealNE) at (current bounding box.north east);
  %           \coordinate (Center) at ($(#1.west)!0.5!(#2.east)$);
  %     
  %           \pgfresetboundingbox 
  % 
  %           \useasboundingbox 
  %             let \p1=(RealSW),\p2=(RealNE),\p3=(Center),
  %             \n1={max(\x3-\x1,\x2-\x3)} in
  %             (\x3-\n1,\y1) rectangle (\x3+\n1,\y2+.25em);
  % 
  %           % \path[use as bounding box] 
  %           %   let \p1=(RealSW),\p2=(RealNE),\p3=(Center),
  %           %   \n1={max(\x3-\x1,\x2-\x3)} in
  %           %   (\x3-\n1,\y1) rectangle (\x3+\n1,\y2+.25em);
  %           % Lo «.25em» discende dalla traslazione imposta all'etichetta della scala verticale mediante «yScaleLabelShift»
  %         } % Scatola delimitante simmetrica
  %       %\endgroup
  % 
  %       %\begingroup Versione valida senza esternalizzazione
  %         % \newcommand\redefineTikZbounds[2]{
  %         %   % Salva le coordinate della cima e del fondo della figura corrente (cosí da mantenere i titoli e le didascalie)
  %         %   \coordinate (Cima) at (current bounding box.north);
  %         %   \coordinate (Fondo) at (current bounding box.south);
  %     
  %         %   \pgfresetboundingbox % Resetta la scatola dei limiti
  %     
  %         %   \useasboundingbox % Definisce la nuova scatola come un rettangolo
  %         %     (#1.west |- Cima)  % con ascissa sinistra data da «Col1 c1r1.west», 
  %         %     rectangle      % ascissa destra data da «Col2 c1r1.east»
  %         %     (#2.east |- Fondo); % ordinate date dalla cima e fondo originali
  %         %     % L'operatore «|-» in «(A |- B)» prende come ascissa quella del punto A mentre come coordinata quella del punto B
  %         % }
  %       %\endgroup
  %     %\endgroup
  % 
  %     %\begingroup TikZBox
  %       % Parzialmente ispirato da https://gist.github.com/Caramdir/899107
  % 
  %       \makeatletter
  %       \newsavebox\tikzBufferBox
  %       \xdef\offsetX{0pt}
  %       \xdef\offsetY{0pt}
  %       
  %       % \tikzbox[#1]{#2}{#3}[#4]{#5}
  %       % #1 -> Ancora rispetto cui posizionare l'intera scatola (di base è «south west», ossia l'angolo in basso a sinistra)
  %       % #2 -> Nodo dove posizionare l'origine (per la natura di come LaTeX riceve gli argomenti è buon pratica passare questo circondato da parentesi graffe)
  %       % #3 -> Nome della scatola limitante definita dall'ambiente «scope»
  %       % #4 -> Argomenti opzionali da passare all'ambiente «scope»; di base sono vuoti
  %       % #5 -> Contenuto della scatola circondato dall'ambiente «scope»
  %       % Es. \tikzbox{{(0,0)}}{figuraTikZ}{blocco TikZ}
  % 
  %       % Macro di apertura che controlla se v'è l'argomento facoltativo dell'ancora
  %       \def\tikzbox{\pgfutil@ifnextchar[{\tikzbox@anchor}{\tikzbox@anchor[south west]}}
  %       % La funzione «\pgfutil@ifnextchar» serve per controllare se il prossimo carattere è l'apertura di una paretensi graffa «[»
  %       
  %       % Macro intermedia per cogliere gli argomenti obbligatori e controllare quelli successivi facoltativi
  %       \def\tikzbox@anchor[#1]#2#3{\pgfutil@ifnextchar[{\tikzbox@opt[#1]#2#3}{\tikzbox@opt[#1]#2#3[]}}
  % 
  %       % Macro degli argomenti obbligatori
  %       \long\def\tikzbox@opt[#1]#2#3[#4]#5{%
  %         \node[ % Misurazione della scatola in un nodo invisibile
  %           overlay,%
  %           opacity=0,%
  %           inner sep=0pt,%
  %           outer sep=0pt%
  %         ] {%
  %           \sbox{\tikzBufferBox}{%
  %             \begin{tikzpicture}%
  %               \begin{scope}[local bounding box=measuringTikZBox]%
  %                 #5%
  %               \end{scope}%
  %               \pgfpointanchor{measuringTikZBox}{#1}%
  %               \pgfgetlastxy{\tempX}{\tempY}% 
  %               \xdef\offsetX{\tempX}%
  %               \xdef\offsetY{\tempY}%
  %             \end{tikzpicture}%
  %           }%
  %         };%
  %         %
  %         % Disegno reale traslato
  %         \begin{scope}[%
  %           local bounding box=#3,%
  %           shift={#2},%
  %           xshift=-\offsetX,%
  %           yshift=-\offsetY,%
  %           #4%
  %         ]%
  %           #5%
  %         \end{scope}%
  %       }%
  %       \makeatother
  %     %\endgroup
  % 
  %     %\begingroup Chiavi-Valori
  %       %\begingroup Testo
  %         % La dimensione del testo di base è 11pt
  %         \newcommand\titleSize{11}
  %         \newcommand\subCaptionSize{10}
  %         \newcommand\ParameterSize{9}
  %         \newcommand\labelSize{10}
  %         \newcommand\tickSize{9}
  %         \newcommand\legendSize{9}
  %       %\endgroup
  %     
  %       %\begingroup Infopagine
  %         %\begingroup Chiavi condivise
  %         \newcounter{lpc} % «loopcounter»
  %         \newcommand\infopageSharedKeys{
  %           %\begingroup Testo
  %             \renewcommand\titleSize{10}
  %             \renewcommand\subCaptionSize8
  %             \renewcommand\ParameterSize7
  %             \renewcommand\labelSize9
  %             \renewcommand\tickSize6
  %             \renewcommand\legendSize7
  %           %\endgroup
  %     
  %           %\begingroup Dimensioni
  %             \pgfmathsetmacro\plotWidth7
  %             \pgfmathsetmacro\plotHeight{.618*\plotWidth}
  %             \pgfmathsetmacro\halfPlotHeight{\plotHeight/2}
  %             %
  %             \pgfmathparse{0.825*\plotWidth}
  %             \pgfmathsetmacro\plotWidth\pgfmathresult
  %             \pgfmathsetmacro\halfWidth{\plotWidth/2}
  %             \pgfmathsetmacro\plotLastWidth{1.25*\plotWidth}
  %     
  %             \pgfmathsetmacro\xPlotSep{1}
  %             \pgfmathsetmacro\yPlotSep{1.25}
  %             \pgfmathsetmacro\yGroupSep{2.75}
  %           %\endgroup
  %     
  %           %\begingroup Legende
  %             \newcommand\placeLegend[5]{%
  %               /pgfplots/plotLegend,
  %               at={([shift={(##4,##5)}]axis description cs:##1,##2)},
  %               anchor=##3,
  %             }
  %             %
  %             \pgfplotsset{%
  %               plotLegend/.style={%
  %                 font=\textsize{\legendSize},
  %                 /pgfplots/legend cell align=left,
  %                 % align=left,
  %                 % draw=white!15!black,
  %                 draw=none, % Bordo invisibile
  %                 fill=white,
  %                 fill opacity=0.5,
  %                 text opacity=1,
  %                 % text width=2cm,
  %               },
  %               plotLegendNW/.style={legend style={\placeLegend{0}{1}{north west}{0em}{0em}}},
  %               plotLegendNE/.style={legend style={\placeLegend{1}{1}{north east}{0em}{0em}}},
  %               plotLegendSW/.style={legend style={\placeLegend{0}{0}{south west}{0em}{0em}}},
  %               plotLegendSE/.style={legend style={\placeLegend{1}{0}{south east}{0em}{0em}}},
  %               %
  %               % plotLegendW/.style={legend style={plotLegend,at={(0,.5)},anchor=west}},
  %               %
  %               % plotLegendNW/.style={legend style={plotLegend,legend pos=north west}},
  %               % plotLegendNE/.style={legend style={plotLegend,legend pos=north east}},
  %               % plotLegendSW/.style={legend style={plotLegend,legend pos=south west}},
  %               % plotLegendSE/.style={legend style={plotLegend,legend pos=south east}},
  %               %
  %               % plotLegendNWcloser/.style={legend style={plotLegend,at={(0,1)},anchor=north west}},
  %               % plotLegendNEcloser/.style={legend style={plotLegend,at={(1,1)},anchor=north east}},
  %               % plotLegendSWcloser/.style={legend style={plotLegend,at={(0,0)},anchor=south west}},
  %               % plotLegendSEcloser/.style={legend style={plotLegend,at={(0,1)},anchor=south east}},
  %             }
  %             %
  %             \pgfmathsetmacro\lineLegendLength{1}
  %             \pgfmathsetmacro\areaLegendHeigth{.25}
  %             \pgfplotsset{
  %               legend image code/.code={
  %                 \draw [mark repeat=2,mark phase=2,##1]
  %                   plot coordinates {
  %                     (0em,0em)
  %                     (.5*\lineLegendLength em,0em)
  %                     (\lineLegendLength em,0em)
  %                   };
  %               },
  %               /pgfplots/area legend/.style={
  %                 legend image code/.code={
  %                   \draw [##1,fill,draw=none] 
  %                     (0em,-\areaLegendHeigth em) 
  %                     rectangle 
  %                     (\lineLegendLength em,\areaLegendHeigth em);
  %                 },
  %               },
  %             }
  %             %
  %             \newcommand\ExactHistogramLegend{Exact histogram}
  %             \newcommand\ApproxHistogramLegend{Approx. histogram}
  %             \newcommand\RealHistogramLegend{Real histogram}
  %             %
  %             \newcommand\ExactBLNFitLegend{Exact BLN fit}
  %             \newcommand\ApproxBLNFitLegend{Approx BLN fit}
  %             \newcommand\RealBLNFitLegend{Real BLN fit}
  %             %
  %             \newcommand\ExactEmpiricCCDFLegend{Exact empiric CCDF}
  %             \newcommand\ApproxEmpiricCCDFLegend{Approx. empiric CCDF}
  %             \newcommand\RealEmpiricCCDFLegend{Real empiric CCDF}
  %             %
  %             \newcommand\ExactParetoFitLegend{Exact Pareto fit}
  %             \newcommand\ApproxParetoFitLegend{Approx. Pareto fit}
  %             \newcommand\RealParetoFitLegend{Real Pareto fit}
  %             %
  %             \newcommand\ExactScatterLegend{Exact scatter}
  %             \newcommand\ApproxScatterLegend{Approx. scatter}
  %             \newcommand\RealScatterLegend{Real scatter}
  %             %
  %             \newcommand\ClasskILegend{Class k=283}
  %             \newcommand\ClasskIILegend{Class k=83}
  %             \newcommand\ClasskIIILegend{Class k=43}
  %             \newcommand\ClasskIVLegend{Class k=10}
  %           %\endgroup
  %     
  %           %\begingroup Etichette vuote
  %             \pgfplotsset{
  %               yLabelEmpty/.style={ylabel={}},
  %               yTickLabelsEmpty/.style={yticklabels={}},
  %               yTickScaleLabelEmpty/.style={ytick scale label code/.code={}},
  %               xTickScaleLabelEmpty/.style={xtick scale label code/.code={}},
  %               emptyLeged/.style={
  %                 legend style={draw=none,fill=none,nodes={draw=none}},
  %                 legend image post style={opacity=0},
  %               },
  %             }
  %           %\endgroup
  %     
  %           %\begingroup Estetica
  %             \newcommand\smallminus{\scalebox{0.5}[1]{$-$}\hspace{-.11em}}
  %             \pgfplotsset{
  %               /pgfplots/log number format basis/.code 2 args={%
  %                 \pgfmathsetmacro\exponent{##2}%
  %                 \pgfmathparse{sign(\exponent)}%
  %                 \ifnum\pgfmathresult<0%
  %                   \pgfmathparse{abs(\exponent)}%
  %                   $##1^{\smash{\smallminus\pgfmathprintnumber{\pgfmathresult}}}$%
  %                 \else%
  %                   $##1^{\smash{\pgfmathprintnumber{\exponent}}}$%
  %                 \fi%
  %               },
  %               ytick scale label code/.code={%
  %                 \pgfmathsetmacro\exponent{##1}%
  %                 \pgfmathparse{sign(\exponent)}%
  %                 \ifnum\pgfmathresult<0
  %                   \pgfmathparse{abs(\exponent)}%
  %                   $\cdot 10^{\smash{\smallminus\pgfmathprintnumber{\pgfmathresult}}}$%
  %                 \else
  %                   $\cdot 10^{\smash{\pgfmathprintnumber{\exponent}}}$%
  %                 \fi
  %               },
  %             }
  %           %\endgroup
  %     
  %           %\begingroup Colori
  %             \definecolor{blueish}{rgb}{0.00000,0.44700,0.74100}%
  %             \definecolor{reddish}{rgb}{0.83529,0.36863,0.00000}%
  %           %\endgroup
  %     
  %           %\begingroup Commenti
  %             %\begingroup Impostazioni
  %             \tikzset{
  %               title/.style={
  %                 align=center,
  %                 anchor=south,
  %                 font=\bfseries\textsize{\titleSize},
  %               },
  %               fittedParameters/.style={
  %                 align=center,
  %                 font=\textsize{\ParameterSize},
  %               },
  %               varyingParameter/.style={
  %                 align=center,
  %                 % rotate=90,
  %                 font=\textsize{\subCaptionSize},
  %               },
  %               subCaption/.style={
  %                 % text width=\plotWidth cm,
  %                 align=center,
  %                 font=\textsize{\subCaptionSize},
  %               },
  %             }
  %             \newcommand\codetocolor[1]{\ifcase##1blueish\or reddish\else black\fi}
  %             %\endgroup
  %     
  %             %\begingroup Parametri Sardegna
  %             \newcommand\betaReal{\num{1.273}}
  %             \newcommand\xiReal{\num{0.929}}
  %             \newcommand\mIReal{\num{3.227}}
  %             \newcommand\sIReal{\num{0.408}}
  %             \newcommand\mIIReal{\num{4.099}}
  %             \newcommand\sIIReal{\num{0.556}}
  %             %\endgroup
  %     
  %             %\begingroup Parametri Pareto
  %             \newcommand\codetoParetovalue[1]{
  %               \ifcase##1\betaExact\or\betaAprx\else\betaReal\fi
  %             }
  %             \newcommand\writeParetoParameters[7][center]{
  %               \node[fittedParameters,anchor=##1] at ([shift={(##5,##6)}]##4) {%
  %                 \(\begin{aligned}
  %                   \color{\codetocolor{##2}}\beta
  %                   &\color{\codetocolor{##2}}=\codetoParetovalue{##2}\\[##7]
  %                   \color{\codetocolor{##3}}\beta
  %                   &\color{\codetocolor{##3}}=\codetoParetovalue{##3}
  %                 \end{aligned}\)
  %               };
  %             }
  %             %\endgroup
  % 
  %             %\begingroup % Vecchia versione in cui si scrivevano esplicitamente tutte le 9 combinazioni possibili
  %             % \newcommand\writeParameters[6]{%
  %             %   \node[parameter] at ([shift={(##4,##5)}]##1) {%
  %             %     \ifcase##2% Caso 0
  %             %       \ifcase##3% Caso 0
  %             %         \(\begin{aligned}
  %             %           \color{blueish}\beta&\color{blueish}=\betaExact\\[##6]
  %             %           \color{blueish}\beta&\color{blueish}=\betaExact
  %             %         \end{aligned}\)
  %             %       \or% Caso 1
  %             %         \(\begin{aligned}
  %             %           \color{blueish}\beta&\color{blueish}=\betaExact\\[##6]
  %             %           \color{reddish}\beta&\color{reddish}=\betaAprx% Caso 1
  %             %         \end{aligned}\)
  %             %       \else% Caso 2
  %             %         \(\begin{aligned}
  %             %           \color{blueish}\beta&\color{blueish}=\betaExact\\[##6]
  %             %           \color{black}\beta&\color{black}=\betaReal% Caso 2
  %             %         \end{aligned}\)
  %             %       \fi
  %             %     \or% Caso 1
  %             %       \ifcase##3% Caso 0
  %             %         \(\begin{aligned}
  %             %           \color{reddish}\beta&\color{reddish}=\betaAprx\\[##6]
  %             %           \color{blueish}\beta&\color{blueish}=\betaExact
  %             %         \end{aligned}\)
  %             %       \or% Caso 1
  %             %         \(\begin{aligned}
  %             %           \color{reddish}\beta&\color{reddish}=\betaAprx\\[##6]
  %             %           \color{reddish}\beta&\color{reddish}=\betaAprx
  %             %         \end{aligned}\)
  %             %       \else% Caso 2
  %             %         \(\begin{aligned}
  %             %           \color{reddish}\beta&\color{reddish}=\betaAprx\\[##6]
  %             %           \color{black}\beta&\color{black}=\betaReal
  %             %         \end{aligned}\)
  %             %       \fi
  %             %     \else% Caso 2
  %             %       \ifcase##3% Caso 0
  %             %         \(\begin{aligned}
  %             %           \color{black}\beta&\color{black}=\betaReal\\[##6]
  %             %           \color{blueish}\beta&\color{blueish}=\betaExact
  %             %         \end{aligned}\)
  %             %       \or% Caso 1
  %             %         \(\begin{aligned}
  %             %           \color{black}\beta&\color{black}=\betaReal\\[##6]
  %             %           \color{reddish}\beta&\color{reddish}=\betaAprx
  %             %         \end{aligned}\)
  %             %       \else% Caso 2
  %             %         \(\begin{aligned}
  %             %           \color{black}\beta&\color{black}=\betaReal\\[##6]
  %             %           \color{black}\beta&\color{black}=\betaReal
  %             %         \end{aligned}\)
  %             %       \fi
  %             %     \fi
  %             %   };
  %             % }
  %             %\endgroup
  %     
  %             %\begingroup Parametri bilognormale
  %             \newcommand\codetoparameter[1]{
  %               \ifcase##1\xi  % Caso 0
  %               \or\mu_1     % Caso 1
  %               \or\sigma_1   % Caso 2
  %               \or\mu_2     % Caso 3
  %               \else\sigma_2\fi % Caso 4
  %             }
  %             \newcommand\codetoBiLognormalvalue[2]{
  %               \ifcase##1 % Caso 0
  %                 \ifcase##2\xiExact % Caso 0
  %                 \or\mIExact    % Caso 1
  %                 \or\sIExact    % Caso 2
  %                 \or\mIIExact    % Caso 3
  %                 \else\sIIExact\fi % Caso 4
  %               \or % Caso 1
  %                 \ifcase##2\xiAprx % Caso 0
  %                 \or\mIAprx     % Caso 1
  %                 \or\sIAprx     % Caso 2
  %                 \or\mIIAprx    % Caso 3
  %                 \else\sIIAprx\fi  % Caso 4
  %               \else % Caso 2
  %                 \ifcase##2\xiReal % Caso 0
  %                 \or\mIReal     % Caso 1
  %                 \or\sIReal     % Caso 2
  %                 \or\mIIReal    % Caso 3
  %                 \else\sIIReal\fi  % Caso 4
  %               \fi
  %             }
  %             \newcommand\matrixline[2]{\codetoparameter{##2}&=\codetoBiLognormalvalue{##1}{##2}}
  %             \def\lineloop##1##2{%
  %               \matrixline{##1}{\number\value{lpc}}\stepcounter{lpc}%
  %               \ifnum\number\value{lpc}<5\\[##2]\lineloop{##1}{##2}\fi%
  %             } % https://tex.stackexchange.com/questions/491358/tikz-foreach-to-work-with-the-aligned-environment
  %             \newcommand\writeBiLognormalParameters[6][center]{
  %               \node[fittedParameters,anchor=##1] at ([shift={(##4,##5)}]##3) {%
  %                 \setcounter{lpc}{0}
  %                 \(\color{\codetocolor{##2}}
  %                 \begin{aligned}
  %                   \lineloop{##2}{##6}
  %                 \end{aligned}\)
  %               };
  %             }
  %             %\endgroup
  %     
  %             %\begingroup Altre etichette
  %             \newcommand\writeSizeEvolutionDescriptions[1][5]{%
  %               \foreach \col/\clr/\txt in {%
  %                 1/blueish/Esatto,
  %                 2/reddish/Approssimato
  %               }{
  %                 \node[
  %                   fittedParameters,
  %                   anchor=south,
  %                   yshift=.1em,
  %                 ] at (Row##1 c\col r1.south) {\color{\clr}\txt};
  %               }
  %             }
  %     
  %             \newcommand\writeVaryingParameter[5]{%
  %               \foreach \r in {1,...,5}{
  %                 \node[
  %                   varyingParameter,
  %                   anchor=##3,
  %                 ] at ([shift={(##4,##5)}]Row\r\space c##1r1.##2)
  %                 {\(\varyingParameterSymbol=\varyingParameterList\r\)};
  %               }
  %             }
  %             %\endgroup
  %           %\endgroup
  %     
  %           %\begingroup Sottodidascalie
  %             \newcommand\writeSubCaption[5]{%
  %               \node[subCaption,anchor=##3]
  %               at ([shift={(##4,##5)}]##1.##2)
  %               {\newsubcaption[]{}};
  %             }
  %             \newcommand\writeSubCaptions[5][3]{%
  %               \foreach \r in {1,...,5}{
  %                 \foreach \c in {1,...,##1}{
  %                   \writeSubCaption{Row\r\space c\c r1}{##2}{##3}{##4}{##5}
  %                 }
  %               }
  %             }
  %           %\endgroup
  %         }
  %         %\endgroup
  %     
  %         %\begingroup Chiavi del primo tipo d'infopagina
  %         \newcommand\infopageKeysI{
  %           \infopageSharedKeys
  %     
  %           %\begingroup Gruppi
  %             \ifpageisodd
  %               \pgfplotsset{/pgfplots/group/yDescPos/.style={y descriptions at=edge right}}
  %             \else
  %               \pgfplotsset{/pgfplots/group/yDescPos/.style={y descriptions at=edge left}}
  %             \fi
  %             \pgfplotsset{
  %               /pgfplots/group/plotGroup/.style={
  %                 horizontal sep=\xPlotSep em,
  %                 vertical sep=\yPlotSep em,
  %                 yDescPos,
  %                 group size=3 by 1,
  %               },
  %               /pgfplots/group/plotLastGroup/.style={
  %                 plotGroup,
  %                 x descriptions at=edge bottom,
  %               },
  %             }
  %           %\endgroup
  %     
  %           %\begingroup Stili
  %             %\begingroup Selezione «recto»/«verso»
  %             \ifpageisodd
  %               \pgfplotsset{%
  %                 yLineAxis/.style={axis y line*=right},
  %                 /tikz/yLabelShift/.style={%
  %                   yshift=-1.25em,
  %                   at={(axis description cs:1,0.5)},
  %                   anchor=north,
  %                   % rotate=180,
  %                 },
  %                 /tikz/yScaleLabelShift/.style={%
  %                   yshift=.25em,
  %                   at={(1,1)},
  %                   anchor=south east,
  %                 },
  %                 /tikz/yTickLabelShift/.style={xshift=-0.15em},
  %               }
  %             \else
  %               \pgfplotsset{%
  %                 yLineAxis/.style={axis y line*=left},
  %                 /tikz/yLabelShift/.style={%
  %                   yshift=1.25em,
  %                   at={(axis description cs:0,0.5)},
  %                   anchor=south,
  %                   % rotate=180,
  %                 },
  %                 /tikz/yScaleLabelShift/.style={%
  %                   yshift=.25em,
  %                   at={(0,1)},
  %                   anchor=south west,
  %                 },
  %                 /tikz/yTickLabelShift/.style={xshift=+0.15em},
  %               }
  %             \fi
  %             %\endgroup
  %             %
  %             %\begingroup Grafici
  %             \pgfplotsset{
  %               plotRow/.style={
  %                 scale only axis,
  %                 width=\plotWidth cm,height=\plotHeight cm,
  %                 xlabel style={font=\color{white!15!black}\textsize{\labelSize}},
  %                 ylabel style={font=\color{white!15!black}\textsize{\labelSize},yLabelShift},
  %                 xticklabel style={font=\textsize{\tickSize}},
  %                 yticklabel style={font=\textsize{\tickSize},yTickLabelShift},
  %                 y tick scale label style={yScaleLabelShift},
  %                 % axis x line*=bottom,yLineAxis,
  %                 enlarge x limits=0.1,enlarge y limits=0.1,
  %                 xmajorgrids,ymajorgrids,
  %                 xminorticks=true,yminorticks=true,
  %                 scale ticks above exponent=2,
  %                 axis background/.style={fill=white},
  %                 grid style={dashed}
  %               },
  %               % xlabel123/.style={xlabel={\(h\)}},
  %               xlabel3/.style={xlabel={\(h\)}},
  %               ylabel23/.style={ylabel={\(\FRC(h)\)}},
  %               row1Keys/.style={
  %                 % ylabel={\(\underhat{\bar l}_N(h)\)},
  %                 ylabel={\(\bar l_N(h)\)},
  %                 % enlarge x limits={lower,value=.05},
  %                 enlarge x limits=.05,
  %                 enlarge y limits={upper,value=.05},
  %                 xmode=linear,
  %                 log origin=infty,
  %               },
  %               row2KeysWithLabels/.style={
  %                 ylabel23,
  %                 xmode=linear,ymode=log
  %               },
  %               row2Keys/.style={
  %                 ylabel23,
  %                 xmode=linear,ymode=log,
  %                 xticklabels={},
  %               },
  %               row3Keys/.style={
  %                 xlabel3,ylabel23,
  %                 xmode=linear,ymode=log
  %               },
  %               row4Keys/.style={
  %                 xlabel={\(\xi\)},ylabel={\(s(\xi)\)},
  %                 xmode=log,ymode=log
  %               },
  %               row5KeysLin/.style={
  %                 width=\plotLastWidth cm,
  %                 xlabel={\(t\)},ylabel={\(s(t)\)},
  %                 enlarge x limits=0,
  %                 enlarge y limits=0.05,
  %                 x tick scale label style={
  %                   yshift=-1em,
  %                   at={(.95,0)},
  %                   anchor=north,
  %                 },
  %               },
  %               row5Keys/.style={
  %                 row5KeysLin,
  %                 ymode=log,
  %               },
  %             }
  %             %\endgroup
  %           %\endgroup
  %         }
  %         %\endgroup
  %     
  %         %\begingroup Chiavi del secondo tipo d'infopagina
  %         \newcommand\infopageKeysII{
  %           \infopageSharedKeys
  %     
  %           %\begingroup Gruppi
  %             \ifpageisodd
  %               \pgfplotsset{/pgfplots/group/yDescPos/.style={y descriptions at=edge right}}
  %             \else
  %               \pgfplotsset{/pgfplots/group/yDescPos/.style={y descriptions at=edge left}}
  %             \fi
  %             \pgfplotsset{
  %               /pgfplots/group/plotGroup/.style={
  %                 horizontal sep=\xPlotSep em,
  %                 vertical sep=\yPlotSep em,
  %                 y descriptions at=all,
  %                 group size=3 by 1,
  %               },
  %               /pgfplots/group/plotLastGroup/.style={
  %                 plotGroup,
  %                 x descriptions at=edge bottom,
  %               },
  %             }
  %           %\endgroup
  %     
  %           %\begingroup Stili
  %             %\begingroup Selezione «recto»/«verso»
  %             \ifpageisodd
  %               \pgfplotsset{
  %                 yLabel/.style={
  %                   ylabel style={
  %                     font=\color{white!15!black}\textsize{\labelSize},
  %                     yshift=-1.25em,
  %                     at={(axis description cs:1,0.5)},
  %                     anchor=north,
  %                     % rotate=180,
  %                   },
  %                 },
  %                 yLabelInv/.style={yLabelEmpty},
  %                 yScaleLabelShift/.style={
  %                   y tick scale label style={
  %                     yshift=.25em,
  %                     % xshift=1em,
  %                     at={(1,1)},
  %                     anchor=south east,
  %                   },
  %                 },
  %                 yScaleLabelShiftInv/.style={
  %                   y tick scale label style={
  %                     yshift=.25em,
  %                     % xshift=-1em,
  %                     at={(0,1)},
  %                     anchor=south west,
  %                   },
  %                 },
  %                 % yLineAxis/.style={axis y line*=right},
  %               }
  %             \else
  %               \pgfplotsset{
  %                 yLabel/.style={yLabelEmpty},
  %                 yLabelInv/.style={
  %                   ylabel style={
  %                     font=\color{white!15!black}\textsize{\labelSize},
  %                     yshift=1.25em,
  %                     at={(axis description cs:0,0.5)},
  %                     anchor=south,
  %                     % rotate=180,
  %                   },
  %                 },
  %                 yScaleLabelShift/.style={
  %                   y tick scale label style={
  %                     yshift=.25em,
  %                     % xshift=1em,
  %                     at={(1,1)},
  %                     anchor=south east,
  %                   },
  %                 },
  %                 yScaleLabelShiftInv/.style={
  %                   y tick scale label style={
  %                     yshift=.25em,
  %                     % xshift=-1em,
  %                     at={(0,1)},
  %                     anchor=south west,
  %                   },
  %                 },
  %                 % yLineAxis/.style={axis y line*=left},
  %               }
  %             \fi
  %             %\endgroup
  %             %
  %             %\begingroup Grafici
  %             \pgfplotsset{
  %               plotRow/.style={
  %                 scale only axis,
  %                 width=\plotWidth cm,height=\plotHeight cm,
  %                 xlabel style={font=\color{white!15!black}\textsize{\labelSize}},
  %                 % ylabel style={font=\color{white!15!black}\textsize{\labelSize},yLabelShift},
  %                 xticklabel style={font=\textsize{\tickSize}},
  %                 % axis x line*=bottom,yLineAxis,
  %                 enlarge x limits=0.1,enlarge y limits=0.1,
  %                 xmajorgrids,ymajorgrids,
  %                 xminorticks=true,yminorticks=true,
  %                 scale ticks above exponent=2,
  %                 axis background/.style={fill=white},
  %                 grid style={dashed},
  %               },
  %               plotLeftCol/.style={
  %                 yLabelInv,
  %                 yticklabel style={font=\textsize{\tickSize},xshift=0.15em},
  %                 yticklabel pos=left,
  %                 yScaleLabelShiftInv,
  %               },
  %               plotMiddleCol/.style={
  %                 yLabelEmpty,
  %                 yTickLabelsEmpty,
  %                 yTickScaleLabelEmpty,
  %               },
  %               plotRightCol/.style={
  %                 yLabel,
  %                 yticklabel style={font=\textsize{\tickSize},xshift=-0.15em},
  %                 yticklabel pos=right,
  %                 yScaleLabelShift,
  %               },
  %               xlabel3/.style={xlabel={\(h\)}},
  %               ylabel23/.style={ylabel={\(\FRC(h)\)}},
  %               row1Keys/.style={
  %                 ylabel={\(\bar l_N(h)\)},
  %                 enlarge x limits=.05,
  %                 enlarge y limits={upper,value=.05},
  %                 xmode=linear,
  %                 log origin=infty,
  %               },
  %               row2Keys/.style={
  %                 ylabel23,
  %                 xmode=linear,ymode=log,
  %                 % xticklabels={},
  %               },
  %               row3Keys/.style={
  %                 xlabel3,ylabel23,
  %                 xmode=linear,ymode=log,
  %               },
  %               row4KeysLin/.style={
  %                 xlabel={\(\xi\)},ylabel={\(s(\xi)\)},
  %                 xmode=log,
  %               },
  %               row4Keys/.style={row4KeysLin,ymode=log},
  %               row5KeysLin/.style={
  %                 width=\plotLastWidth cm,
  %                 xlabel={\(t\)},ylabel={\(s(t)\)},
  %                 enlarge x limits=0,
  %                 enlarge y limits=0.05,
  %                 x tick scale label style={
  %                   yshift=-1em,
  %                   at={(.95,0)},
  %                   anchor=north,
  %                 },
  %               },
  %               row5Keys/.style={
  %                 row5KeysLin,
  %                 ymode=log,
  %               },
  %             }
  %             %\endgroup
  %           %\endgroup
  %         }
  %         %\endgroup
  %       %\endgroup
  %     %\endgroup
  % 
  %     %\begingroup Diagrammi dei nodi
  %     \newcommand\nodeDiagramSettings{
  %       %\begingroup Dimensioni
  %         \pgfmathsetmacro\xPlotSep{10}
  %         \pgfmathsetmacro\yPlotSep{4}
  %         \pgfmathsetmacro\plotScale{.8}
  % 
  %         \pgfmathsetmacro\titleSize{11}
  %         \pgfmathsetmacro\commentSize{10}
  % 
  %         \pgfmathsetmacro\nodeSize{20pt} % Diametro del nodo
  %         \pgfmathsetmacro\starSize{2.25*\nodeSize}
  %         \pgfmathsetmacro\numberSize{.53*\nodeSize}
  %       %\endgroup
  % 
  %       %\begingroup Colori
  %         % \definecolor{cRed}{HTML}{D68490}
  %         % \definecolor{cPurple}{HTML}{7E6EAC}
  %         % \definecolor{cTeal}{HTML}{87BDB1}
  %         % \definecolor{cOrange}{HTML}{F2D0A4}
  %         % \definecolor{cEdgeLight}{HTML}{EEE8D5} 
  %         % \definecolor{cEdgeDark}{HTML}{998888}
  %         % \definecolor{cStar}{HTML}{FDF6E3}
  %         %
  %         \definecolor{color1}{HTML}{800000} % Deep Red
  %         \definecolor{color2}{HTML}{4682B4} % Steel Blue
  %         \definecolor{color3}{HTML}{6B8E23} % Olive Green
  %         \definecolor{color4}{HTML}{708090} % Slate Gray
  %         \definecolor{color5}{HTML}{BC8F8F} % Rosy Brown
  %         \definecolor{color6}{HTML}{008B8B} % Dark Cyan
  %         \definecolor{color7}{HTML}{B8860B} % Dark Goldenrod
  %       %\endgroup
  % 
  %       %\begingroup Testo
  %         \tikzset{
  %           titleStyle/.style={
  %             font=\textsize{\titleSize},
  %             align=center,
  %             anchor=south,
  %           },
  %           commentStyle/.style={font=\textsize{\commentSize}},
  %           arrowCommentStyle/.style={
  %             midway,
  %             commentStyle,
  %             left,
  %             sloped,
  %             below
  %           },
  %           numberStyle/.style={font=\textsize{\numberSize}}
  %         }
  %       %\endgroup
  % 
  %       %\begingroup Frecce
  %         \tikzset{
  %           transformingArrow/.style={
  %             % =Latex,
  %             ->,
  %             >=Stealth,
  %             dashed,
  %             line width=.8pt,
  %             black!90,
  %             shorten >=5pt,
  %             shorten <=10pt
  %           },
  %           indirectArrow/.style={
  %             draw=black!95,
  %             line width=.5pt,
  %           },
  %           nodeArrowExact/.style={
  %             indirectArrow,
  %             % ->,
  %             % -{Stealth[scale=1.2]},
  %             -{Triangle[scale=.75]},
  %             % draw=cEdgeLight,
  %             % draw=red!50,
  %             % shorten >=2pt,
  %             % shorten <=2pt
  %           },
  %           nodeArrowApprox/.style={
  %             indirectArrow,
  %             -,
  %             dashed,
  %           },
  %           transparentNodeArrow/.style={
  %             opacity=.4,
  %           },
  %           arrowAxis/.style={
  %             ->,
  %             line width=.3pt,
  %             >=Latex,
  %           },
  %         }
  %       %\endgroup
  % 
  %       %\begingroup Nodi
  %         \tikzset{
  %           nodeStyle/.style={
  %             circle,
  %             minimum size=\nodeSize,
  %             % draw=black!80,
  %             draw=black,
  %             line width=.5pt,
  %             inner sep=0pt,
  %           },
  %           starStyle/.style={
  %             star,
  %             star points=5,
  %             star point ratio=2,
  %             % fill=cStar,
  %             fill=black!50,
  %             fill opacity=.25,
  %             minimum size=\starSize,
  %             draw=none
  %           },
  %         }
  %       %\endgroup
  %     }
  %     %\endgroup
  % 
  %     %\begingroup GNU-Plot
  %       % \IfFileExists{\sourceauxpath.aux}{
  %       %   \externaldocument{\sourceauxpath}
  %       % }{\relax}
  %       
  %       \pgfkeys{/pgf/plot/gnuplot call={cd /d "\sourceoutdir" && gnuplot}}
  %     %\endgroup
  %   %\endgroup
  % %\endgroup
  % 
  % %\begingroup For the paper
  %   \def\tsc#1{\csdef{#1}{\textsc{\lowercase{#1}}\xspace}}
  %   \tsc{WGM}
  %   \tsc{QE}
  %   
  %   \newcommand\etal{et al.}
  %   \newcommand\KTMAS{\hyperlink{acrKineticTheoryMultiAgentSystems}{KTMAS}}
  %   
  %   % \newcommand\MathSpacing[3]{
  %   %   \thinmuskip = #1mu\relax
  %   %   \medmuskip = #2mu\relax
  %   %   \thickmuskip = #3mu\relax
  %   % }
  %   % \MathSpacing123
  %   
  %   \newcommand\RPlus{\mathbb R_+}
  %   \newcommand\NPlus{\mathbb N_+}
  %   \DeclareMathAlphabet\pazocal{OMS}{zplm}{m}{n}
  %   
  %   \DeclareMathOperator\Lognormal{Lognormal}
  %   \DeclareMathOperator\BiLognormal{BiLognormal}
  %   \DeclareMathOperator\Bernoulli{Bernoulli}
  %   \DeclareMathOperator\Exponential{Exp}
  %   \DeclareMathOperator\Student{Student}
  %   \DeclareMathOperator\Variance{var}
  %   \DeclareMathOperator\IC{IC}
  %   \DeclareMathOperator\Histogram{\mathcal B}
  %   \DeclareMathOperator\CCDF{\mathcal R}
  %     
  %   \newcommand\ThreeQuarters{{3\hspace{-0.15em}/\hspace{-0.15em}4}}
  %   \newcommand\OneQuarter{{1\hspace{-0.15em}/\hspace{-0.15em}4}}
  %   
  %   \newcommand\de[1][]{\mathrm{d}#1}
  %   \newcommand\derS[3][d]{\frac{\mathrm{#1}#2}{\mathrm{#1}#3}}
  %   \newcommand\derP[2]{\frac{\partial#1}{\partial#2}}
  %   \newcommand\underhat[1]{\underaccent{\hat}{#1}}
  %   \renewcommand\undertilde[1]{\underaccent{\tilde}{#1}}
  %   
  %   \makeatletter
  %   \newcommand\eqtag{\@ifstar{\eqtag@star}{\eqtag@nostar}}
  %   
  %   \newcounter{eqtag}
  %   \newcommand\eqtag@nostar[2]{%
  %     \begin{equation*}%
  %       \refstepcounter{eqtag}%
  %       % \protected@edef\@currentlabeltype{eqtag}%
  %       \protected@edef\@currentlabel{#1}%
  %       #2%
  %       \tag*{#1}%
  %     \end{equation*}%
  %   }
  %   \newcommand\eqtag@star[2]{%
  %     \begin{equation*}%
  %       \refstepcounter{eqtag}%
  %       % \protected@edef\@currentlabeltype{eqtag}%
  %       \protected@edef\@currentlabel{#1}%
  %       \let\label\ltx@label%
  %       #2%
  %       % \notag%
  %     \end{equation*}%
  %   }
  %   \makeatother
  %   
  %   \newcommand\crefsubfigInt[4][-]{Figures~\ref{#2}#3#1\ref{#2}#4}
  %   \newcommand\crefsubfigSep[4]{Figures~\ref{#1}#2\foreach \l in#3{,~\ref{#1}\l}~and~\ref{#1}#4}
  %   
  %   \crefname{section}{§}{§§}
  %   \crefrangeformat{section}{§§~#3#1#4--#5#2#6}
  %   \crefmultiformat{section}{§§~#2#1#3}{ and~#2#1#3}{, #2#1#3}{ and~#2#1#3}
  %   \crefrangemultiformat{section}{§§~#3#1#4--#5#2#6}{ and~#3#1#4--#5#2#6}{, #3#1#4--#5#2#6}{ and~#3#1#4--#5#2#6}
  %   
  %   \crefname{equation}{}{}
  %   \crefrangeformat{equation}{(#3#1#4--#5#2#6)}
  %   \crefmultiformat{equation}{(#2#1#3}{,~#2#1#3)}{,~#2#1#3}{,~#2#1#3)}
  %   \crefrangemultiformat{equation}{(#3#1#4--#5#2#6}{,~#3#1#4--#5#2#6)}{,~#3#1#4--#5#2#6}{,~#3#1#4--#5#2#6)}
  %   
  %   \crefname{figure}{Figure}{Figures}
  %   \crefrangeformat{figure}{Figures~#3#1#4--#5#2#6}
  %   \crefmultiformat{figure}{Figures~#2#1#3}{ and~#2#1#3}{, #2#1#3}{ and~#2#1#3}
  %   \crefrangemultiformat{figure}{Figures~#3#1#4--#5#2#6}{ and~#3#1#4--#5#2#6}{, #3#1#4--#5#2#6}{ and~#3#1#4--#5#2#6}
  %   
  %   \crefname{table}{Table}{Tables}
  %   \crefrangeformat{table}{Tables~#3#1#4--#5#2#6}
  %   \crefmultiformat{table}{Tables~#2#1#3}{ and~#2#1#3}{, #2#1#3}{ and~#2#1#3}
  %   \crefrangemultiformat{table}{Tables~#3#1#4--#5#2#6}{ and~#3#1#4--#5#2#6}{, #3#1#4--#5#2#6}{ and~#3#1#4--#5#2#6}
  %   
  %   \crefname{algorithm}{Algorithm}{Algorithms}
  %   \crefrangeformat{algorithm}{Algorithms~#3#1#4--#5#2#6}
  %   \crefmultiformat{algorithm}{Algorithms~#2#1#3}{ and~#2#1#3}{, #2#1#3}{ and~#2#1#3}
  %   \crefrangemultiformat{algorithm}{Algorithms~#3#1#4--#5#2#6}{ and~#3#1#4--#5#2#6}{, #3#1#4--#5#2#6}{ and~#3#1#4--#5#2#6}
  %   
  %   %\newtheorem{theorem}{Theorem}
  %   %\newtheorem{lemma}[theorem]{Lemma}
  %   \newdefinition{dfn}{Definition}
  %   \newdefinition{rmk}{Remark}
  %   \newdefinition{hyp}{Hypotesis}
  %   %\newproof{pf}{Proof}
  %   %\newproof{pot}{Proof of Theorem \ref{thm}}
  %   
  %   \crefname{dfn}{Definition}{Definitions}
  %   \crefrangeformat{dfn}{Definitions~#3#1#4--#5#2#6}
  %   \crefmultiformat{dfn}{Definitions~#2#1#3}{ and~#2#1#3}{, #2#1#3}{ and~#2#1#3}
  %   \crefrangemultiformat{dfn}{Definitions~#3#1#4--#5#2#6}{ and~#3#1#4--#5#2#6}{, #3#1#4--#5#2#6}{ and~#3#1#4--#5#2#6}
  %   \crefname{hyp}{Hypothesis}{Hypotheses}
  %   
  %   \crefrangeformat{hyp}{Hypotheses~#3#1#4--#5#2#6}
  %   \crefmultiformat{hyp}{Hypotheses~#2#1#3}{ and~#2#1#3}{, #2#1#3}{ and~#2#1#3}
  %   \crefrangemultiformat{hyp}{Hypotheses~#3#1#4--#5#2#6}{ and~#3#1#4--#5#2#6}{, #3#1#4--#5#2#6}{ and~#3#1#4--#5#2#6}
  %   
  %   \crefname{rmk}{Remark}{Remarks}
  %   \crefrangeformat{rmk}{Remarks~#3#1#4--#5#2#6}
  %   \crefmultiformat{rmk}{Remarks~#2#1#3}{ and~#2#1#3}{, #2#1#3}{ and~#2#1#3}
  %   \crefrangemultiformat{rmk}{Remarks~#3#1#4--#5#2#6}{ and~#3#1#4--#5#2#6}{, #3#1#4--#5#2#6}{ and~#3#1#4--#5#2#6}
  %   
  %   % \newcommand\tikzfigure[5][0em]{%
  %   %   \begin{figure}[%
  %   %     pos=#2,%
  %   %     width=#4,%
  %   %     align=\centering%
  %   %   ]%
  %   %     \centering%
  %   %     \vspace*{#1}%
  %   %     \IfSubStr{#2}{p}{%
  %   %       \leavevmode\stepcounter{figposcount}%
  %   %       \zsaveposx{figpos-\thefigposcount}%
  %   %       \makebox[0pt][l]{%
  %   %         \hspace*{-\dimexpr\zposx{figpos-\thefigposcount}sp\relax}%
  %   %         \hspace*{\dimexpr0.5\paperwidth\relax}%
  %   %         \makebox[0pt][c]{%
  %   %           \scalebox{0.86}{\includegraphics{figs/#5.pdf}}%
  %   %         }%
  %   %       }
  %   %     }{%
  %   %       \includegraphics{figs/#5.pdf}%
  %   %     }%
  %   %     \caption{%
  %   %       \hyphenpenalty=0%
  %   %       \exhyphenpenalty=0%
  %   %       #3%
  %   %     }%
  %   %     \protect\label{#5}%
  %   %   \end{figure}%
  %   % }
  % %\endgroup
%\endgroup


% !TEX root = ../../Esperimenti/.tex/MEF.tex

% TikZFigureExternalized
\newcommand\compilerpath{../Altro/tikz2pdf.tex}
\newcommand\pdfdirectory{../../Figure/.fnl/Articolo/TikZpdf/}
\newcommand\auxpath{../Altro/mainCopy}
\newcommand\outdirectory\pdfdirectory

\makeatletter

% \tikzexternalizedfigure[#1]{#2}[#3]{#4}{#5}{#6}{#7}
% #1 -> Tipo di compilatore (di base è «pdflatex»)
% #2 -> Posizione della figura flottante
% #3 -> Tralazione verticale (di base è nulla)
% #4 -> Percorso assoluto o relativo della filza .tex da compilare esternamente
% #5 -> Didascalia
% #6 -> Lunghezza della didascalia
% #7 -> Nome della filza «.pdf» della figura tikz compilata

% https://tex.stackexchange.com/questions/57788/understanding-ifnextchar
\def\tikzexternalizedfigure{%
  \@ifnextchar[{\tikz@ext@figure@compiler}{\tikz@ext@figure@compiler[pdflatex]}
}
\def\tikz@ext@figure@compiler[#1]#2{%
  \@ifnextchar[{\tikz@ext@figure@translation[#1]#2}{\tikz@ext@figure@translation[#1]#2[0em]}
}
\def\tikz@ext@figure@translation[#1]#2[#3]#4#5#6#7{%
  % \newcommand\outdirectory{\pdfdirectory#6.}
  \begin{figure}[%
   pos=#2,%
   width=#6,%
   align=\centering%
  ]%
    \centering%
    % Recupera la pagina dall'etichetta (default a 1 se non esiste)
    \checkoddpage%
    \Ifthispageodd{\def\pagetype{odd}}{\def\pagetype{even}}%
    %
    % Alternativa con «\getpagerefnumber»
    % \edef\currentfigpage{\getpagerefnumber{#7}}
    % \ifx\currentfigpage\relax\def\currentfigpage{1} \fi
    % \pgfmathparse{int(mod(\currentfigpage,2) == 1 ? 1 : 0)}
    % \ifnum\pgfmathresult=1\def\pagetype{odd}
    % \else\def\pagetype{even}\fi
    %
    \def\shlcmd{%
      \immediate\ShellEscape{%
        #1 %
        --shell-escape %
        -interaction=batchmode %
        -output-directory="\pdfdirectory" %
        -jobname="#7" %
        "\string\def\string\isodd{\pagetype}%
        \string\def\string\sourcefilepath{#4}%
        \string\def\string\sourceoutdir{\outdirectory}%
        \string\def\string\sourceauxpath{\auxpath}%
        \string\input{\compilerpath}"%
      }%
    }%
    \IfSubStr{#2}{p}{% 
      \leavevmode\stepcounter{figposcount}\zsaveposx{figpos-\thefigposcount}%
      \makebox[0pt][l]{%
        \hspace*{-\dimexpr\zposx{figpos-\thefigposcount}sp\relax}%
        % \hspace*{\dimexpr0.5\paperwidth\relax}%
        \def\hsep{0cm}\Ifthispageodd{\def\offset{+\hsep}}{\def\offset{-\hsep}}% per far toccare il bordo sinistro/destro dei grafici al margine interno/esterno si trasla di -/+1.9 cm
        \hspace*{\dimexpr0.5\paperwidth+\offset\relax}%
        \makebox[0pt][c]{%
          \begin{minipage}\paperwidth%
            \centering%
            \vspace*{#3}%
            \addtocounter{figure}{1}%
            \setcounter{subfigure}{0}%
            \IfFileExists{\pdfdirectory#7.pdf}{%
              \ifbool{forcetikzrecompile}{\shlcmd}{\relax}%
            }{\shlcmd}%
            \scalebox{0.86}{%
              \includegraphics{\pdfdirectory#7.pdf}%
            }%
            \addtocounter{figure}{-1}%
            \newcaption{%
              \hyphenpenalty=0%
              \exhyphenpenalty=0%
              #5%
            }{#6}%
            \protect\label{#7}%
          \end{minipage}%
        }%
      }%
    }{%
      \makebox[\textwidth][c]{%
        \begin{minipage}\linewidth
          \centering%
          \addtocounter{figure}{1}\setcounter{subfigure}{0}%
          \IfFileExists{\pdfdirectory#7.pdf}{%
            \ifbool{forcetikzrecompile}{\shlcmd}{\relax}
          }{\shlcmd}%
          \includegraphics{\pdfdirectory#7.pdf}%
          \addtocounter{figure}{-1}%
          \newcaption{%
            \hyphenpenalty=0%
            \exhyphenpenalty=0%
            #5%
          }{#6}%
          \protect\label{#7}%
        \end{minipage}%
      }%
    }%
  \end{figure}%
}
\makeatother

%\begingroup Vecchia versione con un solo argomento facoltativo
  % \newcommand\tikzexternalizedfigure[6][0em]{%
  %   % \newcommand\outdirectory{\pdfdirectory#6.}
  %   \begin{figure}[#2]%
  %     \centering%
  %     % Recupera la pagina dall'etichetta (default a 1 se non esiste)
  %     \checkoddpage
  %     \Ifthispageodd{\def\pagetype{odd}}{\def\pagetype{even}}
  %     %
  %     % Alternativa con «\getpagerefnumber»
  %     % \edef\currentfigpage{\getpagerefnumber{#6}}
  %     % \ifx\currentfigpage\relax\def\currentfigpage{1} \fi
  %     % \pgfmathparse{int(mod(\currentfigpage,2) == 1 ? 1 : 0)}
  %     % \ifnum\pgfmathresult=1\def\pagetype{odd}
  %     % \else\def\pagetype{even}\fi
  %     %
  %     \edef\shlcmd{
  %       \immediate\ShellEscape{%
  %         lualatex %
  %         --shell-escape %
  %         -interaction=batchmode %
  %         -output-directory="\pdfdirectory" %
  %         -jobname="#6" %
  %         "\string\def\string\isodd{\pagetype}%
  %         \string\def\string\sourcefilepath{#3}%
  %         \string\def\string\sourceoutdir{\outdirectory}%
  %         \string\def\string\sourceauxpath{\auxpath}%
  %         \string\input{\compilerpath}"%
  %       }
  %     }
  %     \IfSubStr{#2}{p}{% 
  %       \leavevmode\stepcounter{figposcount}\zsaveposx{figpos-\thefigposcount}%
  %       \makebox[0pt][l]{%
  %         \hspace*{-\dimexpr\zposx{figpos-\thefigposcount}sp\relax}%
  %         \hspace*{0.5\paperwidth}%
  %         \makebox[0pt][c]{%
  %           \begin{minipage}{\paperwidth}%
  %             \vspace*{#1}\centering
  %             \addtocounter{figure}{1}
  %             \setcounter{subfigure}{0}%
  %             \IfFileExists{\pdfdirectory#6.pdf}{%
  %               \ifbool{forcetikzrecompile}{\shlcmd}{\relax}
  %             }{\shlcmd}%
  %             \includegraphics{\pdfdirectory#6.pdf}%
  %             \addtocounter{figure}{-1}
  %             \didascalia{#4}{#5}
  %             \protect\label{#6}
  %           \end{minipage}
  %         }%
  %       }%
  %     }{%
  %       \makebox[\textwidth][c]{
  %         \begin{minipage}{\linewidth}
  %           \centering
  %           \addtocounter{figure}{1}\setcounter{subfigure}{0}
  %           \IfFileExists{\pdfdirectory#6.pdf}{%
  %             \ifbool{forcetikzrecompile}{\shlcmd}{\relax}
  %           }{\shlcmd}%
  %           \includegraphics{\pdfdirectory#6.pdf}%
  %           \addtocounter{figure}{-1}
  %           \didascalia{#4}{#5}
  %           \protect\label{#6}
  %         \end{minipage}%
  %       }%
  %     }%
  %   \end{figure}%
  % }
%\endgroup


\newcommand{\N}{\mathbb{N}}
\newcommand{\R}{\mathbb{R}}

\setbool{forcetikzrecompile}{false}
% \setbool{forcetikzrecompile}{true}

\allowdisplaybreaks

\title{Modelling population distribution across cities on spatial networks using the kinetic theory of multi-agent systems}
\author{Valerio Taralli}
\author{Andrea Tosin}
\affil{{\small Department of Mathematical Sciences ``G. L. Lagrange'', Politecnico di Torino, Italy}}
\date{}
		
\begin{document}
\maketitle
	
\begin{abstract}
We investigate the urban population distribution of the Italian region of Sardinia through the kinetic theory of multi-agent systems on spatial networks, interpreting cities both as agents and vertices of a graph, and representing interurban links as edges. We consider various emigration rules to model the interactions between agents, obtaining two main results: the bimodal lognormal fit of the total population distribution among cities and the Pareto fit of the tail. Overall, our model provides results consistent with the actual distribution derived from ISTAT data of Sardinia; furthermore, the bimodal lognormal fit generally follows the distribution tail more accurately than the Pareto one, although the two coincide in the first half.

\medskip

\noindent{\bf Keywords:} multi-agent systems, Boltzmann-type equations on graphs, spatial networks, population distribution, Pareto distribution

\medskip

\noindent{\bf Mathematics Subject Classification:} 82C40, 90B10, 91D25, 05C90
\end{abstract}

\section{Introduction}
\label{secIntroduction}
%
Population distribution across cities has been sporadically analysed in the past. The first was Auerbach \cite{Auerbach1913} who, at the beginning of the 20th century, noted a feature later formalized by Zipf\footnote{Actually, as Zipf himself admits in~\cite{Zipf2016}., he was not the first to discover such a law.} \cite{Zipf2016} almost fifty years later, albeit initially in a linguistic context. Zipf's empirical law states the following relationship between the frequency $\nu$ of a word in a language and its rank $r$ within the most frequent words of that language:
\begin{equation}
    \nu\propto\frac{1}{r}.
    \label{eqZipfsLaw}
\end{equation}
To give a concrete example, this means that there exists a constant $C>0$ such that if the tenth most frequent word appears $2000$ times, then $2000\approx C/10$. The same law described in~\eqref{eqZipfsLaw} can be found in many other contexts, including the distribution of population across cities if one replaces $\nu$ with the \textit{population size} $s$~\cite[Fig. 9-2]{Zipf2016}. However, this is true only for $s\gg 1$, namely in the tail of the distribution, contrary to Zipf's original linguistic context, where~\eqref{eqZipfsLaw} is instead expected to hold always.
        
More generally, Zipf's law can be described by a Pareto distribution: let $S\in\R_+$ be the random variable describing a city size and let $f_S=f_S(s)$ be its probability distribution, which allows one to interpret $f_S(s)\,ds$ as the fraction of cities with $s\in [s,\,s+ds]$. Then, $S$ is said to follow a Pareto distribution in the tail if
\begin{equation}
    \int_{s}^{+\infty}f_S(z)\,dz \approx \frac{c}{s^\beta}, \quad s\gg 1,
    \label{eqComplementaryCumulativePareto}
\end{equation}
where $\beta,\,c\in\R_+$ are parameters. In particular, the exponent $\beta$ is called the Pareto index. For $\beta=1$,~\eqref{eqComplementaryCumulativePareto} is analogous to~\eqref{eqZipfsLaw} with $s$ playing the role of $r$.
        
More recently, it has been discovered~\cite{Eeckhout2004,Gualandi2019HBLD} that, apart from the tail, $f_S$ is well fitted by the probability density function of a lognormal random variable $X\sim\Lognormal(\mu,\,\sigma)$, specifically
\begin{equation}
    f_X(x)=\frac1{x\sigma\sqrt{2\pi}}\exp\!\left(-\frac{(\log{x}-\mu)^2}{2\sigma^2}\right).
    \label{eqLognormalDistribution}
\end{equation}
As the name suggests, the peculiarity of this distribution is that $\log{X}\sim\mathcal{N}(\mu,\,\sigma)$, where $\mathcal{N}(\mu,\,\sigma)$ denotes the Gaussian distribution with mean $\mu$ and standard deviation $\sigma$.
        
Furthermore, Gualandi and coauthors~\cite{Gualandi2019HBLD,Gualandi2019SDC} realised that the entire distribution is well fitted by a bimodal lognormal (here also denoted as bilognormal), namely a convex combination of~\eqref{eqLognormalDistribution}:
\begin{equation}
    f_S(s)\approx \zeta f_{X_1}(s)+(1-\zeta)f_{X_2}(s),
    \label{eqBimodalLognormalDistribution}
\end{equation}
with $\zeta\in(0,\,1)$, $X_i\sim\Lognormal(\mu_i,\,\sigma_i)$ for $i=1,\,2$, whence $S\approx\BiLognormal(\zeta,\,\mu_1,\,\sigma_1,\,\mu_2,\,\sigma_2)$, cf. Figure~\ref{figLognormalFunctionExample}.
        
Building on these preliminary observations, in this paper we aim to model, by means of the Kinetic Theory of Multi-Agent Systems (KTMAS), the interactions between cities connected on a spatial graph, thereby attempting to reproduce a population distribution with characteristics analogous to those actually observed in the Italian region of Sardinia -- mainly bilognormal body and Pareto tail. Upon achieving such a goal, we interpret and explain the migration phenomenon through the emigration laws at the basis of the interactions between cities in our kinetic model.
        
In the present work, we regard cities as agents characterized by their population size $s\in\R_+$, and assume that they interact, i.e. they exchange individuals, through an underlying spatial graph structure. Except for the network, such a description is analogous to the classical one of gas molecules characterized by their position and velocity in the kinetic theory of gases. As a matter of fact, in the literature, at least to the authors' knowledge, there are no contributions dealing simultaneously with KTMAS, population distribution across cities, and graphs, despite the natural modelling representation that these elements provide. More precisely:
\begin{itemize}
\item some authors either apply the KTMAS to cities without considering any underlying graph topology~\cite{Gualandi2019SDC} or, conversely, consider the graph structure without applying it~\cite{Berliant2018};
\item other authors view graph vertices as places where agents live and interact, thus creating a large system of coupled Boltzmann-type equations~\cite{Loy2021};
\item the contribution closest to the present paper applies the aforesaid theory to the context of social networks~\cite{Nurisso2024}.
\end{itemize}

This gap is probably due to the fact that developments in graph theory applied to the KTMAS have been quite recent and primarily focused on the vertex-place perspective rather than the vertex-agent one. Therefore, the primary contribution of this work is showing that the second perspective is as valid as the first, despite being less common. Therefore, this paper presents a formal application-oriented exploration of the vertex-agent perspective; for preliminary analytical investigations, we refer the interested reader to~\cite{taricco2026PREPRINT}.

Overall, the paper is organised as follows. The second section summarises the essential notions of the graph-mediated KTMAS and contextualises them within the problem of population distribution across cities. The third section presents numerical results for various emigration rules applied to the Italian region of Sardinia. Finally, the fourth section draws conclusions and outlines potential directions for future research.

\section{Outline of the KTMAS on graphs}
\label{secNotesKTMAS}
%
In this section, we briefly discuss some assumptions on the type of graph used and the two Boltzmann-type equations on graphs that we use to formalise our mathematical model. For the general principles of KTMAS, we refer the interested reader to~\cite{Loy2025EKTMAS,PareschiIMS2013}.

\subsection{Networks and cities}
\label{secNetworksAndCities}
%
We begin by recalling the definition of graph\footnote{Other related concepts are instead taken for granted. They can be found in any reference about graph theory, such as e.g.,~\cite{Barthelemy2011}.}:
            
\begin{dfn}[Graph] \label{defGraph}
A \emph{graph} is a set $\mathcal G:=(\mathcal I,\,\mathcal E)$, where $\mathcal I\subseteq\N_+$ is the set of vertex indices and $\mathcal E$ is the set of edges, which are pairs of indices in $\mathcal I$. Two vertices $i,\,j\in\mathcal I$ are connected if and only if $(i,\,j)\in\mathcal E\subseteq\N_+\times\N_+$.
\end{dfn}
            
Next, we discuss how to interpret a link between two vertices, viz. cities, given the multitude of possible urban connections~\cite{Barthelemy2011}. In fact, a city can be represented at many levels by means of a graph:
\begin{enumerate*}[label=(\roman*)]
\item \textit{intraurban} level, representing either streets as edges and their intersections as vertices or transport networks (tram, bus, and metro);
\item \textit{interurban} level, taking individually into consideration various transport networks (railway, naval, air, \dots).
\end{enumerate*}
Recalling the introductory section, it is clear that in order to predict the population distribution across cities the following observations hold: on the one hand, intraurban representations should be discarded, because in that context people's movement take place within a single city; on the other hand, interurban representations should be considered simultaneously and not individually, since each individual may choose different transportation modes to move.
    
This is why the graph we refer to in our model is the one related to commuting movements across cities~\cite{Barthelemy2011,DeMontis2007}, which holistically shows all their possible connections regardless of the chosen transportation mode. Moreover, it shows realistic connections associated with daily rather than extraordinary movements, such as e.g., holidays, medical visits, \dots. Upon choosing such a representation, we need to establish the type of graph we deal with:
\begin{hyp} \label{hypGraphType}
The graph $\mathcal G$ is \emph{undirected}, hence its adjacency matrix is symmetric.
\end{hyp}

A comment on Assumption~\ref{hypGraphType} is in order. A city may interact, viz. exchange population, with another one but not vice versa: this means that an interaction between two cities may not be symmetric. On the other hand, if a city can interact with another one, then the opposite is always possible, because the two cities are connected: this means that the structure of the graph connections is undirected, i.e. symmetric.

In addition to this, we are only interested in predicting the distribution of the population with respect to an \textit{a priori} fixed graph topology, which leads to the following further assumption:
\begin{hyp} \label{hypStaticGraph}
The graph $\mathcal G$ is \emph{non-coevolving}, i.e. the sets $\mathcal I$ and $\mathcal E$ do not vary in time.
\end{hyp}
Clearly, this is a strong simplification, since in reality the graph topology, especially the interurban connections in $\mathcal E$, clearly co-evolves following the development of the cities themselves.
            
Finally, we observe that this type of graph is spatial~\cite{Barthelemy2011,DeMontis2007}, meaning that its vertices are points in an Euclidean space, even though, for the moment, this is a somewhat formal characteristic. However, it will be useful to keep it in mind when interpreting certain emigration rules in Section~\ref{secInterpretations}. Moreover, contrary to what one might think at first glance, our graph is not scale-free \cite{Barabasi1999}, precisely because of its spatial nature~\cite{Barthelemy2011,Broido2019}. This, however, does not rule out the existence of vertices that are more central than others, but only implies that the maximum degree of a vertex is bounded above due to physical constraints.
    
\subsection{Boltzmann-type equation on the exact graph topology}
\label{secExactBoltzmannType}
%
The KTMAS is centred around agents interacting with one another, thereby changing their microscopic state. However, usually the precise knowledge of the microscopic state of a single agent is neither useful nor necessary, which is why a statistical approach is taken by regarding microscopic states as time-evolving random variables forming a stochastic process. To derive an evolution equation for the probability distribution, say $f$, of the microscopic state of a representative agent, four assumptions are formulated, cf.~\cite{Loy2025EKTMAS}:
\begin{hyp} \label{hypBoltzmannType}
\begin{enumerate}[before=\leavevmode,label=(\roman*)]
\item\label{hypBoltzmannTypeHomogeneousDensity} The probability distribution $f$ is homogeneous in space.
\item\label{hypBoltzmannTypeBinaryInteractions} Non-pairwise interactions between agents are neglected.
\item\label{hypBoltzmannTypeIndistinguishableAgents} Agents are indistinguishable, and each of them represents the ensemble.
\item\label{hypBoltzmannTypeMolecularChaos} Two interacting agents are statistically independent at the moment of the interaction.
\end{enumerate}
\end{hyp}

Apparently, an inconsistency arises from Assumption~\ref{hypBoltzmannType}\ref{hypBoltzmannTypeIndistinguishableAgents} if agents are confused with graph vertices: the former are indistinguishable, but the underlying network makes the latter distinct due to the indices in $\mathcal I$. To bypass such a difficulty of the theory, it is sufficient to include the vertex index in the agent's microscopic state, considering it as a random variable. As a result, the microscopic state of a generic agent of the system is the pair $(I,\,S_t)$, where $S_t\in\R_+$ is the random variable expressing the population size at time $t\geq 0$ of a representative agent, viz. city, of the system. Due to~Assumption~\ref{hypStaticGraph}, the random variable $I$ is constant in time.

\begin{rmk} \label{rmkPopulationNumericalNature}
It is natural that fractions of people do not exist, hence in principle $S_t\in\mathbb N$. Nevertheless, the most realistic choice is not always the easiest modelling-wise. At any rate, if necessary, it is possible to gauge the actual population size by rounding it to the nearest natural number.
\end{rmk}

\tikzexternalizedfigure{thb}
  {../../Figure/.fnl/Elaborato/Capitolo3/RappresentazioneGraficaTCSMAReti.tex}
  {Graphical representation of the KTMAS on graphs}
  {\linewidth}{figGraphicRepresentationKTMASNetworks}

Reasoning in this way, we describe statistically the microscopic state $(I,\,S_t)$ of a representative city with the joint probability distribution $f:\mathcal I\times\R_+\times [0,\,+\infty)\to\R_+$, which is discrete in the first variable $i\in\mathcal I$ and continuous in the second variable $s\in\R_+$, and takes therefore the form
\begin{equation}
    f(i,s,t):=\frac{1}{N}\sum_{j\in\mathcal I}f_j(s,t)\otimes\delta(i-j),
    \label{eqUrbanJointDensity}
\end{equation}
where $N:=\vert\mathcal I\vert$ is the total number of cities in the graph, while $\delta(\cdot)$ denotes the Dirac delta centred at the origin. On the other hand, $f_i:\R_+\times [0,\,+\infty)\to\R_+$ is the probability distribution of $S_t$ conditioned to the $i$-th city, cf. Figure~\ref{figGraphicRepresentationKTMASNetworks}. This way, $f$ integrates to $1$ over its domain:
$$ \int_{\mathcal I}\int_{\R_+}f(i,s,t)\,ds\,di=\frac{1}{N}\sum_{i\in\mathcal I}\int_{\R_+}f_i(s,t)\,ds=1, \qquad \forall\,t\geq 0. $$
            
When two cities interact they exchange population, thereby altering their respective microscopic states. This is modelled by the following general interaction rules:
\begin{equation}
    S'_t=\psi(I,S_t,I_\ast,S_{\ast,t}), \qquad S_{\ast,t}'=\psi_\ast(I,S_t,I_\ast,S_{\ast,t}),
    \label{eqUrbanInteractionRules}
\end{equation}
where the functions $\psi,\,\psi_\ast$ describe the post-interaction microscopic states. However, two randomly picked cities can interact only if they are connected by an edge of the graph $\mathcal G$. This can be taken into account via the Bernoulli random variable
\begin{equation}
    \Theta\sim\Bernoulli(a_{I,I_\ast}\Delta t), \qquad
        a_{I,I_\ast}:=
        \begin{cases}
            1 & \text{if } (I,\,I_\ast)\in\mathcal E \\
            0 & \text{otherwise},
        \end{cases}
    \label{eqUrbanInteractionBernoulli}
\end{equation}
where $a_{I,I_\ast}$ is the element of the adjacency matrix of $\mathcal G$ corresponding to the pair of indices $(I,\,I_\ast)$ and $\Delta{t}>0$ is a discrete time step corresponding to the characteristic duration of an interaction. For consistency, we require $\Delta{t}\leq 1$, so that the parameter of the random variable $\Theta$ is well defined.

Combining~\eqref{eqUrbanInteractionRules} and~\eqref{eqUrbanInteractionBernoulli}, the following agent-based interaction algorithm is obtained:
\begin{equation}
    \begin{cases}
        S_{t+\Delta t}=(1-\Theta)S_t+\Theta S'_t \\
        S_{\ast,t+\Delta t}=(1-\Theta)S_{\ast,t}+\Theta S'_{\ast,t},
    \end{cases}
    \label{eqUrbanInteractionAlgorithm}
\end{equation}
consistently with the fact that, in an urban context, if a city exchange population with another both change their population size.

Let now $\Phi:\mathcal I\times\R_+\to\R$ be an arbitrary observable quantity. By averaging~\eqref{eqUrbanInteractionAlgorithm} over all possible observable quantities and taking the continuous-time limit $\Delta{t}\to 0^+$, one obtains, cf.~\cite{Nurisso2024}:
\begin{multline}
    \frac{d}{dt}\int_{\mathcal I}\int_{\R_+}\Phi(i,s)f_i(s,t)\,ds \\
    =\frac{1}{N}\sum_{i,i_\ast\in\mathcal I}a_{i,i_\ast}\iint_{\R_+^2}\frac{\Phi(i,s')+\Phi(i_\ast,s_\ast')-\Phi(i,s)-\Phi(i_\ast,s_\ast)}{2}
        f_i(s,t)f_{i_\ast}(s_\ast,t)\,ds\,ds_\ast,
    \label{eqUrbanAsymmetricHomogeneousWeakFormBoltzmannType}
\end{multline}
which is a system of Boltzmann-type equations in weak form for the in-vertex probability distributions $\{f_i\}_{i=1}^{N}$, based on the exact topology of the connections of $\mathcal G$ as encoded in the $N\times N$ adjacency matrix $\mathbf{A}:=(a_{i,i_\ast})_{i,i_\ast=1}^{N}$.

In~\eqref{eqUrbanAsymmetricHomogeneousWeakFormBoltzmannType}, $s',\,s_\ast'$ are the post-interaction population sizes of two interacting cities, which exchange inhabitants on the basis of the interaction rules~\eqref{eqUrbanInteractionRules}. In particular, they are realisations of the random variables $S_t',\,S_{\ast,t}'$.

\subsection{Boltzmann-type equation on an approximated graph topology}
\label{secApproximatedBoltzmannType}
%
In~\cite{Nurisso2024}, an approximation of~\eqref{eqUrbanAsymmetricHomogeneousWeakFormBoltzmannType} has been proposed based on a corresponding approximation of the adjacency matrix $\mathbf A$. Such an approximation is done so as to make it possible to pass to the \textit{dense graph limit} $N\to\infty$, thereby obtaining a single Boltzmann-type equation, which does not require the precise knowledge of the graph connections.

To this purpose, we recall that $\mathbf A$ is assumed to be symmetric, cf. Assumption~\ref{hypGraphType}, and we give the following
\begin{dfn} \label{defDegreeMatrix}
Let $\mathbf{1}:=(1,\,1,\,\dots,\,1)^\top\in\R^N$, and define
$$ \boldsymbol{\xi}:=\mathbf{A}\mathbf{1}, $$
the vector of the \emph{degrees} of the vertices of $\mathcal G$; i.e., $\xi_i:=\sum_{i_\ast\in\mathcal I}a_{i,i_\ast}\in\{0,\,\dots,\,N\}$ is the number of edges incident on vertex $i\in\mathcal I$, for $i=1,\,\dots,\,N$. Define moreover
$$ \mathbf{B}:=\frac{\boldsymbol{\xi}\boldsymbol{\xi}^\top}{d_N}, \qquad d_N:=\mathbf{1}^\top\mathbf{A}\mathbf{1}. $$
Then, since $\mathbf{B}\mathbf{1}=\mathbf{A}\mathbf{1}$, $\mathbf{B}\in\R^{N\times N}$ is a rank-one matrix with the same degrees as $\mathbf{A}$.
\end{dfn}

Replacing $\mathbf A$ with $\mathbf B$ corresponds to approximating $\mathcal G$ with a \textit{random graph} with the same degree sequence as $\mathcal G$, see~\cite{bollobas1980EJC,molloy1995RSA}. Using $\mathbf B$ instead of $\mathbf A$ in~\eqref{eqUrbanAsymmetricHomogeneousWeakFormBoltzmannType} and, at the same time, introducing the \textit{normalised} degree variable
$$ x:=\frac{\xi}{N}\in [0,\,1], $$
it can be shown, see~\cite{Nurisso2024,taricco2026PREPRINT}, that~\eqref{eqUrbanAsymmetricHomogeneousWeakFormBoltzmannType} converges, in the limit $N\to\infty$, to the following weak Boltzmann-type equation:
\begin{multline}
    \frac{d}{dt}\int_{\R_+}\varphi(s)g(x,s,t)\,ds \\
    =\int_0^1\iint_{\R_+^2}\frac{x x_\ast}{d}\cdot\frac{\varphi(s')+\varphi(s_\ast')-\varphi(s)-\varphi(s_\ast)}{2}
        g(x,s,t)g(x_\ast,s_\ast,t)\,ds\,ds_\ast\,dx_\ast,
    \label{eqRelaxedGraphBoltzmannType}
\end{multline}
where:
\begin{itemize}
\item $g$ is the limit probability distribution that $f$ in~\eqref{eqUrbanJointDensity} converges to, up to a suitable scaling, in the dense graph limit;
\item the observable $\Phi$ has been reduced to an arbitrary observable $\varphi:\R_+\to\R$, which depends only on the population size $s$;
\item the constant $d$ is the dense-graph-limit analogue of $d_N$ in Definition~\ref{defDegreeMatrix} and is defined as $d:=\int_0^1\int_{\R_+}xg(x,s,t)\,ds\,dx$, see~\cite{Nurisso2024}.
\end{itemize}

It is worth stressing that, in the limit, $x\in [0,\,1]$ plays the role of a (continuous) \textit{latent variable}, which accounts for the topology of the connections among the cities.

\tikzexternalizedfigure{htb}
  {../../Figure/.fnl/Articolo/Capitolo2/RappresentazioneGraficaApprossimazione.tex}
  {Graphical representation of the KTMAS approximation on graphs}
  {.85\linewidth}{figGraphicalRepresentationApproximation}

The implications of~\eqref{eqRelaxedGraphBoltzmannType} are profound: approximating $\mathbf A$ with $\mathbf B$ implies losing the distinguishability of the vertices of the graph, as it is well represented by the passage from the index $i$ to the latent variable $x$. This way, the graph topology is relaxed, cf. Figure~\ref{figGraphicalRepresentationApproximation}; specifically, it remains solely as a degree distribution. Such a reduction is akin to blurring the connection details of the vertices: one transitions from a specific graph to a class of random graphs sharing the same degree distribution. Therefore,~\eqref{eqRelaxedGraphBoltzmannType} is closer to the original statistical spirit of the KTMAS, and, as such, it is a good candidate to approximate~\eqref{eqUrbanAsymmetricHomogeneousWeakFormBoltzmannType}.

\subsection{Emigration rules}
\label{secEmigrationRules}
%
We recall that, in our context, the agents are cities and interactions among them consist in exchanges of inhabitants due to emigration patterns. An emigration model can be specified by choosing particular forms for the functions $\psi,\,\psi_\ast$ in~\eqref{eqUrbanInteractionRules}, replacing the indices $i,\,i_\ast$ by $x,\,x_\ast$ consistently with the dense graph limit.

We consider the class of emigration models described by the rules
\begin{equation}
    s'=s\bigl(1-E(x,s,x_\ast,s_\ast)\bigr), \qquad s_\ast'=s_\ast+sE(x,s,x_\ast,s_\ast),
    \label{eqExplicitUrbanInteractionRules}
\end{equation}
partially inspired by~\cite{Gualandi2019SDC}. Here, $E$ is the \textit{emigration rate}.
            
\begin{rmk}
One could also define
$$ s'=s\bigl(1-E(x,s,x_\ast,s_\ast)+\eta\bigr), $$
where $\eta$ is a random variable expressing the natural stochastic fluctuations of the population size in a city due e.g., to birth and death rates. Rules~\eqref{eqExplicitUrbanInteractionRules} just assume that the birth and death rates always compensate.
\end{rmk}
            
We further characterise $E$ by means of the following
\begin{hyp} \label{hypEmigrationRateLimits}
To ensure physically meaningful, i.e. non-negative, post-interaction states $s',\,s_\ast'$, we assume that there exists a constant $\bar E\in (0,\,1)$ such that
$$ 0\leq E(x,s,x_\ast,s_\ast)\leq\bar E, \qquad \forall s,\,s_\ast\in\R_+,\ \forall\,x,\,x_\ast\in [0,\,1]. $$
\end{hyp}

It is not difficult to see that rules~\eqref{eqExplicitUrbanInteractionRules} conserve the total population at each exchange of people between cities:
$$ s'+s_\ast'=s+s_\ast. $$
This is reasonable, since every pattern of emigration from one city corresponds to a pattern of immigration into another, leaving the combined population of the two interacting cities unchanged.

\section{Implementation and simulations of the kinetic model}
\label{secImplementationSimulation}
%
In this section, we consider several possible models for the emigration rate $E$ in~\eqref{eqExplicitUrbanInteractionRules} and we analyse the corresponding predictions of the Boltzmann-type equations on both the exact graph topology, cf.~\eqref{eqUrbanAsymmetricHomogeneousWeakFormBoltzmannType}, and the approximate random graph topology, cf.~\eqref{eqRelaxedGraphBoltzmannType}. Given the lack of well-established physical theories of emigration patterns, our main aim is to identify a minimal set of empirical factors that may influence the distribution of the population across cities, by comparing the model predictions with real-world data.

\subsection{Monte Carlo method}
\label{secMonteCarloMethod}
%
\begin{algorithm}[!t]
    \caption{Nanbu--Babovsky-type Monte Carlo algorithm for~\eqref{eqUrbanInteractionAlgorithm}}
    \label{algMonteCarlo}
    \KwData{%
        \(N\in\mathbb{N}_+\) (number of cities),
        \(\Delta t\leq 1\) (time step),
        \(T>0\) (final time),
        \(P\) (initial total population), %
        \(\mathbf A=(a_{i,j})_{i,j=1}^{N}\) (adjacency matrix)%
    }
    \(\boldsymbol\xi\gets(\xi_1,\xi_2,\dots,\xi_N)\sim\operatorname{Uniform}(\{0,\dots,N\}^N)\)\;
    \(\mathcal{S}^0\gets(s_1^0,s_2^0,\dots,s_N^0)=\frac{P}{N}\mathbf 1\in\mathbb{R}_+^N\)\;\label{algMonteCarloDistribuzioneIniziale}
    \For{\(n=0,1,2,\dots,\left\lfloor T/\Delta t\right\rfloor-1\)}{
        \(p\gets\) random permutation of \(\{1,2,\dots,N\}\)\;
        \For{\(k=1,2,\dots,\left\lfloor N/2\right\rfloor\)}{
            \(i\gets p(k)\); \(i_\ast\gets p(\left\lfloor N/2\right\rfloor+k)\)\;
            \eIf{the exact graph topology is used}{
                sample \(\Theta\sim\operatorname{Bernoulli}(a_{i,i_\ast}\Delta t)\)\;
                \(x_i\gets\frac{1}{N}\sum_{i_\ast=1}^{N}a_{i,i_\ast}\)\;
                \(x_{i_\ast}\gets\frac{1}{N}\sum_{i=1}^{N}a_{i,i_\ast}\)\;
            }{
                sample \(\Theta\sim\operatorname{Bernoulli}\!\left(\frac{\xi_i\xi_{i_\ast}}{d_N}\Delta t\right)\)\;
                \(x_i\gets\frac{\xi_i}{N}\)\;
                \(x_{i_\ast}\gets\frac{\xi_{i_\ast}}{N}\)\;
            }
            \eIf{\(\Theta=1\)}{
                \(E\gets E(x_i,s_i^n,x_{i_\ast},s_{i_\ast}^n)\)\;
                \(s_i^{n+1}\gets s_i^n(1-E)\)\;
                \(s_{i_\ast}^{n+1}\gets s_{i_\ast}^n+s_i^nE\)\;
            }{
                \(s_i^{n+1}\gets s_i^n\)\;
                \(s_{i_\ast}^{n+1}\gets s_{i_\ast}^n\)\;
            }
        }
        \(\mathcal{S}^{n+1}\gets(s_1^{n+1},s_2^{n+1},\dots,s_N^{n+1})\)\;
    }
\end{algorithm}

To solve~\eqref{eqUrbanAsymmetricHomogeneousWeakFormBoltzmannType},~\eqref{eqRelaxedGraphBoltzmannType} numerically, we use a Nanbu--Babovsky-type Monte Carlo scheme~\cite{babovsky1986M2AS,nanbu1980JPSJ}, which consists in simulating directly the binary interactions described by the agent-based stochastic algorithm~\eqref{eqUrbanInteractionAlgorithm} supplemented by the interaction rules~\eqref{eqUrbanInteractionRules},~\eqref{eqExplicitUrbanInteractionRules}. We describe the scheme in detail in Algorithm~\ref{algMonteCarlo}.

It is worth stressing that, in line~\ref{algMonteCarloDistribuzioneIniziale} of Algorithm~\ref{algMonteCarlo}, the array $\mathcal{S}^0$ of the initial city sizes may in principle be sampled from whatever probability distribution. However, due to the lack of information on a typical probability distribution of city size, we decided to sample $\mathcal{S}^0$ from a uniform probability distribution with respect to a given initial total population size $P$.
    
\subsection{Numerical results}
\label{secNumericalResults}
%
All real-world data used to inform the numerical simulations, as well as those used for comparison with the numerical results, were obtained from the Italian National Institute of Statistics (ISTAT). The data refer to $1991$ and to the Italian region of Sardinia, a choice primarily inspired by~\cite{DeMontis2007}. In more detail,~\cite{ISTATPendolarismo1991} provided the data to build the graph and its adjacency matrix, while~\cite{ISTATCensimento1991} provided the actual population distributions to which the numerical results were then compared.

\begin{table}[!t]
\centering
\caption{Data and parameters of Sardinia in $1991$~\cite{ISTATCensimento1991,ISTATPendolarismo1991}}
\begin{tabular}{cccccccc}
    \toprule
    $N$ & $P$ & $\beta$ & $\zeta$ & $\mu_1$ & $\sigma_1$ & $\mu_2$ & $\sigma_2$ \\
    \midrule
    $375$ & $1.648\cdot 10^6$ & $1.273$ & $0.929$ & $3.227$ & $0.408$ & $4.099$ & $0.556$ \\
    \bottomrule
    \end{tabular}
\label{tabSardiniaDataParameters}
\end{table}

The main data ($N$, $P$) of the Italian region of Sardinia are listed in Table~\ref{tabSardiniaDataParameters} along with the fitted parameters of the Pareto ($\beta$, cf.~\eqref{eqComplementaryCumulativePareto}) and bilognormal distributions ($\zeta$, $\mu_1$, $\sigma_1$, $\mu_2$, $\sigma_2$, cf.~\eqref{eqLognormalDistribution},~\eqref{eqBimodalLognormalDistribution}).
                        
Concerning the plots, their complex definitions are not necessary to understand the results shown hereafter. Hence, they are omitted here and deferred to Appendix~\ref{apxPlotDefinitions} for the interested reader. Nonetheless, three features common to all numerical tests have to be kept in mind:
\begin{enumerate}[label=(\roman*)]
\item due to the stochastic nature of Algorithm~\ref{algMonteCarlo} and to the relative smallness of the number $N$ of cities considered, cf. Table~\ref{tabSardiniaDataParameters}, typically in each test $R=100$ simulations have been performed and subsequently averaged;
\item the time step is kept fixed to $\Delta t=10^{-2}$ while the number $N_t$ of time steps is variable depending on the desired final time $T$;
\item for the sake of brevity, only the best fitting configuration empirically found is analysed in each test.
\end{enumerate}

In the following, we consider various possible choices for the emigration rate $E$ to specify the interaction model~\eqref{eqExplicitUrbanInteractionRules}, and discuss the resulting outcomes.

\subsubsection{Size rule}
\label{secSER}
%
A first possibility consists in choosing
\begin{equation}
    E(s,s_\ast):=\lambda\frac{\left(s_\ast/s\right)^\alpha}{1+\left(s_\ast/s\right)^\alpha},
    \label{eqSizeEmigrationRule}
\end{equation}
independent of the latent variables $x,\,x_\ast$, where $\lambda\in (0,\,1)$, consistently with Assumption~\ref{hypEmigrationRateLimits}, and $\alpha\in\R_+$. This is a Hill function of order $\alpha$, partially inspired by~\cite{Gualandi2019SDC}, which is also largely applied in biomedical contexts.

The rationale behind~\eqref{eqSizeEmigrationRule} is that the larger the relative population size of the destination city compared to that of the origin city the larger the emigration rate.

\newsavebox{\stabularbox}
\begin{table}[!t]
\caption{Parameters of all the studied reference configurations}
\label{tabParmetersReferenceConfigurations}
%
\newcommand\subtabular[4]{
\sbox{\stabularbox}{\begin{tabular}{@{} #1 @{}}#2\end{tabular}}%
                    \begin{subtable}[b]{\wd\stabularbox}
                        \centering
                        \begin{tabular}{@{} #1 @{}}#2\end{tabular}
                        \caption{#3}
                        \label{#4}
                    \end{subtable}
}
\newcommand\tableindent{\hskip 0pt plus .4fill}
\tableindent
\subtabular{lcccc}{
\toprule
& $\lambda$ & $\alpha$ & $N_t$ & $R$ \\
\midrule
Stable & $0.1$ & $0.5$ & $10^7$ & $100$ \\
Unstable & $0.1$ & $1.5$ & $10^6$ & $10$ \\
\bottomrule
}
{Figures~\ref{figSERStableReferenceConfiguration},~\ref{figSERUnstableReferenceConfiguration}}{tabSERParameters}
\hfill
\subtabular{ccc}{
\toprule
$\lambda$ & $\alpha$ & $N_t$ \\
\midrule
$0.15$ & $0.63$ & $3\cdot 10^6$ \\
\bottomrule
}
{Figure~\ref{figSDERReferenceConfiguration}}{tabSDERParameters}
\hfill
\subtabular{ccc}{
\toprule
$\lambda$ & $\alpha$ & $N_t$ \\
\midrule
$0.18$ & $0.44$ & $3\cdot 10^6$ \\
\bottomrule
}
{Figure~\ref{figSSERReferenceConfiguration}}{tabSSERParameters}
\tableindent
\mbox{}
\end{table}
    
\tikzexternalizedfigure[lualatex]{p}[-.5em]
  {../../Figure/.fnl/Articolo/Capitolo3/20RETConfigurazioneRiferimento/Stabile/Infopagina.tex}
  {Study of the \textit{stable} reference configuration of~\eqref{eqSizeEmigrationRule} with parameters from Table~\ref{tabSERParameters}}
  {.6\linewidth}{figSERStableReferenceConfiguration}

\tikzexternalizedfigure[lualatex]{p}[-1em]
  {../../Figure/.fnl/Articolo/Capitolo3/20RETConfigurazioneRiferimento/Instabile/Infopagina.tex}
  {Study of the \textit{unstable} reference configuration of~\eqref{eqSizeEmigrationRule} with parameters from Table~\ref{tabSERParameters}}
  {.62\linewidth}{figSERUnstableReferenceConfiguration}

We studied two reference configurations described by the parameters in Table~\ref{tabSERParameters}. The corresponding results are illustrated in Figures~\ref{figSERStableReferenceConfiguration},~\ref{figSERUnstableReferenceConfiguration}. The latter describe two trends: a stable one, which is not distant from the initial uniform distribution, and an unstable one, in which, on average, the medium-size cities depopulate.

We stress that the latter conclusion holds on average, since in single realisations it is much more likely that the larger cities depopulate, whereas over multiple realisations at least one with a large size manifests, thereby balancing the average. This is also the reason why only $10$ runs were considered instead of the usual $100$ in Table~\ref{tabSERParameters}, given that the average tends to mask the depopulation, which is nonetheless inevitable albeit slower. In addition to Figures~\ref{figSERUnstableReferenceConfiguration}j,l, the histograms in Figures~\ref{figSERUnstableReferenceConfiguration}a,c also confirm this tendency. Specifically, their tails elongate towards negative values, corresponding to the exponents of the cities that depopulate.
    
More generally, we have verified, albeit only numerically, that for every $\lambda>0$ there exists $\bar\alpha\in\R_+$ such that $\alpha<\bar\alpha$ results in a stable configuration while $\alpha\geq\bar\alpha$ yields an unstable one.

\subsubsection{Size-degree rule}
\label{secSDER}
%
To improve~\eqref{eqSizeEmigrationRule}, we introduce the relative degree ratio:
\begin{equation}
    E(x,s,x_\ast,s_*):=\lambda\frac{\left[\left(s_\ast/s\right)\cdot \left(x_\ast/x\right)\right]^\alpha}{1+\left[\left(s_\ast/s\right)\cdot \left(x_\ast/x\right)\right]^\alpha}.
    \label{eqSizeDegreeEmigrationRule}
\end{equation}
The rationale is that the more populous and \textit{better connected} the destination city $(x_\ast,\,s_\ast)$ is, the more attractive it is relative to the origin city $(x,\,s)$.

\tikzexternalizedfigure[lualatex]{p}
  {../../Figure/.fnl/Articolo/Capitolo3/20RETGConfigurazioneRiferimento/Infopagina.tex}
  {Study of the reference configuration of~\eqref{eqSizeDegreeEmigrationRule} with parameters from Table~\ref{tabSDERParameters}}
  {.58\linewidth}{figSDERReferenceConfiguration}

The studied reference configuration is described by the parameters in Table~\ref{tabSDERParameters}, while the results are illustrated in Figure~\ref{figSDERReferenceConfiguration} (cf. also related figures in the original source for comparison where appropriate). Overall, these results improve those obtained using~\eqref{eqSizeEmigrationRule} in several aspects:
\begin{itemize}
\item Figures~\ref{figSDERReferenceConfiguration}a, d, j show very good agreement between the results obtained from the exact graph topology and its random-graph-approximation;
\item in Figures~\ref{figSDERReferenceConfiguration}b, c, the simulated distributions are fairly similar to the real ones; in Figures~\ref{figSDERReferenceConfiguration}e, f the Pareto indices are reasonably close;
\item in Figures~\ref{figSDERReferenceConfiguration}g--i, the bilognormal profile follows the terminal part of the tail better than the Pareto profile.
\end{itemize}

On the other hand, in Figures~\ref{figSDERReferenceConfiguration}k, l we notice that the scatter plots present a similar positive correlation, albeit much tighter than the real one. It is nearly a one-to-one correlation, a clearly quite unrealistic outcome.

\subsubsection{Size-strength rule}
\label{secSSER}
%
In~\eqref{eqSizeDegreeEmigrationRule}, we may interpret the latent variable $x$ as the \textit{strength} of a city, i.e. the total number of commuters associated with that city, instead of the number of connections. This way, the same emigration rule~\eqref{eqSizeDegreeEmigrationRule} can be used to model the fact that the more populous and \textit{heavily trafficked} the destination city is, the more attractive it is relative to the origin city.

\tikzexternalizedfigure[lualatex]{p}
  {../../Figure/.fnl/Articolo/Capitolo3/20RETFConfigurazioneRiferimento/Infopagina.tex}
  {Study of the reference configuration of~\ref{eqSizeDegreeEmigrationRule}, in which the latent variable $x$ is reinterpreted as the strength of a city, with parameters from Table~\ref{tabSSERParameters}}
  {.577\linewidth}{figSSERReferenceConfiguration}

The reference configuration analysed here is now described by the parameters listed in Table~\ref{tabSSERParameters}, while the new results are shown in Figure~\ref{figSSERReferenceConfiguration}. Note that, in this case, we are not changing the mathematical model itself, but rather the type of real-world data used to validate its results. Indeed, these results are the best according to several criteria:
\begin{itemize}
\item Figures~\ref{figSSERReferenceConfiguration}a, d, j confirm that the simulation outcomes obtained using the random-graph approximation of the graph topology correspond almost exactly to those obtained using the exact graph topology;
\item in Figures~\ref{figSSERReferenceConfiguration}b, c, the simulated distributions are almost identical to the real one, and the relative Pareto indices in Figures~\ref{figSSERReferenceConfiguration}e, f also correspond;
\item Figures~\ref{figSSERReferenceConfiguration}h, i do not deviate significantly from the previous configuration;
\item the scatter plots in Figures~\ref{figSSERReferenceConfiguration}k, l show a larger spread than in the previous case while maintaining the same positive correlation with the real-world data. This captures correctly the size differentiation for a given degree.
\end{itemize}
    
The only drawback is visible in Figures~\ref{figSSERReferenceConfiguration}k, l: the emigration rule fails to correctly predict, on average, the city sizes in the lower portion of the real-world scatter plot. Specifically, the model prediction overestimates the size of the smaller cities and slightly underestimates that of the larger cities. Yet, on second thought, this is an issue already present in the previous case, cf. the slope of the line described by the scatter plot in Figures~\ref{figSDERReferenceConfiguration}k, l. We believe that it is possibly caused by the absence of intra-urban commuters in the reinterpretation of $x$ as the strength of a city, which considers only inter-urban commuters. Including also this aspect may significantly increase the strength of the medium-sized and larger cities and only marginally that of the smaller ones, thereby presumably correcting the slope.

\subsection{Interpretations of the results}
\label{secInterpretations}
%
In this section, we discuss the previous results in more detail, with the aim of identifying the microscopic, i.e. city-level, factors that are most likely to influence migration.
        
\subsubsection{Size rule~\texorpdfstring{\eqref{eqSizeEmigrationRule}}{}}
From the results of~\eqref{eqSizeEmigrationRule}, it is clear that the relative population size alone is insufficient to accurately describe the city-size distribution. Indeed, if $s\leq s_\ast$, the destination city will invariably attract people from the origin city (and vice versa). Consequently, as soon as a city loses even a small portion of its population relative to all its neighbouring cities, it is bound to continue depopulating. This mechanism can explain the depopulation, \textit{on average}, of medium-sized cities: large cities draw population away from medium-sized cities, making the latter relatively unattractive vertices from the outset. In the short term, medium-sized cities depopulate to such an extent that they become relatively unattractive even compared to small cities. The latter do not interact directly with large cities and therefore do not initially experience the same depopulation as medium-sized cities. The outcome of these dynamics is shown in Figures~\ref{figSERUnstableReferenceConfiguration}k, l, at least in the unstable case.

\tikzexternalizedfigure{t}
  {../../Figure/.fnl/Elaborato/Capitolo4/FunzioneHillGradoa.tex}
  {Hill function for different values of the degree $\alpha$.}
  {.775\linewidth}{figDegreeAHillFunction}

Conversely, in the stable case, the interpretation is less straightforward. Certainly, $\alpha$ is not sufficiently small for the Hill function in~\eqref{eqSizeEmigrationRule} to become approximately constant (cf. Figure~\ref{figDegreeAHillFunction}) and equal to $\lambda/2$, particularly given that~\eqref{eqSizeDegreeEmigrationRule}, for the reference configuration in Table~\ref{tabSSERParameters}, yields an even lower value of $\alpha$. We can therefore only conclude that the issue lies in the emigration rule itself.

Overall, these results suggest that migration most likely does not depend \textit{solely} on the relative population sizes of the interacting cities.
            
\subsubsection{Size-degree rule~\texorpdfstring{\eqref{eqSizeDegreeEmigrationRule}}{}}
The results obtained with the emigration rule~\eqref{eqSizeDegreeEmigrationRule} are considerably more realistic. Nevertheless, unlike relative population size, the introduction of relative degree through the ratio $x_\ast/x$ is not immediately intuitive. In other words, one might wonder why it is reasonable to assume that a city with a low degree, i.e. one that is less connected, should also have a small population. There is a risk that the model is being forced towards results that, although consistent with the real-world data, are not realistic in terms of the underlying mechanisms. This modelling choice can, however, be motivated by two naturally observable constraints.
\begin{itemize}
\item The first is a physical constraint arising from the spatial nature of the graph. By their nature, edges occupy the spaces between cities and connect to them along portions of their boundary. Therefore, a city can increase its degree only by expanding so as to have a larger boundary, which in turn requires a larger population to sustain it.
\item The second is a social constraint arising from commuter flows: the greater the connectivity, the more complex the transport network, and such complexity can be justified only if there are sufficiently many commuters using it. Such commuter flows, in turn, are more likely to be attracted to a city with a sufficiently large population to sustain them.
\end{itemize}

The social constraint is indeed supported by Figure~\ref{figStrengthVsDegree20}, which shows a positive correlation between degree and strength.
                
The physical constraint justifies this type of interaction rule because the graph is inherently spatial. Similarly, the social constraint supports this interaction rule provided that degree and strength are positively correlated. Together, these considerations suggest the relationship $s\propto x$, which is also supported by the empirical scatter shown in Figures~\ref{figSDERReferenceConfiguration}k, l. This relationship makes the introduction of the ratio $x_\ast/x$ plausible from a modelling perspective.

\tikzexternalizedfigure{t}
  {../../Figure/.fnl/Articolo/Capitolo3/SVsD.tex}
  {Strength vs degree of the Sardinian network}
  {.75\linewidth}{figStrengthVsDegree20}
            
It is also worth considering why the emigration rule~\eqref{eqSizeDegreeEmigrationRule} leads the system to a non-trivial equilibrium configuration at long times. In particular, one may wonder why the population flow ``small cities $\to$ medium-sized cities $\to$ large cities'' does not occur, as this would ultimately lead to the depopulation of small cities. This can only be explained by the presence of a backflow naturally induced by~\eqref{eqSizeDegreeEmigrationRule}, which balances the flow towards relatively more populous and better-connected cities. In other words, introducing degree in~\eqref{eqSizeEmigrationRule} not only makes medium-sized cities more attractive than small cities in the short term, but also stabilises the long-term, i.e. stationary, interactions. However, the precise mechanism producing this stabilising reverse flow remains unclear, since the ratio $r$ in the Hill function (cf. Figure~\ref{figDegreeAHillFunction}) should be smaller in~\eqref{eqSizeDegreeEmigrationRule} than in~\eqref{eqSizeEmigrationRule}.
                
In conclusion, these results suggest that, according to~\eqref{eqSizeDegreeEmigrationRule}, migration is likely to be governed by \textit{at least} two factors: the relative population sizes and the connectivity of the interacting cities.
            
\subsubsection{Size-strength rule~\texorpdfstring{\eqref{eqSizeDegreeEmigrationRule}}{}}
Reinterpreting the latent variables $x$ and $x_\ast$ in~\eqref{eqSizeDegreeEmigrationRule} as the strengths of the origin and destination cities makes it possible to distinguish between cities with the same degree but not necessarily the same commuter traffic. This natural distinction allows us to relax the almost one-to-one relationship between degree and population size observed in Figures~\ref{figSDERReferenceConfiguration}k, l, replacing it with the more realistic positive correlation shown in Figures~\ref{figSSERReferenceConfiguration}k, l.
            
Since these results provide the best agreement both qualitatively and quantitatively, we conclude that the migration is \textit{most likely} governed by the relative population sizes and strengths of the interacting cities.
            
\subsubsection{The roles of~\texorpdfstring{$\boldsymbol{\lambda}$}{} and~\texorpdfstring{$\boldsymbol{\alpha}$}{}}
%
\begin{table}[!t]
\centering
\caption{Parametric studies in the Supplementary Material}
\label{tabRETFParametricStudies}
\begin{tabular}{lcccc}
\toprule
& $N_p$ & $\lambda$ & $\alpha$ & $N_t$ \\
\midrule
Parametric study on $\lambda$: & $5$ & $[0.01,\,0.3]$ & $0.44$ & $10^7$ \\
Parametric study on $\alpha$: & $5$ & $0.18$ & $[0.3,\,0.6]$ & $3\cdot 10^6$ \\
\bottomrule
\end{tabular}
\end{table}
                
All the proposed emigration rules involve the parameters $\lambda$ and $\alpha$. To investigate their roles, we performed two parametric studies, which are presented in the Supplementary Material for conciseness. Both studies are based on the reference configuration for~\eqref{eqSizeDegreeEmigrationRule} reported in Table~\ref{tabSSERParameters}, with only one of the two parameters varied at a time. Table~\ref{tabRETFParametricStudies} summarises the main settings adopted for these parametric studies, where $N_p=5$ denotes the number of values uniformly sampled from the specified intervals for each parameter.

\tikzexternalizedfigure[lualatex]{p}
  {../../Figure/.fnl/Elaborato/Capitolo1/EsempioFunzioneLognormale.tex}
  {Probability density function of a random variable $X\sim\operatorname{Lognormal}(\mu,\,\sigma)$ for different parameters values, with the expected value held constant at $\mathbb{E}[X]=10^{\mu+\sigma^2/2}=10$. The distributions are shown on (a) linear, (b) logarithmic, and (c) log--log scales. Dashed lines indicate the expected value of each distribution}
  {.65\linewidth}{figLognormalFunctionExample}

One effect of $\lambda$ and $\alpha$ is that they influence the parameters $\mu$ and $\sigma^2$ of the lognormal distribution~\eqref{eqLognormalDistribution}. This is because, owing to the conservation of the total population in each binary exchange between cities, the sample mean, fixed initially as $P/N$ in Algorithm~\ref{algMonteCarlo}, cf. line~\ref{algMonteCarloDistribuzioneIniziale}, remains constant over time. Consequently, the mean of the lognormal distribution must also remain constant, i.e. $10^{\mu+\sigma^2/2}=P/N$, independently of the values of $\lambda$ and $\alpha$. This, in turn, implies that $\mu$ and $\sigma$ vary with $\lambda$ and $\alpha$. This behaviour is difficult to discern on either a linear or a logarithmic scale, but becomes more apparent on a log--log scale, since $\mu$ and $\sigma$ are precisely the mean and standard deviation of the corresponding Gaussian distribution after logarithmic transformation. Similar reasoning applies to the bilognormal fit~\eqref{eqBimodalLognormalDistribution}, since it is a convex combination of two lognormal distributions.

In more detail, the parametric study in which $\lambda$ is varied while $\alpha$ is held fixed leads to the following conclusions:
\begin{itemize}
\item as $\lambda$ increases, the variance increases, thereby broadening the distribution, albeit only slightly;
\item the Pareto index remains \textit{de facto} almost constant, with an average variation of only approximately $0.05$;
\item the relationship between the Pareto and bilognormal distributions remains unchanged: the latter provides a better approximation of the far end of the tail, whereas the two distributions coincide near its beginning;
\item the degree of scatter remains essentially unchanged as $\lambda$ varies, and the average behaviour is the same;
\item all considered cases converge to a steady state, with the convergence time decreasing as $\lambda$ increases.
\end{itemize}
We further observe that all the variations described above are small, despite $\lambda$ increasing by as much as an order of magnitude. We therefore conclude that $\lambda$ has little influence on the results, except for the convergence time. Nonetheless, for $\lambda>0.3$, the system reaches an unstable configuration.

Conversely, the parametric study in which $\alpha$ is varied while $\lambda$ is held fixed leads to the following conclusions:
\begin{itemize}
\item as $\alpha$ increases, the variance increases significantly, thereby reducing the mean and broadening the distribution;
\item the Pareto index varies significantly, decreasing by nearly an order of magnitude;
\item the relationship between the Pareto and bilognormal distributions in terms of their approximation of the tail remains unchanged;
\item for small values of $\alpha$, the degree of scatter leads to an overestimation of the size of the minor and medium-sized cities and an underestimation of those of large cities; as $\alpha$ increases, this tendency is reversed;
\item all considered cases converge to a steady state, with the convergence time increasing as $\alpha$ increases.
\end{itemize}
In this case, all the variations described above are significant. In particular, the Pareto index appears to be strongly affected by the value of $\alpha$. Similarly, for $\alpha>0.6$, the system reaches an unstable configuration.

Overall, in light of these parametric studies, we propose the following interpretations of the roles of $\lambda$ and $\alpha$:
\begin{itemize}
\item $\lambda$ is an \textit{attraction} parameter: the greater its value, the larger the fraction of the population that large cities can attract from smaller cities in a single interaction;
\item $\alpha$ is a \textit{repulsion} parameter: the greater its value, the less likely people from large cities are to move towards small cities. This effect is clearly inversely related to the magnitude of the stabilising backflow.
\end{itemize}
Alternatively, these roles can be interpreted in terms of the Hill function (cf. Figure~\ref{figDegreeAHillFunction}): $\lambda$ mainly influences the part of the function corresponding to $r>1$, and hence the attraction, whereas $\alpha$ mainly influences the part corresponding to $r\leq 1$, and hence the repulsion.

\section{Conclusions}
\label{secConclusions}
%
The results presented in this paper support the vertex-agent perspective as a modelling framework in an urban context, both for the city-size distribution and for its Pareto index. The random-graph approximation introduced in Definition~\ref{defDegreeMatrix} produces results that closely match those obtained using the exact graph topology. The majority of the proposed interaction rules faithfully reproduce both the empirical distribution and the Pareto index of its tail. However, the bimodal lognormal distribution provides a better fit to the tail better than the Pareto distribution, confirming the findings of~\cite{Gualandi2019SDC,Gualandi2019HBLD}. Both fits have advantages and disadvantages:
\begin{itemize}
    \item From an applied perspective, the Pareto distribution is the easier one to fit: it depends on a single parameter, which the data determine in a nearly univocal way, so that the same fit is always recovered, which makes it the preferable choice from an engineering standpoint;
    \item From a mathematical perspective, the bimodal lognormal distribution is instead the preferable one: it depends on five parameters, so that fitting it does not always lead to the same result, and two of them in particular are more difficult to compare.
\end{itemize}

The work presented here is by no means exhaustive and opens up several avenues for future research. In particular, possible developments include:
\begin{itemize}
\item incorporating intra-urban commuting into~\eqref{eqSizeDegreeEmigrationRule}, or considering interaction rules involving other graph parameters, such as betweenness centrality or the clustering coefficient;
\item deriving the Fokker--Planck approximation of~\eqref{eqUrbanAsymmetricHomogeneousWeakFormBoltzmannType} as a basis for the analytical study of the stationary distribution;
\item relaxing the simplifying Assumption~\ref{hypStaticGraph} of a non-coevolving network in order to model dynamic connections that evolve alongside population migration, with a possible comparison with empirical data, where available;
\item investigating whether the results can be reproduced for European or international networks, thereby assessing the validity of the present theory at larger scales.
\end{itemize}

\appendix
    
\section{Results reproducibility}

\subsection{Technical specifications}
The code used to perform the numerical simulations is available in the repository~\cite{Taralli2026}. It was written entirely in Python 3.13.7, and the simulations were run on a Windows-based machine equipped with an Intel Core i5-7500 processor (quad-core, 3.40 GHz).

\subsection{Confidence interval}
Let $R\in\N_+$ denote the number of simulations performed, and let $\mathbf S\in\R_+^R$ be the random vector whose components represent the population size of a city. By construction of Algorithm~\ref{algMonteCarlo}, the components $S_1,\,S_2,\,\dots,\,S_R$ of $\mathbf S$ are independent and identically distributed, which allows us to define
$$ \bar S:=\frac{1}{R}\sum_{r=1}^{R}S_r, \qquad V:=\frac{1}{R-1}\sum_{r=1}^R{(S_r-\bar S)}^2 $$
as the sample mean and variance, respectively, of the $S_r$'s, $r=1,\,\dots,\,R$. If $\mu\in\R$ is the true mean of the distribution of the $S_r$'s, Student's \(T\)-distribution with parameter $R-1$ can be written as
$$ T:=\frac{\bar S-\mu}{\sqrt{\frac{V}{R}}}. $$
From this, we can define the (symmetric) confidence interval with significance level $\delta\in (0,\,1)$ by setting
$$ \operatorname{Prob}\!\left(-t_{R-1}^{\delta/2}\leq T\leq t_{R-1}^{\delta/2}\right)=1-\delta, $$
where $t_{R-1}^{\delta/2}\in\R$ is such that
$$ \operatorname{Prob}\!\left(T<t_{R-1}^{\delta/2}\right)=1-\frac{\delta}{2}. $$

Expanding $T$ and isolating the true mean $\mu$, we obtain
$$ \operatorname{Prob}\!\left(\bar{S}-t_{R-1}^{\delta/2}\sqrt{\frac{V}{R}}\leq\mu\leq\bar{S}+t_{R-1}^{\delta/2}\sqrt{\frac{V}{R}}\right)=1-\delta, $$
whence
\begin{equation}
    \IC^\delta_R(\mathbf S):=\left[\bar{S}-t_{R-1}^{\delta/2}\sqrt{\frac{V}{R}},\,\bar{S}+t_{R-1}^{\delta/2}\sqrt{\frac{V}{R}}\right]
    \label{eqConfidenceInterval}
\end{equation}
is the interval estimator defining the desired confidence interval.

The confidence interval essentially represents the statistical error one is prepared to accept: it assesses how likely the interval estimator~\eqref{eqConfidenceInterval} is to contain the parameter $\mu$. In other words, $\IC^\delta_R$ quantifies the uncertainty associated with estimating the mean: for a given sample $\mathbf s$ of the random vector $\mathbf S$, the wider the confidence interval, the greater the uncertainty with which the sample mean $\bar s$ estimates the true mean $\mu$. Conversely, the narrower the confidence interval, the more precise the estimate $\bar s$, in the sense that $\mu$ lies within a small neighbourhood of the sample mean. From this perspective, the confidence interval is conceptually analogous to the precision of a measuring instrument.

In all the tests presented in Section~\ref{secImplementationSimulation}, we set $\delta=0.05$, corresponding to a confidence level of $95\%$.

The interval estimate can then be obtained directly from the interval estimator~\eqref{eqConfidenceInterval} by considering the realisations $\mathbf s$ of the random vector $\mathbf S$, which are precisely the outcomes of the $R$ numerical simulations.
        
\subsection{Data structure}
The data are organised into a fourth-order tensor $\mathbf s=(s^{n,a}_{i,r})\in\R_+^{N_f\times 2\times N\times R}$, whose indices have the following meanings:
\begin{itemize}
\item $n$ denotes the time index;
\item $a$ denotes the type of simulation ($a=1$ for the exact graph topology; $a=2$ for the approximated graph topology);
\item $i$ denotes the city index;
\item $r$ denotes the simulation number.
\end{itemize}

The time index is further characterised by three parameters in $\N_+$:
\begin{itemize}
\item $N_t:=\left\lfloor\frac{T}{\Delta t}\right\rfloor$ is the total number of simulated time steps;
\item $N_s\ll N_t$ is the number of snapshots retained from the simulated time steps;
\item $N_f<N_s$ is the number of reduced time points obtained from the snapshots.
\end{itemize}

Both the snapshots and the reduced time points are sampled at evenly spaced intervals, with steps $\Delta s:=N_t/N_s$ and $\Delta f:=N_s/N_f$, respectively. For simplicity, we assume that $N_t$ is a multiple of $N_s$ and that $N_s$ is a multiple of $N_f$. We make an analogous assumption also for $T$ and $\Delta t$.

Using Algorithm~\ref{algMonteCarlo}, a total of $N_t$ times with time steps of size $\Delta t$ are simulated, of which $N_s$ are retained and stored in the snapshot tensor $\undertilde{\mathbf s}=(\undertilde s^{n,a}_{i,r})\in\R_+^{N_s\times 2\times N\times R}$. The tensor $\undertilde{\mathbf s}$ is then averaged over the time index according to
$$ s^{n,a}_{i,r}:=\frac{1}{\Delta f}\sum_{j=1}^{\Delta f}\undertilde s^{(n-1)\Delta f+j,a}_{i,r}, $$
so that each consecutive block of $\Delta f$ snapshots is replaced by its average. This smoothing reduces the noise in $\mathbf s$ relative to $\undertilde{\mathbf s}$, thereby improving the readability of the resulting plots.

Three notions of physical time can be defined, depending on the interpretation of the index $n$:
\begin{itemize}
\item $t_n:=n\Delta{t}$, $n=1,\,\dots,\,N_t$, is the discretised time;
\item $t^s_n:=t_n\Delta{s}$, $n=1,\,\dots,\,N_s$, is the snapshot time;
\item $t^f_n:=t^s_n\Delta{f}$, $n=1,\,\dots,\,N_f$, is the reduced time.
\end{itemize}
In each case, $t_{N_t}=t^s_{N_s}=t^f_{N_f}=T$.

In all the tests presented in Section~\ref{secImplementationSimulation}, we set $\Delta t=10^{-2}$, $N_s=1000$, and $N_f=50$. Moreover, we choose $N_t$ such that the final simulated time $T$ is slightly greater than the convergence time.

In most plots, only the distribution at the final time $T$ is considered, \textit{without} time averaging. For ease of notation, the resulting third-order tensor $\mathbf s^T:=\undertilde{\mathbf s}^{N_s}$ is therefore denoted, with a slight abuse of notation, by $\mathbf s^T\in\R_+^{2\times N\times R}$.
        
\subsection{Plot definitions}
\label{apxPlotDefinitions}
%
\subsubsection{Histograms}
The first rows of Figures~\ref{figSERStableReferenceConfiguration}--\ref{figSSERReferenceConfiguration} show the normalised histograms of the city-size distribution at the final time $T$.

We first consider the minimum and maximum entries of the tensor $\mathbf s^T$:
$$ s_{\min}:=\min_{a,\,i,\,r}s^{T,a}_{i,r}, \qquad s_{\max}:=\max_{a,\,i,\,r}s^{T,a}_{i,r}, $$
and use them to construct an evenly spaced grid defining the histogram bins. Let $N_c$ be the number of such bins. We then introduce the function $\mathcal{B}:\R^N\to\R^{N_c}$, which returns the normalised bin values of the histogram, and define the tensor $\mathbf b\in\R_+^{2\times N_c\times R}$ by $\mathbf b^a_{\cdot,r}:=\mathcal{B}(\mathbf s^{T,a}_{\cdot,r})$ for every $a$ and $r$. Finally, for every $a$ and $c$, we evaluate the confidence intervals
$$ \IC_R^\delta(\mathbf b^a_c)=\left[\bar{b}^a_c-t^{\delta/2}_{R-1}\sqrt{\frac{v^a_c}{R}},\,\bar{b}^a_c+t^{\delta/2}_{R-1}\sqrt{\frac{v^a_c}{R}}\right]. $$
We then plot the mean histograms for both the ``exact'' ($a=1$) and ``approximated'' ($a=2$) simulations using the mean vectors $\bar{\mathbf b}^a$, together with the associated statistical uncertainty in each bin, represented by adding and subtracting the vector $t^{\delta/2}_{R-1}\sqrt{\frac{\mathbf v^a}{R}}$.

\subsubsection{Bimodal lognormal distribution}
The first rows of Figures~\ref{figSERStableReferenceConfiguration}--\ref{figSSERReferenceConfiguration} also show the bimodal lognormal distributions fitted to the data represented by the histograms at the final time $T$.

To represent these distributions, we follow a procedure similar to that used for the histograms. Let $\mathcal{L}(\cdot;\,\mathbf s^{T,a}_{\cdot,r}):\R_+\to\R_+$ denote the probability density of a bilognormal random variable, cf.~\eqref{eqBimodalLognormalDistribution}, fitted to the vector $\mathbf s^{T,a}_{\cdot,r}$. Moreover, let $\boldsymbol{\mathcal{L}}(\cdot;\,\mathbf s^{T,a}):\R_+\to\R_+^R$ be the vector-valued function defined by
$$ \boldsymbol{\mathcal{L}}(y;\,\mathbf s^{T,a}):=
    {\left[\mathcal{L}(y;\mathbf s^{T,a}_{\cdot,1}),\,\mathcal{L}(y;\,\mathbf s^{T,a}_{\cdot,2}),\,\dots,\,\mathcal{L}(y;\,\mathbf s^{T,a}_{\cdot,R})\right]}^\top, $$
which collects the fits obtained from all $R$ simulations into a single vector. The corresponding confidence intervals are then given by
$$ \IC_R^{\delta}(\boldsymbol{\mathcal{L}}(y;\,\mathbf s^{T,a}))=
    \left[\bar{\mathcal L}^a(y)-t^{\delta/2}_{R-1}\sqrt{\frac{v^a(y)}{R}},\,\bar{\mathcal L}^a(y)+t^{\delta/2}_{R-1}\sqrt{\frac{v^a(y)}{R}}\right],
        \qquad y\in [s_{\min},\,s_{\max}], $$
where, as before, $\bar{\mathcal L}^a(y)$ denotes the mean function, while $\pm t^{\delta/2}_{R-1}\sqrt{\frac{v^a(y)}{R}}$ represents its statistical uncertainty.

Since $y$ varies continuously over the interval $[s_{\min},\,s_{\max}]$, the collection of pointwise confidence intervals defines a pointwise confidence envelope for the fitted functions.

\begin{rmk}
We compute confidence intervals both for the fitted bilognormal density as a whole and for each of its five parameters $\zeta,\,\mu_1,\,\sigma_1,\,\mu_2,\,\sigma_2$. For illustrative purposes, we report only the confidence interval for $\zeta$. Let $\boldsymbol{\zeta}\in\R^{2\times R}$ be the matrix collecting the values of $\zeta$ obtained from all simulations. Then, for each $a$,
$$ \IC_R^\delta(\boldsymbol{\zeta})=\left[\bar\zeta^a-t^{\delta/2}_{R-1}\sqrt{\frac{v^a}{R}},\,\bar\zeta^a+t^{\delta/2}_{R-1}\sqrt{\frac{v^a}{R}}\right]. $$
\label{ossBiLognormalParameterFitting}
\end{rmk}

\subsubsection{Empirical complementary cumulative distribution function}
\label{secEmpiricalCCDF}
%
The second and third rows of Figures~\ref{figSERStableReferenceConfiguration}--\ref{figSSERReferenceConfiguration} show the empirical Complementary Cumulative Distribution Function (CCDF) over the tail of the city-size distribution.

To characterise the tail of the distribution quantitatively, we take the upper quartile as threshold, i.e., that value $s_{3/4}$ of the city size $s$ such that the probability mass for $s\leq s_{3/4}$ is exactly $\frac{3}{4}$. More precisely, $\operatorname{Prob}(S\leq s_{3/4})=\frac{3}{4}$. In practice, this amounts to identifying the element $s_{3/4}$ of the tensor $\mathbf s^{T,a}_{\cdot,r}$ such that
$$ \left\vert\left\{s^{T,a}_{i,r}\geq s_{3/4}\,:\, i=1,\,\dots,\,N\right\}\right\vert=\left\lceil\frac{N}{4}\right\rceil=:N_{1/4},
    \qquad \forall\,a,\,r, $$
so that at least one quarter of the entries of $\mathbf s^{T,a}_{\cdot,r}$ are greater than or equal to $s_{3/4}$. We then regard the values $s\in [s_{3/4},\,s_{\max}]$ as constituting the tail of the distribution.

The empirical CCDF, denoted by $\CCDF_{\mathcal E}(\cdot;\,\mathbf s^{T,a}_{\cdot,r}):\R_+\to\R_+$, provide an empirical estimate of the probability $\operatorname{Prob}(S>y)$ and is defined by
\begin{equation}
    \CCDF_{\mathcal E}(y;\,\mathbf s^{T,a}_{\cdot,r}):=1-\frac{1}{N}\sum_{i=1}^N\chi_{[0,\,y]}(s^{T,a}_{i,r}),
    \label{eqComplementaryCumulativeDistributionFunction}
\end{equation}
where $\chi_{[0,\,y]}$ denotes the characteristic function of the interval $[0,\,y]$. Since~\eqref{eqComplementaryCumulativeDistributionFunction} is evaluated only for $y\in [s_{3/4},\,s_{\max}]$, we introduce the tensor $\mathbf p\in\mathbb R^{2\times N_{1/4}\times R}$ corresponding to the restriction of \(\mathbf s\) to its tail:
$$ p^a_{i,r}:=s^{T,a}_{N-N_{1/4}+i,r}, \qquad i=1,\,2,\,\dots,\,N_{1/4}. $$
It then follows that
\begin{align*}
    \CCDF_{\mathcal E}(y;\,\mathbf s^{T,a}_{\cdot,r}) &= 1-\frac{N-N_{1/4}}{N}-\frac{1}{N}\sum_{i=1}^{N_{1/4}}\chi_{[0,\,y]}(s^{T,a}_{N-N_{1/4}+i,r}) \\
    &= \frac{N_{1/4}}{N}\left[1-\frac{1}{N_{1/4}}\sum_{i=1}^{N_{1/4}}\chi_{[0,\,y]}(p^{T,a}_{i,r})\right] \\
    &= \frac{N_{1/4}}{N}\CCDF_{\mathcal E}(y;\,\mathbf p^a_{\cdot,r})
\end{align*}
for all $a$ and $r$. Thus, $\CCDF_{\mathcal E}(y;\,\mathbf p^a_{\cdot,r})$ and $\CCDF_{\mathcal E}(y;\,\mathbf s^{T,a}_{\cdot,r})$ differ only by the scaling factor $\frac{N_{1/4}}{N}$.
                
\begin{rmk}
Suppose that the city-size distribution exhibits a Pareto tail. Then, by~\eqref{eqComplementaryCumulativePareto}, $\CCDF(s)\approx\frac{c}{s^\beta}$, which, on a logarithmic scale, becomes $\log{\CCDF(s)}\approx-\beta\log{s}+\log{c}$. On a log--log plot, this corresponds to a straight line with slope $-\beta$ and intercept $\log{c}$. Now let $\tilde{\CCDF}:=k\CCDF$ be a rescaling of $\CCDF$ by a factor $k$. Then $\log{\tilde{\CCDF}(s)}=\log{\CCDF(s)}+\log{k}\approx-\beta\log{s}+\log{(kc)}$, which has the same slope as the original straight line. Hence, the Pareto index of the tail is invariant under multiplicative rescaling of the CCDF.
\label{ossLogarithmParetoCCDF}
\end{rmk}

In view of Remark~\ref{ossLogarithmParetoCCDF}, we may therefore plot $\CCDF_{\mathcal E}(y;\,\mathbf s^{T,a}_{\cdot,r})$ on a logarithmic scale to improve readability, since our primary interest lies not in the pointwise values of $\CCDF$, but in the parameter $\beta$. This, however, gives rise to a minor issue: the last element $p^a_{N_{1/4},r}$ has zero ordinate and is therefore not visible on a logarithmic scale. To circumvent this issue, for graphical purposes we shift the empirical CCDF on the linear scale by setting
$$ \CCDF_{\mathcal E}(y;\,\mathbf p^a_{\cdot,r})+\frac{1}{2N_{1/4}} $$
for all $a$ and $r$. On a logarithmic scale, this transformation results in a non-uniform shift of the values, which are therefore slightly distorted. The distortion is small, however, since $\frac{1}{2N_{1/4}}$ is small. This transformation is thus used solely for graphical purposes and not in the computation of the Pareto index $\beta$, which is always estimated directly from the tensor $\mathbf p$.

Overall, in the second and third rows of Figures~\ref{figSERStableReferenceConfiguration}--\ref{figSSERReferenceConfiguration} the empirical CCDF is plotted on a logarithmic scale using the scatter plot of the points
$$ \left(\frac{N}{N_{1/4}}\CCDF_{\mathcal E}(\bar{p}^a_i;\,\mathbf s^{T,a}_{\cdot,r})+\frac{1}{2N_{1/4}},\,\bar{p}^a_i\right)
    =\left(\CCDF_{\mathcal E}(\bar{p}^a_i;\,\mathbf p^a_{\cdot,r})+\frac{1}{2N_{1/4}},\,\bar{p}^a_i\right) $$
for all $a$ and $i$, where $\bar{p}^a_i$ is the sample mean associated with the confidence interval
$$  \IC_R^{\delta}(\mathbf p^a_i)=\left[\bar{p}^a_i-t^{\delta/2}_{R-1}\sqrt{\frac{v^a_i}{R}},\,\bar{p}^a_i+t^{\delta/2}_{R-1}\sqrt{\frac{v^a_i}{R}}\right]. $$

\subsubsection{Pareto CCDF and Pareto index}
\label{secParetoCCDF}
%
The second and third rows of Figures~\ref{figSERStableReferenceConfiguration}--\ref{figSSERReferenceConfiguration} also show the Pareto CCDF fitted over the tail of the city-size distribution.

To fit the Pareto distribution, we proceed analogously to Section~\ref{secEmpiricalCCDF}. Let $\CCDF_{\mathcal P}(\cdot;\,\mathbf p^a_{\cdot,r}):\R_+\to\R_+$ denote the CCDF of the Pareto distribution, cf.~\eqref{eqComplementaryCumulativePareto}, fitted to the vector $\mathbf p^a_{\cdot,r}$, and let $\CCDF_{\mathcal P}(\cdot;\,\mathbf p^a):\R_+\to\R_+$ be the vector-valued function defined by
$$ \CCDF_{\mathcal P}(y;\,\mathbf p^a):=
    \left[\CCDF_{\mathcal P}(y;\mathbf p^a_{\cdot,1}),\,\CCDF_{\mathcal P}(y;\,\mathbf p^a_{\cdot,2}),\,\dots,\,\CCDF_{\mathcal P}(y;\,\mathbf p^a_{\cdot,R})\right]^\top. $$
The corresponding confidence intervals are then given by
$$ \IC_R^{\delta}\!\left(\frac{N}{N_{1/4}}\CCDF_{\mathcal P}(y;\,\mathbf p^a)\right)
    =\left[\bar{\CCDF}_{\mathcal P}^a(y)-t^{\delta/2}_{R-1}\sqrt{\frac{v^a(y)}{R}},\,\bar{\CCDF}_{\mathcal P}^a(y)+t^{\delta/2}_{R-1}\sqrt{\frac{v^a(y)}{R}}\right],
    \qquad y\in [s_{3/4},\,s_{\max}], $$
where $\frac{N}{N_{1/4}}$ is the scaling factor introduced in Section~\ref{secEmpiricalCCDF}, which is required here to compare the fitted Pareto CCDF with the empirical CCDF.

As in Remark~\ref{ossBiLognormalParameterFitting}, confidence intervals are computed both for the fitted Pareto CCDF and for its parameter $\beta$. In particular, if $\boldsymbol\beta\in\R^{2\times R}$ denotes the matrix collecting the Pareto indices obtained from all simulations, then
$$ \IC_R^\delta(\boldsymbol\beta)=\left[\bar\beta^a-t^{\delta/2}_{R-1}\sqrt{\frac{v^a}{R}},\,\bar\beta^a+t^{\delta/2}_{R-1}\sqrt{\frac{v^a}{R}}\right] $$
for each $a$.
        
\subsubsection{Bimodal lognormal CCDF}
The third rows of Figures~\ref{figSERStableReferenceConfiguration}--\ref{figSSERReferenceConfiguration} also show the bilognormal CCDF fitted over the tail of the city-size distribution.

Let $\CCDF_{\mathcal L}(\cdot;\,\mathbf p^a_{\cdot,r}):\R_+\to\R_+$ and $\CCDF_{\mathcal B}(\cdot;\mathbf p^a_{\cdot,r}):\R_+\to\R_+$ denote the CCDFs of the lognormal and bilognormal distributions, respectively, fitted to the vector $\mathbf p^a_{\cdot,r}$. By linearity, it follows from~\eqref{eqBimodalLognormalDistribution} that
$$ \CCDF_{\mathcal B}(y;\,\mathbf p^a_{\cdot,r})=\zeta\CCDF_{\mathcal L}(y;\,\mathbf p^a_{\cdot,r})
    +(1-\zeta)\CCDF_{\mathcal L}(y;\,\mathbf p^a_{\cdot,r}). $$
Hence, analogously to Section~\ref{secParetoCCDF}, we can define the vector-valued function $\CCDF_{\mathcal B}(\cdot;\,\mathbf p^a):\R_+\to\R_+$ by
$$ \CCDF_{\mathcal B}(y;\,\mathbf p^a):=
    \left[\CCDF_{\mathcal B}(y;\,\mathbf p^a_{\cdot,1}),\,\CCDF_{\mathcal B}(y;\,\mathbf p^a_{\cdot,2}),\,\dots,\,\CCDF_{\mathcal B}(y;\,\mathbf p^a_{\cdot,R})\right]^\top. $$
The corresponding confidence intervals are then given by
$$ \IC_R^{\delta}\!\left(\frac{N}{N_{1/4}}\CCDF_{\mathcal B}(y;\,\mathbf p^a)\right)
    =\left[\bar{\CCDF}_{\mathcal B}^a(y)-t^{\delta/2}_{R-1}\sqrt{\frac{v^a(y)}{R}},\,
        \bar{\CCDF}_{\mathcal B}^a(y)+t^{\delta/2}_{R-1}\sqrt{\frac{v^a(y)}{R}}\right],
            \qquad y\in [s_{3/4},\,s_{\max}] $$
for each $a$.
        
\subsubsection{Data transformation}
The plot in the second to fourth rows of Figures~\ref{figSERStableReferenceConfiguration}--\ref{figSSERReferenceConfiguration} are shown on a semilogarithmic scale. This choice is motivated by the fact that a lognormally distributed dataset $\mathbf s$ is transformed into a normally distributed dataset $\mathbf h\in\R_+^{N_f\times2\times N\times R}$ by the transformation
\begin{equation}
    h_{i,r}^{n,a}:=\log_{10}{s_{i,r}^{n,a}},
    \label{eqlog10transformation}
\end{equation}    
which facilitates interpretation, particularly when the parameters are varied, as in parametric studies presented in the Supplementary Material.

Under the transformation~\eqref{eqlog10transformation}, the histogram bins are obtained as $\mathbf b^a_{\cdot,r}:=\mathcal{B}(\mathbf h^{T,a}_{\cdot,r})$, and the confidence intervals are transformed accordingly. Moreover, the probability density $\bar f_N(s):=\bar f_N(s,T)$ is mapped to a new density, denoted by $\bar l_N=\bar l_N(s)$, while preserving the probability mass:
$$ \bar f_N(s)\,ds=\bar l_N(h)\,dh. $$
Since $h=\log_{10}{s}$, it follows that
$$ \bar l_N(h)=10^h\log{10}\bar f_N(10^h). $$

Note that the logarithm in~\eqref{eqLognormalDistribution} is meant in the natural sense, whereas the transformation~\eqref{eqlog10transformation} adopts base 10. The two scales are related by the rescaling factor $\log{10}\approx2.303$, which is the very same Jacobian factor carried by $\bar l_N(h)$ above. Consequently, if the logarithm in~\eqref{eqLognormalDistribution} is interpreted in base 10, then the parameters must be divided by $\log{10}$, i.e. $\mu\mapsto\mu/\log{10}$ and $\sigma\mapsto\sigma/\log{10}$, whose transformation changes the prefactor of~\eqref{eqLognormalDistribution} to $\left(x\sigma\sqrt{2\pi}\,\log{10}\right)^{-1}$. The $\mu$ and $\sigma$ of the fits reported in this work are therefore base-10 parameters and are not directly comparable with those of~\eqref{eqLognormalDistribution}.

It is this transformed density $\bar l_N$ that is plotted in the first rows of Figures~\ref{figSERStableReferenceConfiguration}--\ref{figSSERReferenceConfiguration}.

By contrast, the empirical CCDFs remain unchanged in value, with their domain simply mapped from $[s_{\min},\,s_{\max}]$ to $[\log_{10}{s_{\min}},\,\log_{10}{s_{\max}}]$. This follows from the fact that the logarithm is strictly increasing and therefore preserves the ordering of the data. In particular:
$$ \operatorname{Prob}(S<s)=\operatorname{Prob}(\log_{10}{S}<\log_{10}{s}), \qquad \forall\,s\in\R_+. $$

Finally, we observe that, on a semilogarithmic scale, the relationship discussed in Remark~\ref{ossLogarithmParetoCCDF} takes the form $\log_{10}{\tilde\CCDF(h)}\approx-\beta h+\log{(kc)}$, which remains linear in $h$.

\subsubsection{Mean city size}
We compute the mean city-size tensor $\bar{\mathbf s}\in\R_+^{N_f\times2\times R}$ as
$$ \bar s^{n,a}_r:=\frac{1}{N}\sum_{i=1}^Ns^{n,a}_{i,r} $$
and the corresponding confidence intervals as
$$ \IC_R^{\delta}(\bar{\mathbf s}^{n,a})=\left[\bar{\bar s}^{n,a}-t^{\delta/2}_{R-1}\sqrt{\frac{v^{n,a}}{R}},\,\bar{\bar s}^{n,a}+t^{\delta/2}_{R-1}\sqrt{\frac{v^{n,a}}{R}}\right] $$
for all $n$ and $a$.

The evolution of the ensemble-averaged mean city size is then plotted by linearly interpolating the points $\{(t^f_n,\,\bar{\bar s}^{n,a})\}_{n,\,a}$ between consecutive time instants, together with confidence band $\pm t^{\delta/2}_{R-1}\sqrt{\frac{v^{n,a}}{R}}$.
        
\subsubsection{Mean city size against latent variable}
\label{secMeanCitySizevsLatentVariable}
%
The fourth rows of Figures~\ref{figSERStableReferenceConfiguration}--\ref{figSSERReferenceConfiguration} show scatter plots of the mean city size against the city latent variable (for instance, the city degree).

To this end, we compute the mean city size by averaging over the ensemble of simulations:
$$ \bar s^a_i:=\frac{1}{R}\sum_{r=1}^Rs^{T,a}_{i,r} $$
for all $a$ and $i$, and then plot the set of points $\{(x_i,\,\bar s^a_i)\}_{a,\,i}$.

In this case, we do not show the corresponding confidence intervals, as the they would make the plots, which are already rather dense, difficult to read.
        
\subsubsection{Evolutions of the mean city size by degree class}
These are the plots shown in the fifth rows of Figures~\ref{figSERStableReferenceConfiguration}--\ref{figSSERReferenceConfiguration}.

Let $\mathcal X\subseteq\{1,\,2,\,\dots,\,N\}$ denote the set of distinct degrees in the graph and, let us continue to denote by $\bar{\mathbf s}\in\R^{N_f\times 2\times\vert\mathcal X\vert}$ the tensor whose entries are defined by
$$ \bar s^{n,a}_k:=\frac1{\vert\mathcal I_k\vert}\sum_{i\in\mathcal I_k}\frac{1}{R}\sum_{r=1}^Rs^{n,a}_{i,r} $$
for all $n$, $a$, and $k$, where $\mathcal I_k:=\{i\in\mathcal I\,:\,x_i=k\}\subseteq\mathcal I$ denotes the set of indices of the vertices of degree $k$.

By linearly interpolating the points $\{(t^f_n,\,\bar s^{n,a}_k)\}_{n,\,a,\,k}$ between consecutive time instants, we represent the time evolution of the mean population size of cities associated with vertices of degree $k$, averaged over all simulations. We use this plot to verify that the system has reached a steady state by the prescribed final simulation time $T$.
               
As in Section~\ref{secMeanCitySizevsLatentVariable}, we do not show the corresponding confidence intervals, for the same reason given there.

\section*{CRediT authorship contribution statement}
Valerio Taralli: Conceptualization, Methodology, Software, Formal analysis, Investigation, Writing - Original Draft. Andrea Tosin: Supervision.

\bibliographystyle{../Corpo/cas-model2-names}
\bibliography{../Corpo/cas-refs}
\end{document}

%#region Bozza
  % Tutto ciò che segue è la prima versione dell'articolo 
  % Essa è conservata solo come riferimento storico

  % !LW recipe=2xpdfLaTeX
  % !LW recipe=pdfLaTeX
  % !LW recipe=pdfLaTeX ➞ bibTeX ➞ pdfLaTeX

  %% Copyright 2019-2024 Elsevier Ltd
  %% 
  %% This file is part of the 'CAS Bundle'.
  %% --------------------------------------
  %% 
  %% It may be distributed under the conditions of the LaTeX Project Public
  %% License, either version 1.3c of this license or (at your option) any
  %% later version.  The latest version of this license is in
  %%    http://www.latex-project.org/lppl.txt
  %% and version 1.3c or later is part of all distributions of LaTeX
  %% version 1999/12/01 or later.
  %% 
  %% The list of all files belonging to the 'CAS Bundle' is
  %% given in the file `manifest.txt'.
  %% 
  %% Template article for cas-sc documentclass for 
  %% double column output.

  % \documentclass[a4paper,fleqn]{cas-sc}

  % 
%\begingroup For image generation 
  \usepackage[strict]{changepage}

  \usepackage{shellesc}

  \usepackage{etoolbox}
  \newbool{forcetikzrecompile}

  \usepackage{tikz}
  \usetikzlibrary{
    positioning,
    shapes.geometric,
    plotmarks,
    arrows.meta,
    calc,
    decorations.pathreplacing,
    backgrounds,
    patterns.meta,
    fpu, % «floating point unit»
    % external,
  } 

  \usepackage{pgfplots}
  \pgfplotsset{compat=newest}
  \usepgfplotslibrary{
    patchplots,
    groupplots,
    colorbrewer
  }
  \pgfdeclarelayer{background}
  \pgfdeclarelayer{foreground}
  \pgfsetlayers{background,main,foreground}
%\endgroup

%\begingroup For the paper
  \usepackage{microtype}

  \usepackage{amsmath}
  \usepackage{amsthm}
  \usepackage{accents}

  \usepackage{zref-savepos}
  \newcounter{figposcount}

  \usepackage{xstring}

  %\usepackage[numbers]{natbib}
  %\usepackage[authoryear]{natbib}
  % \usepackage[authoryear,longnamesfirst]{natbib}
  \usepackage[numbers,square]{natbib}

  \usepackage{
    graphicx,
    caption,
    subcaption,
    stfloats,
    float
  }

  \usepackage[%
    ruled,%
    vlined,%
    linesnumbered%
  ]{algorithm2e}

  \usepackage{ragged2e}

  \usepackage{enumitem}
  \setlist[enumerate]{%
    label=\arabic*.,
    topsep=0.5em,%
    parsep=0em,%
    itemsep=0.2em,%
    % leftmargin=2em,%
    rightmargin=1.5em,%
  }
  \setlist[itemize]{%
    label=\(\diamond\),
    topsep=0.5em,%
    parsep=0em,%
    itemsep=0.2em,%
    % leftmargin=2em,%
    rightmargin=1.5em,%
  }

  \usepackage{hyperref}
  \usepackage{xurl}

  \usepackage[english]{cleveref}

  \usepackage{siunitx}
  \sisetup{%
    exponent-mode=threshold,%
    exponent-thresholds=-3:3,%
    round-mode=places,% figures
    round-precision=3,%
    round-pad=false,%
    separate-uncertainty=true,%
    print-unity-mantissa=false,%
    tight-spacing=true,%
  }
%\endgroup

  % 
\newif\ifcasscclass
\makeatletter
\@ifclassloaded{cas-sc}{\casscclasstrue}{\casscclassfalse}
\makeatother

\newif\ifpageisodd % È usato (nelle macro delle infopagini) dalla meccanica

\ifcasscclass
%\begingroup For image generation
  %\begingroup Tipografia
    \newcommand\textsize[1]{%
      \pgfmathparse{#1*1.2}%
      \fontsize{#1pt}{\pgfmathresult pt}\selectfont%
    }

    \DeclareMathOperator\FRC{\mathcal R}

    % (il rilevamento della classe è in cima al file)

    \providecommand\Ifthispageodd[2]{\checkoddpage\ifoddpage#1\else#2\fi}
  %\endgroup
\fi

  %\begingroup TikZ
    %\begingroup Didascalie
      \newcommand\newcaption[2]{%
        \captionsetup{width=#2} % aboveskip=-.2cm,
        \caption{#1}%
      }
      
      \newcommand\newsubcaption[2][ ]{%
        \refstepcounter{subfigure}%
        (\thesubfigure)#1#2%
      }
    %\endgroup
    
\ifcasscclass
    %\begingroup RedefineTikZbounds
      % La differenza è che colla prima bisogna catturare anche le etichette perché si crea un pdf isolato mentre colla non è importante perché viene tutto compilato sul momento

      %\begingroup Versione valida con esternalizzazione [ma anche senza]
        \newcommand\redefineTikZbounds[2]{
          \coordinate (RealSW) at (current bounding box.south west);
          \coordinate (RealNE) at (current bounding box.north east);
          \coordinate (Center) at ($(#1.west)!0.5!(#2.east)$);
    
          \pgfresetboundingbox 

          \useasboundingbox 
            let \p1=(RealSW),\p2=(RealNE),\p3=(Center),
            \n1={max(\x3-\x1,\x2-\x3)} in
            (\x3-\n1,\y1) rectangle (\x3+\n1,\y2+.25em);

          % \path[use as bounding box] 
          %   let \p1=(RealSW),\p2=(RealNE),\p3=(Center),
          %   \n1={max(\x3-\x1,\x2-\x3)} in
          %   (\x3-\n1,\y1) rectangle (\x3+\n1,\y2+.25em);
          % Lo «.25em» discende dalla traslazione imposta all'etichetta della scala verticale mediante «yScaleLabelShift»
        } % Scatola delimitante simmetrica
      %\endgroup

      %\begingroup Versione valida senza esternalizzazione
        % \newcommand\redefineTikZbounds[2]{
        %   % Salva le coordinate della cima e del fondo della figura corrente (cosí da mantenere i titoli e le didascalie)
        %   \coordinate (Cima) at (current bounding box.north);
        %   \coordinate (Fondo) at (current bounding box.south);
    
        %   \pgfresetboundingbox % Resetta la scatola dei limiti
    
        %   \useasboundingbox % Definisce la nuova scatola come un rettangolo
        %     (#1.west |- Cima)  % con ascissa sinistra data da «Col1 c1r1.west», 
        %     rectangle      % ascissa destra data da «Col2 c1r1.east»
        %     (#2.east |- Fondo); % ordinate date dalla cima e fondo originali
        %     % L'operatore «|-» in «(A |- B)» prende come ascissa quella del punto A mentre come coordinata quella del punto B
        % }
      %\endgroup
    %\endgroup

    %\begingroup TikZBox
      % Parzialmente ispirato da https://gist.github.com/Caramdir/899107

      \makeatletter
      \newsavebox\tikzBufferBox
      \xdef\offsetX{0pt}
      \xdef\offsetY{0pt}
      
      % \tikzbox[#1]{#2}{#3}[#4]{#5}
      % #1 -> Ancora rispetto cui posizionare l'intera scatola (di base è «south west», ossia l'angolo in basso a sinistra)
      % #2 -> Nodo dove posizionare l'origine (per la natura di come LaTeX riceve gli argomenti è buon pratica passare questo circondato da parentesi graffe)
      % #3 -> Nome della scatola limitante definita dall'ambiente «scope»
      % #4 -> Argomenti opzionali da passare all'ambiente «scope»; di base sono vuoti
      % #5 -> Contenuto della scatola circondato dall'ambiente «scope»
      % Es. \tikzbox{{(0,0)}}{figuraTikZ}{blocco TikZ}

      % Macro di apertura che controlla se v'è l'argomento facoltativo dell'ancora
      \def\tikzbox{\pgfutil@ifnextchar[{\tikzbox@anchor}{\tikzbox@anchor[south west]}}
      % La funzione «\pgfutil@ifnextchar» serve per controllare se il prossimo carattere è l'apertura di una paretensi graffa «[»
      
      % Macro intermedia per cogliere gli argomenti obbligatori e controllare quelli successivi facoltativi
      \def\tikzbox@anchor[#1]#2#3{\pgfutil@ifnextchar[{\tikzbox@opt[#1]#2#3}{\tikzbox@opt[#1]#2#3[]}}

      % Macro degli argomenti obbligatori
      \long\def\tikzbox@opt[#1]#2#3[#4]#5{%
        \node[ % Misurazione della scatola in un nodo invisibile
          overlay,%
          opacity=0,%
          inner sep=0pt,%
          outer sep=0pt%
        ] {%
          \sbox{\tikzBufferBox}{%
            \begin{tikzpicture}%
              \begin{scope}[local bounding box=measuringTikZBox]%
                #5%
              \end{scope}%
              \pgfpointanchor{measuringTikZBox}{#1}%
              \pgfgetlastxy{\tempX}{\tempY}% 
              \xdef\offsetX{\tempX}%
              \xdef\offsetY{\tempY}%
            \end{tikzpicture}%
          }%
        };%
        %
        % Disegno reale traslato
        \begin{scope}[%
          local bounding box=#3,%
          shift={#2},%
          xshift=-\offsetX,%
          yshift=-\offsetY,%
          #4%
        ]%
          #5%
        \end{scope}%
      }%
      \makeatother
    %\endgroup

    %\begingroup Chiavi-Valori
      %\begingroup Testo
        % La dimensione del testo di base è 11pt
        \newcommand\titleSize{11}
        \newcommand\subCaptionSize{10}
        \newcommand\ParameterSize{9}
        \newcommand\labelSize{10}
        \newcommand\tickSize{9}
        \newcommand\legendSize{9}
      %\endgroup
    
      %\begingroup Infopagine
        %\begingroup Chiavi condivise
        \newcounter{lpc} % «loopcounter»
        \newcommand\infopageSharedKeys{
          %\begingroup Testo
            \renewcommand\titleSize{10}
            \renewcommand\subCaptionSize8
            \renewcommand\ParameterSize7
            \renewcommand\labelSize9
            \renewcommand\tickSize6
            \renewcommand\legendSize7
          %\endgroup
    
          %\begingroup Dimensioni
            \pgfmathsetmacro\plotWidth7
            \pgfmathsetmacro\plotHeight{.618*\plotWidth}
            \pgfmathsetmacro\halfPlotHeight{\plotHeight/2}
            %
            \pgfmathparse{0.825*\plotWidth}
            \pgfmathsetmacro\plotWidth\pgfmathresult
            \pgfmathsetmacro\halfWidth{\plotWidth/2}
            \pgfmathsetmacro\plotLastWidth{1.25*\plotWidth}
    
            \pgfmathsetmacro\xPlotSep{1}
            \pgfmathsetmacro\yPlotSep{1.25}
            \pgfmathsetmacro\yGroupSep{2.75}
          %\endgroup
    
          %\begingroup Legende
            \newcommand\placeLegend[5]{%
              /pgfplots/plotLegend,
              at={([shift={(##4,##5)}]axis description cs:##1,##2)},
              anchor=##3,
            }
            %
            \pgfplotsset{%
              plotLegend/.style={%
                font=\textsize{\legendSize},
                /pgfplots/legend cell align=left,
                % align=left,
                % draw=white!15!black,
                draw=none, % Bordo invisibile
                fill=white,
                fill opacity=0.5,
                text opacity=1,
                % text width=2cm,
              },
              plotLegendNW/.style={legend style={\placeLegend{0}{1}{north west}{0em}{0em}}},
              plotLegendNE/.style={legend style={\placeLegend{1}{1}{north east}{0em}{0em}}},
              plotLegendSW/.style={legend style={\placeLegend{0}{0}{south west}{0em}{0em}}},
              plotLegendSE/.style={legend style={\placeLegend{1}{0}{south east}{0em}{0em}}},
              %
              % plotLegendW/.style={legend style={plotLegend,at={(0,.5)},anchor=west}},
              %
              % plotLegendNW/.style={legend style={plotLegend,legend pos=north west}},
              % plotLegendNE/.style={legend style={plotLegend,legend pos=north east}},
              % plotLegendSW/.style={legend style={plotLegend,legend pos=south west}},
              % plotLegendSE/.style={legend style={plotLegend,legend pos=south east}},
              %
              % plotLegendNWcloser/.style={legend style={plotLegend,at={(0,1)},anchor=north west}},
              % plotLegendNEcloser/.style={legend style={plotLegend,at={(1,1)},anchor=north east}},
              % plotLegendSWcloser/.style={legend style={plotLegend,at={(0,0)},anchor=south west}},
              % plotLegendSEcloser/.style={legend style={plotLegend,at={(0,1)},anchor=south east}},
            }
            %
            \pgfmathsetmacro\lineLegendLength{1}
            \pgfmathsetmacro\areaLegendHeigth{.25}
            \pgfplotsset{
              legend image code/.code={
                \draw [mark repeat=2,mark phase=2,##1]
                  plot coordinates {
                    (0em,0em)
                    (.5*\lineLegendLength em,0em)
                    (\lineLegendLength em,0em)
                  };
              },
              /pgfplots/area legend/.style={
                legend image code/.code={
                  \draw [##1,fill,draw=none] 
                    (0em,-\areaLegendHeigth em) 
                    rectangle 
                    (\lineLegendLength em,\areaLegendHeigth em);
                },
              },
            }
            %
            \newcommand\ExactHistogramLegend{Exact histogram}
            \newcommand\ApproxHistogramLegend{Approx. histogram}
            \newcommand\RealHistogramLegend{Real histogram}
            %
            \newcommand\ExactBLNFitLegend{Exact BLN fit}
            \newcommand\ApproxBLNFitLegend{Approx BLN fit}
            \newcommand\RealBLNFitLegend{Real BLN fit}
            %
            \newcommand\ExactEmpiricCCDFLegend{Exact empiric CCDF}
            \newcommand\ApproxEmpiricCCDFLegend{Approx. empiric CCDF}
            \newcommand\RealEmpiricCCDFLegend{Real empiric CCDF}
            %
            \newcommand\ExactParetoFitLegend{Exact Pareto fit}
            \newcommand\ApproxParetoFitLegend{Approx. Pareto fit}
            \newcommand\RealParetoFitLegend{Real Pareto fit}
            %
            \newcommand\ExactScatterLegend{Exact scatter}
            \newcommand\ApproxScatterLegend{Approx. scatter}
            \newcommand\RealScatterLegend{Real scatter}
            %
            \newcommand\ClasskILegend{Class k=283}
            \newcommand\ClasskIILegend{Class k=83}
            \newcommand\ClasskIIILegend{Class k=43}
            \newcommand\ClasskIVLegend{Class k=10}
          %\endgroup
    
          %\begingroup Etichette vuote
            \pgfplotsset{
              yLabelEmpty/.style={ylabel={}},
              yTickLabelsEmpty/.style={yticklabels={}},
              yTickScaleLabelEmpty/.style={ytick scale label code/.code={}},
              xTickScaleLabelEmpty/.style={xtick scale label code/.code={}},
              emptyLeged/.style={
                legend style={draw=none,fill=none,nodes={draw=none}},
                legend image post style={opacity=0},
              },
            }
          %\endgroup
    
          %\begingroup Estetica
            \newcommand\smallminus{\scalebox{0.5}[1]{$-$}\hspace{-.11em}}
            \pgfplotsset{
              /pgfplots/log number format basis/.code 2 args={%
                \pgfmathsetmacro\exponent{##2}%
                \pgfmathparse{sign(\exponent)}%
                \ifnum\pgfmathresult<0%
                  \pgfmathparse{abs(\exponent)}%
                  $##1^{\smash{\smallminus\pgfmathprintnumber{\pgfmathresult}}}$%
                \else%
                  $##1^{\smash{\pgfmathprintnumber{\exponent}}}$%
                \fi%
              },
              ytick scale label code/.code={%
                \pgfmathsetmacro\exponent{##1}%
                \pgfmathparse{sign(\exponent)}%
                \ifnum\pgfmathresult<0
                  \pgfmathparse{abs(\exponent)}%
                  $\cdot 10^{\smash{\smallminus\pgfmathprintnumber{\pgfmathresult}}}$%
                \else
                  $\cdot 10^{\smash{\pgfmathprintnumber{\exponent}}}$%
                \fi
              },
            }
          %\endgroup
    
          %\begingroup Colori
            \definecolor{blueish}{rgb}{0.00000,0.44700,0.74100}%
            \definecolor{reddish}{rgb}{0.83529,0.36863,0.00000}%
          %\endgroup
    
          %\begingroup Commenti
            %\begingroup Impostazioni
            \tikzset{
              title/.style={
                align=center,
                anchor=south,
                font=\bfseries\textsize{\titleSize},
              },
              fittedParameters/.style={
                align=center,
                font=\textsize{\ParameterSize},
              },
              varyingParameter/.style={
                align=center,
                % rotate=90,
                font=\textsize{\subCaptionSize},
              },
              subCaption/.style={
                % text width=\plotWidth cm,
                align=center,
                font=\textsize{\subCaptionSize},
              },
            }
            \newcommand\codetocolor[1]{\ifcase##1blueish\or reddish\else black\fi}
            %\endgroup
    
            %\begingroup Parametri Sardegna
            \newcommand\betaReal{\num{1.273}}
            \newcommand\zetaReal{\num{0.929}}
            \newcommand\mIReal{\num{3.227}}
            \newcommand\sIReal{\num{0.408}}
            \newcommand\mIIReal{\num{4.099}}
            \newcommand\sIIReal{\num{0.556}}
            %\endgroup
    
            %\begingroup Parametri Pareto
            \newcommand\codetoParetovalue[1]{
              \ifcase##1\betaExact\or\betaAprx\else\betaReal\fi
            }
            \newcommand\writeParetoParameters[7][center]{
              \node[fittedParameters,anchor=##1] at ([shift={(##5,##6)}]##4) {%
                \(\begin{aligned}
                  \color{\codetocolor{##2}}\beta
                  &\color{\codetocolor{##2}}=\codetoParetovalue{##2}\\[##7]
                  \color{\codetocolor{##3}}\beta
                  &\color{\codetocolor{##3}}=\codetoParetovalue{##3}
                \end{aligned}\)
              };
            }
            %\endgroup

            %\begingroup % Vecchia versione in cui si scrivevano esplicitamente tutte le 9 combinazioni possibili
            % \newcommand\writeParameters[6]{%
            %   \node[parameter] at ([shift={(##4,##5)}]##1) {%
            %     \ifcase##2% Caso 0
            %       \ifcase##3% Caso 0
            %         \(\begin{aligned}
            %           \color{blueish}\beta&\color{blueish}=\betaExact\\[##6]
            %           \color{blueish}\beta&\color{blueish}=\betaExact
            %         \end{aligned}\)
            %       \or% Caso 1
            %         \(\begin{aligned}
            %           \color{blueish}\beta&\color{blueish}=\betaExact\\[##6]
            %           \color{reddish}\beta&\color{reddish}=\betaAprx% Caso 1
            %         \end{aligned}\)
            %       \else% Caso 2
            %         \(\begin{aligned}
            %           \color{blueish}\beta&\color{blueish}=\betaExact\\[##6]
            %           \color{black}\beta&\color{black}=\betaReal% Caso 2
            %         \end{aligned}\)
            %       \fi
            %     \or% Caso 1
            %       \ifcase##3% Caso 0
            %         \(\begin{aligned}
            %           \color{reddish}\beta&\color{reddish}=\betaAprx\\[##6]
            %           \color{blueish}\beta&\color{blueish}=\betaExact
            %         \end{aligned}\)
            %       \or% Caso 1
            %         \(\begin{aligned}
            %           \color{reddish}\beta&\color{reddish}=\betaAprx\\[##6]
            %           \color{reddish}\beta&\color{reddish}=\betaAprx
            %         \end{aligned}\)
            %       \else% Caso 2
            %         \(\begin{aligned}
            %           \color{reddish}\beta&\color{reddish}=\betaAprx\\[##6]
            %           \color{black}\beta&\color{black}=\betaReal
            %         \end{aligned}\)
            %       \fi
            %     \else% Caso 2
            %       \ifcase##3% Caso 0
            %         \(\begin{aligned}
            %           \color{black}\beta&\color{black}=\betaReal\\[##6]
            %           \color{blueish}\beta&\color{blueish}=\betaExact
            %         \end{aligned}\)
            %       \or% Caso 1
            %         \(\begin{aligned}
            %           \color{black}\beta&\color{black}=\betaReal\\[##6]
            %           \color{reddish}\beta&\color{reddish}=\betaAprx
            %         \end{aligned}\)
            %       \else% Caso 2
            %         \(\begin{aligned}
            %           \color{black}\beta&\color{black}=\betaReal\\[##6]
            %           \color{black}\beta&\color{black}=\betaReal
            %         \end{aligned}\)
            %       \fi
            %     \fi
            %   };
            % }
            %\endgroup
    
            %\begingroup Parametri bilognormale
            \newcommand\codetoparameter[1]{
              \ifcase##1\zeta  % Caso 0
              \or\mu_1     % Caso 1
              \or\sigma_1   % Caso 2
              \or\mu_2     % Caso 3
              \else\sigma_2\fi % Caso 4
            }
            \newcommand\codetoBiLognormalvalue[2]{
              \ifcase##1 % Caso 0
                \ifcase##2\zetaExact % Caso 0
                \or\mIExact    % Caso 1
                \or\sIExact    % Caso 2
                \or\mIIExact    % Caso 3
                \else\sIIExact\fi % Caso 4
              \or % Caso 1
                \ifcase##2\zetaAprx % Caso 0
                \or\mIAprx     % Caso 1
                \or\sIAprx     % Caso 2
                \or\mIIAprx    % Caso 3
                \else\sIIAprx\fi  % Caso 4
              \else % Caso 2
                \ifcase##2\zetaReal % Caso 0
                \or\mIReal     % Caso 1
                \or\sIReal     % Caso 2
                \or\mIIReal    % Caso 3
                \else\sIIReal\fi  % Caso 4
              \fi
            }
            \newcommand\matrixline[2]{\codetoparameter{##2}&=\codetoBiLognormalvalue{##1}{##2}}
            \def\lineloop##1##2{%
              \matrixline{##1}{\number\value{lpc}}\stepcounter{lpc}%
              \ifnum\number\value{lpc}<5\\[##2]\lineloop{##1}{##2}\fi%
            } % https://tex.stackexchange.com/questions/491358/tikz-foreach-to-work-with-the-aligned-environment
            \newcommand\writeBiLognormalParameters[6][center]{
              \node[fittedParameters,anchor=##1] at ([shift={(##4,##5)}]##3) {%
                \setcounter{lpc}{0}
                \(\color{\codetocolor{##2}}
                \begin{aligned}
                  \lineloop{##2}{##6}
                \end{aligned}\)
              };
            }
            %\endgroup
    
            %\begingroup Altre etichette
            \newcommand\writeSizeEvolutionDescriptions[1][5]{%
              \foreach \col/\clr/\txt in {%
                1/blueish/Esatto,
                2/reddish/Approssimato
              }{
                \node[
                  fittedParameters,
                  anchor=south,
                  yshift=.1em,
                ] at (Row##1 c\col r1.south) {\color{\clr}\txt};
              }
            }
    
            \newcommand\writeVaryingParameter[5]{%
              \foreach \r in {1,...,5}{
                \node[
                  varyingParameter,
                  anchor=##3,
                ] at ([shift={(##4,##5)}]Row\r\space c##1r1.##2)
                {\(\varyingParameterSymbol=\varyingParameterList\r\)};
              }
            }
            %\endgroup
          %\endgroup
    
          %\begingroup Sottodidascalie
            \newcommand\writeSubCaption[5]{%
              \node[subCaption,anchor=##3]
              at ([shift={(##4,##5)}]##1.##2)
              {\newsubcaption[]{}};
            }
            \newcommand\writeSubCaptions[5][3]{%
              \foreach \r in {1,...,5}{
                \foreach \c in {1,...,##1}{
                  \writeSubCaption{Row\r\space c\c r1}{##2}{##3}{##4}{##5}
                }
              }
            }
          %\endgroup
        }
        %\endgroup
    
        %\begingroup Chiavi del primo tipo d'infopagina
        \newcommand\infopageKeysI{
          \infopageSharedKeys
    
          %\begingroup Gruppi
            \ifpageisodd
              \pgfplotsset{/pgfplots/group/yDescPos/.style={y descriptions at=edge right}}
            \else
              \pgfplotsset{/pgfplots/group/yDescPos/.style={y descriptions at=edge left}}
            \fi
            \pgfplotsset{
              /pgfplots/group/plotGroup/.style={
                horizontal sep=\xPlotSep em,
                vertical sep=\yPlotSep em,
                yDescPos,
                group size=3 by 1,
              },
              /pgfplots/group/plotLastGroup/.style={
                plotGroup,
                x descriptions at=edge bottom,
              },
            }
          %\endgroup
    
          %\begingroup Stili
            %\begingroup Selezione «recto»/«verso»
            \ifpageisodd
              \pgfplotsset{%
                yLineAxis/.style={axis y line*=right},
                /tikz/yLabelShift/.style={%
                  yshift=-1.25em,
                  at={(axis description cs:1,0.5)},
                  anchor=north,
                  % rotate=180,
                },
                /tikz/yScaleLabelShift/.style={%
                  yshift=.25em,
                  at={(1,1)},
                  anchor=south east,
                },
                /tikz/yTickLabelShift/.style={xshift=-0.15em},
              }
            \else
              \pgfplotsset{%
                yLineAxis/.style={axis y line*=left},
                /tikz/yLabelShift/.style={%
                  yshift=1.25em,
                  at={(axis description cs:0,0.5)},
                  anchor=south,
                  % rotate=180,
                },
                /tikz/yScaleLabelShift/.style={%
                  yshift=.25em,
                  at={(0,1)},
                  anchor=south west,
                },
                /tikz/yTickLabelShift/.style={xshift=+0.15em},
              }
            \fi
            %\endgroup
            %
            %\begingroup Grafici
            \pgfplotsset{
              plotRow/.style={
                scale only axis,
                width=\plotWidth cm,height=\plotHeight cm,
                xlabel style={font=\color{white!15!black}\textsize{\labelSize}},
                ylabel style={font=\color{white!15!black}\textsize{\labelSize},yLabelShift},
                xticklabel style={font=\textsize{\tickSize}},
                yticklabel style={font=\textsize{\tickSize},yTickLabelShift},
                y tick scale label style={yScaleLabelShift},
                % axis x line*=bottom,yLineAxis,
                enlarge x limits=0.1,enlarge y limits=0.1,
                xmajorgrids,ymajorgrids,
                xminorticks=true,yminorticks=true,
                scale ticks above exponent=2,
                axis background/.style={fill=white},
                grid style={dashed}
              },
              % xlabel123/.style={xlabel={\(h\)}},
              xlabel3/.style={xlabel={\(h\)}},
              ylabel23/.style={ylabel={\(\FRC(h)\)}},
              row1Keys/.style={
                % ylabel={\(\underhat{\bar l}_N(h)\)},
                ylabel={\(\bar l_N(h)\)},
                % enlarge x limits={lower,value=.05},
                enlarge x limits=.05,
                enlarge y limits={upper,value=.05},
                xmode=linear,
                log origin=infty,
              },
              row2KeysWithLabels/.style={
                ylabel23,
                xmode=linear,ymode=log
              },
              row2Keys/.style={
                ylabel23,
                xmode=linear,ymode=log,
                xticklabels={},
              },
              row3Keys/.style={
                xlabel3,ylabel23,
                xmode=linear,ymode=log
              },
              row4Keys/.style={
                xlabel={\(\xi\)},ylabel={\(s(\xi)\)},
                xmode=log,ymode=log
              },
              row5KeysLin/.style={
                width=\plotLastWidth cm,
                xlabel={\(t\)},ylabel={\(s(t)\)},
                enlarge x limits=0,
                enlarge y limits=0.05,
                x tick scale label style={
                  yshift=-1em,
                  at={(.95,0)},
                  anchor=north,
                },
              },
              row5Keys/.style={
                row5KeysLin,
                ymode=log,
              },
            }
            %\endgroup
          %\endgroup
        }
        %\endgroup
    
        %\begingroup Chiavi del secondo tipo d'infopagina
        \newcommand\infopageKeysII{
          \infopageSharedKeys
    
          %\begingroup Gruppi
            \ifpageisodd
              \pgfplotsset{/pgfplots/group/yDescPos/.style={y descriptions at=edge right}}
            \else
              \pgfplotsset{/pgfplots/group/yDescPos/.style={y descriptions at=edge left}}
            \fi
            \pgfplotsset{
              /pgfplots/group/plotGroup/.style={
                horizontal sep=\xPlotSep em,
                vertical sep=\yPlotSep em,
                y descriptions at=all,
                group size=3 by 1,
              },
              /pgfplots/group/plotLastGroup/.style={
                plotGroup,
                x descriptions at=edge bottom,
              },
            }
          %\endgroup
    
          %\begingroup Stili
            %\begingroup Selezione «recto»/«verso»
            \ifpageisodd
              \pgfplotsset{
                yLabel/.style={
                  ylabel style={
                    font=\color{white!15!black}\textsize{\labelSize},
                    yshift=-1.25em,
                    at={(axis description cs:1,0.5)},
                    anchor=north,
                    % rotate=180,
                  },
                },
                yLabelInv/.style={yLabelEmpty},
                yScaleLabelShift/.style={
                  y tick scale label style={
                    yshift=.25em,
                    % xshift=1em,
                    at={(1,1)},
                    anchor=south east,
                  },
                },
                yScaleLabelShiftInv/.style={
                  y tick scale label style={
                    yshift=.25em,
                    % xshift=-1em,
                    at={(0,1)},
                    anchor=south west,
                  },
                },
                % yLineAxis/.style={axis y line*=right},
              }
            \else
              \pgfplotsset{
                yLabel/.style={yLabelEmpty},
                yLabelInv/.style={
                  ylabel style={
                    font=\color{white!15!black}\textsize{\labelSize},
                    yshift=1.25em,
                    at={(axis description cs:0,0.5)},
                    anchor=south,
                    % rotate=180,
                  },
                },
                yScaleLabelShift/.style={
                  y tick scale label style={
                    yshift=.25em,
                    % xshift=1em,
                    at={(1,1)},
                    anchor=south east,
                  },
                },
                yScaleLabelShiftInv/.style={
                  y tick scale label style={
                    yshift=.25em,
                    % xshift=-1em,
                    at={(0,1)},
                    anchor=south west,
                  },
                },
                % yLineAxis/.style={axis y line*=left},
              }
            \fi
            %\endgroup
            %
            %\begingroup Grafici
            \pgfplotsset{
              plotRow/.style={
                scale only axis,
                width=\plotWidth cm,height=\plotHeight cm,
                xlabel style={font=\color{white!15!black}\textsize{\labelSize}},
                % ylabel style={font=\color{white!15!black}\textsize{\labelSize},yLabelShift},
                xticklabel style={font=\textsize{\tickSize}},
                % axis x line*=bottom,yLineAxis,
                enlarge x limits=0.1,enlarge y limits=0.1,
                xmajorgrids,ymajorgrids,
                xminorticks=true,yminorticks=true,
                scale ticks above exponent=2,
                axis background/.style={fill=white},
                grid style={dashed},
              },
              plotLeftCol/.style={
                yLabelInv,
                yticklabel style={font=\textsize{\tickSize},xshift=0.15em},
                yticklabel pos=left,
                yScaleLabelShiftInv,
              },
              plotMiddleCol/.style={
                yLabelEmpty,
                yTickLabelsEmpty,
                yTickScaleLabelEmpty,
              },
              plotRightCol/.style={
                yLabel,
                yticklabel style={font=\textsize{\tickSize},xshift=-0.15em},
                yticklabel pos=right,
                yScaleLabelShift,
              },
              xlabel3/.style={xlabel={\(h\)}},
              ylabel23/.style={ylabel={\(\FRC(h)\)}},
              row1Keys/.style={
                ylabel={\(\bar l_N(h)\)},
                enlarge x limits=.05,
                enlarge y limits={upper,value=.05},
                xmode=linear,
                log origin=infty,
              },
              row2Keys/.style={
                ylabel23,
                xmode=linear,ymode=log,
                % xticklabels={},
              },
              row3Keys/.style={
                xlabel3,ylabel23,
                xmode=linear,ymode=log,
              },
              row4KeysLin/.style={
                xlabel={\(\xi\)},ylabel={\(s(\xi)\)},
                xmode=log,
              },
              row4Keys/.style={row4KeysLin,ymode=log},
              row5KeysLin/.style={
                width=\plotLastWidth cm,
                xlabel={\(t\)},ylabel={\(s(t)\)},
                enlarge x limits=0,
                enlarge y limits=0.05,
                x tick scale label style={
                  yshift=-1em,
                  at={(.95,0)},
                  anchor=north,
                },
              },
              row5Keys/.style={
                row5KeysLin,
                ymode=log,
              },
            }
            %\endgroup
          %\endgroup
        }
        %\endgroup
      %\endgroup
    %\endgroup

    %\begingroup Diagrammi dei nodi
    \newcommand\nodeDiagramSettings{
      %\begingroup Dimensioni
        \pgfmathsetmacro\xPlotSep{10}
        \pgfmathsetmacro\yPlotSep{4}
        \pgfmathsetmacro\plotScale{.8}

        \pgfmathsetmacro\titleSize{11}
        \pgfmathsetmacro\commentSize{10}

        \pgfmathsetmacro\nodeSize{20pt} % Diametro del nodo
        \pgfmathsetmacro\starSize{2.25*\nodeSize}
        \pgfmathsetmacro\numberSize{.53*\nodeSize}
      %\endgroup

      %\begingroup Colori
        % \definecolor{cRed}{HTML}{D68490}
        % \definecolor{cPurple}{HTML}{7E6EAC}
        % \definecolor{cTeal}{HTML}{87BDB1}
        % \definecolor{cOrange}{HTML}{F2D0A4}
        % \definecolor{cEdgeLight}{HTML}{EEE8D5} 
        % \definecolor{cEdgeDark}{HTML}{998888}
        % \definecolor{cStar}{HTML}{FDF6E3}
        %
        \definecolor{color1}{HTML}{800000} % Deep Red
        \definecolor{color2}{HTML}{4682B4} % Steel Blue
        \definecolor{color3}{HTML}{6B8E23} % Olive Green
        \definecolor{color4}{HTML}{708090} % Slate Gray
        \definecolor{color5}{HTML}{BC8F8F} % Rosy Brown
        \definecolor{color6}{HTML}{008B8B} % Dark Cyan
        \definecolor{color7}{HTML}{B8860B} % Dark Goldenrod
      %\endgroup

      %\begingroup Testo
        \tikzset{
          titleStyle/.style={
            font=\textsize{\titleSize},
            align=center,
            anchor=south,
          },
          commentStyle/.style={font=\textsize{\commentSize}},
          arrowCommentStyle/.style={
            midway,
            commentStyle,
            left,
            sloped,
            below
          },
          numberStyle/.style={font=\textsize{\numberSize}}
        }
      %\endgroup

      %\begingroup Frecce
        \tikzset{
          transformingArrow/.style={
            % =Latex,
            ->,
            >=Stealth,
            dashed,
            line width=.8pt,
            black!90,
            shorten >=5pt,
            shorten <=10pt
          },
          indirectArrow/.style={
            draw=black!95,
            line width=.5pt,
          },
          nodeArrowExact/.style={
            indirectArrow,
            % ->,
            % -{Stealth[scale=1.2]},
            -{Triangle[scale=.75]},
            % draw=cEdgeLight,
            % draw=red!50,
            % shorten >=2pt,
            % shorten <=2pt
          },
          nodeArrowApprox/.style={
            indirectArrow,
            -,
            dashed,
          },
          transparentNodeArrow/.style={
            opacity=.4,
          },
          arrowAxis/.style={
            ->,
            line width=.3pt,
            >=Latex,
          },
        }
      %\endgroup

      %\begingroup Nodi
        \tikzset{
          nodeStyle/.style={
            circle,
            minimum size=\nodeSize,
            % draw=black!80,
            draw=black,
            line width=.5pt,
            inner sep=0pt,
          },
          starStyle/.style={
            star,
            star points=5,
            star point ratio=2,
            % fill=cStar,
            fill=black!50,
            fill opacity=.25,
            minimum size=\starSize,
            draw=none
          },
        }
      %\endgroup
    }
    %\endgroup

    %\begingroup GNU-Plot
      % \IfFileExists{\sourceauxpath.aux}{
      %   \externaldocument{\sourceauxpath}
      % }{\relax}
      
      \pgfkeys{/pgf/plot/gnuplot call={cd /d "\sourceoutdir" && gnuplot}}
    %\endgroup
  %\endgroup
\fi
%\endgroup

%\begingroup For the paper
  \def\tsc#1{\csdef{#1}{\textsc{\lowercase{#1}}\xspace}}
  \tsc{WGM}
  \tsc{QE}
  
  \newcommand\etal{et al.}
  \newcommand\KTMAS{\hyperlink{acrKineticTheoryMultiAgentSystems}{KTMAS}}
  
  % \newcommand\MathSpacing[3]{
  %   \thinmuskip = #1mu\relax
  %   \medmuskip = #2mu\relax
  %   \thickmuskip = #3mu\relax
  % }
  % \MathSpacing123
  
  \newcommand\RPlus{\mathbb R_+}
  \newcommand\NPlus{\mathbb N_+}
  \DeclareMathAlphabet\pazocal{OMS}{zplm}{m}{n}
  
  \DeclareMathOperator\Lognormal{Lognormal}
  \DeclareMathOperator\BiLognormal{BiLognormal}
  \DeclareMathOperator\Bernoulli{Bernoulli}
  \DeclareMathOperator\Exponential{Exp}
  \DeclareMathOperator\Student{Student}
  \DeclareMathOperator\Variance{var}
  \DeclareMathOperator\IC{IC}
  \DeclareMathOperator\Histogram{\mathcal B}
  \DeclareMathOperator\CCDF{\mathcal R}
    
  \newcommand\ThreeQuarters{{3\hspace{-0.15em}/\hspace{-0.15em}4}}
  \newcommand\OneQuarter{{1\hspace{-0.15em}/\hspace{-0.15em}4}}
  
  \newcommand\de[1][]{\mathrm{d}#1}
  \newcommand\derS[3][d]{\frac{\mathrm{#1}#2}{\mathrm{#1}#3}}
  \newcommand\derP[2]{\frac{\partial#1}{\partial#2}}
  \newcommand\underhat[1]{\underaccent{\hat}{#1}}
  \renewcommand\undertilde[1]{\underaccent{\tilde}{#1}}
  
  \makeatletter
  \newcommand\eqtag{\@ifstar{\eqtag@star}{\eqtag@nostar}}
  
  \newcounter{eqtag}
  \newcommand\eqtag@nostar[2]{%
    \begin{equation*}%
      \refstepcounter{eqtag}%
      % \protected@edef\@currentlabeltype{eqtag}%
      \protected@edef\@currentlabel{#1}%
      #2%
      \tag*{#1}%
    \end{equation*}%
  }
  \newcommand\eqtag@star[2]{%
    \begin{equation*}%
      \refstepcounter{eqtag}%
      % \protected@edef\@currentlabeltype{eqtag}%
      \protected@edef\@currentlabel{#1}%
      \let\label\ltx@label%
      #2%
      % \notag%
    \end{equation*}%
  }
  \makeatother
  
  \newcommand\crefsubfigInt[4][-]{Figures~\ref{#2}#3#1\ref{#2}#4}
  \newcommand\crefsubfigSep[4]{Figures~\ref{#1}#2\foreach \l in#3{,~\ref{#1}\l}~and~\ref{#1}#4}
  
  \crefname{section}{§}{§§}
  \crefrangeformat{section}{§§~#3#1#4--#5#2#6}
  \crefmultiformat{section}{§§~#2#1#3}{ and~#2#1#3}{, #2#1#3}{ and~#2#1#3}
  \crefrangemultiformat{section}{§§~#3#1#4--#5#2#6}{ and~#3#1#4--#5#2#6}{, #3#1#4--#5#2#6}{ and~#3#1#4--#5#2#6}
  
  \crefname{equation}{}{}
  \crefrangeformat{equation}{(#3#1#4--#5#2#6)}
  \crefmultiformat{equation}{(#2#1#3}{,~#2#1#3)}{,~#2#1#3}{,~#2#1#3)}
  \crefrangemultiformat{equation}{(#3#1#4--#5#2#6}{,~#3#1#4--#5#2#6)}{,~#3#1#4--#5#2#6}{,~#3#1#4--#5#2#6)}
  
  \crefname{figure}{Figure}{Figures}
  \crefrangeformat{figure}{Figures~#3#1#4--#5#2#6}
  \crefmultiformat{figure}{Figures~#2#1#3}{ and~#2#1#3}{, #2#1#3}{ and~#2#1#3}
  \crefrangemultiformat{figure}{Figures~#3#1#4--#5#2#6}{ and~#3#1#4--#5#2#6}{, #3#1#4--#5#2#6}{ and~#3#1#4--#5#2#6}
  
  \crefname{table}{Table}{Tables}
  \crefrangeformat{table}{Tables~#3#1#4--#5#2#6}
  \crefmultiformat{table}{Tables~#2#1#3}{ and~#2#1#3}{, #2#1#3}{ and~#2#1#3}
  \crefrangemultiformat{table}{Tables~#3#1#4--#5#2#6}{ and~#3#1#4--#5#2#6}{, #3#1#4--#5#2#6}{ and~#3#1#4--#5#2#6}
  
  \crefname{algorithm}{Algorithm}{Algorithms}
  \crefrangeformat{algorithm}{Algorithms~#3#1#4--#5#2#6}
  \crefmultiformat{algorithm}{Algorithms~#2#1#3}{ and~#2#1#3}{, #2#1#3}{ and~#2#1#3}
  \crefrangemultiformat{algorithm}{Algorithms~#3#1#4--#5#2#6}{ and~#3#1#4--#5#2#6}{, #3#1#4--#5#2#6}{ and~#3#1#4--#5#2#6}
  
  %\newtheorem{theorem}{Theorem}
  %\newtheorem{lemma}[theorem]{Lemma}
  \ifcasscclass
    % «cas-sc» fornisce \newdefinition
    \newdefinition{dfn}{Definition}
    \newdefinition{rmk}{Remark}
    \newdefinition{hyp}{Hypotesis}
  \else
    % Con «article» (Articolo) si usano gli ambienti del kernel, con
    % contatore condiviso e «Assumption» come titolo di «hyp»
    \newtheorem{dfn}{Definition}
    \newtheorem{rmk}[dfn]{Remark}
    \newtheorem{hyp}[dfn]{Assumption}
  \fi
  %\newproof{pf}{Proof}
  %\newproof{pot}{Proof of Theorem \ref{thm}}
  
  \crefname{dfn}{Definition}{Definitions}
  \crefrangeformat{dfn}{Definitions~#3#1#4--#5#2#6}
  \crefmultiformat{dfn}{Definitions~#2#1#3}{ and~#2#1#3}{, #2#1#3}{ and~#2#1#3}
  \crefrangemultiformat{dfn}{Definitions~#3#1#4--#5#2#6}{ and~#3#1#4--#5#2#6}{, #3#1#4--#5#2#6}{ and~#3#1#4--#5#2#6}
  \crefname{hyp}{Hypothesis}{Hypotheses}
  
  \crefrangeformat{hyp}{Hypotheses~#3#1#4--#5#2#6}
  \crefmultiformat{hyp}{Hypotheses~#2#1#3}{ and~#2#1#3}{, #2#1#3}{ and~#2#1#3}
  \crefrangemultiformat{hyp}{Hypotheses~#3#1#4--#5#2#6}{ and~#3#1#4--#5#2#6}{, #3#1#4--#5#2#6}{ and~#3#1#4--#5#2#6}
  
  \crefname{rmk}{Remark}{Remarks}
  \crefrangeformat{rmk}{Remarks~#3#1#4--#5#2#6}
  \crefmultiformat{rmk}{Remarks~#2#1#3}{ and~#2#1#3}{, #2#1#3}{ and~#2#1#3}
  \crefrangemultiformat{rmk}{Remarks~#3#1#4--#5#2#6}{ and~#3#1#4--#5#2#6}{, #3#1#4--#5#2#6}{ and~#3#1#4--#5#2#6}
  
  % \newcommand\tikzfigure[5][0em]{%
  %   \begin{figure}[%
  %     pos=#2,%
  %     width=#4,%
  %     align=\centering%
  %   ]%
  %     \centering%
  %     \vspace*{#1}%
  %     \IfSubStr{#2}{p}{%
  %       \leavevmode\stepcounter{figposcount}%
  %       \zsaveposx{figpos-\thefigposcount}%
  %       \makebox[0pt][l]{%
  %         \hspace*{-\dimexpr\zposx{figpos-\thefigposcount}sp\relax}%
  %         \hspace*{\dimexpr0.5\paperwidth\relax}%
  %         \makebox[0pt][c]{%
  %           \scalebox{0.86}{\includegraphics{figs/#5.pdf}}%
  %         }%
  %       }
  %     }{%
  %       \includegraphics{figs/#5.pdf}%
  %     }%
  %     \caption{%
  %       \hyphenpenalty=0%
  %       \exhyphenpenalty=0%
  %       #3%
  %     }%
  %     \protect\label{#5}%
  %   \end{figure}%
  % }
%\endgroup

%\begingroup Versione per la classe «cas-sc»
  % 
  % %\begingroup For image generation
  %   %\begingroup Tipografia
  %     \newcommand\textsize[1]{%
  %       \pgfmathparse{#1*1.2}%
  %       \fontsize{#1pt}{\pgfmathresult pt}\selectfont%
  %     }
  % 
  %     \DeclareMathOperator\FRC{\mathcal R}
  % 
  %     \newif\ifcasscclass
  %     \makeatletter
  %     \@ifclassloaded{cas-sc}{\casscclasstrue}{\casscclassfalse}
  %     \makeatother
  % 
  %     \providecommand\Ifthispageodd[2]{\checkoddpage\ifoddpage#1\else#2\fi}
  %   %\endgroup
  % 
  %   %\begingroup TikZ
  %     %\begingroup Didascalie
  %       \newcommand\newcaption[2]{%
  %         \captionsetup{width=#2} % aboveskip=-.2cm,
  %         \caption{#1}%
  %       }
  %       
  %       \newcommand\newsubcaption[2][ ]{%
  %         \refstepcounter{subfigure}%
  %         (\thesubfigure)#1#2%
  %       }
  %     %\endgroup
  %     
  %     %\begingroup RedefineTikZbounds
  %       % La differenza è che colla prima bisogna catturare anche le etichette perché si crea un pdf isolato mentre colla non è importante perché viene tutto compilato sul momento
  % 
  %       %\begingroup Versione valida con esternalizzazione [ma anche senza]
  %         \newcommand\redefineTikZbounds[2]{
  %           \coordinate (RealSW) at (current bounding box.south west);
  %           \coordinate (RealNE) at (current bounding box.north east);
  %           \coordinate (Center) at ($(#1.west)!0.5!(#2.east)$);
  %     
  %           \pgfresetboundingbox 
  % 
  %           \useasboundingbox 
  %             let \p1=(RealSW),\p2=(RealNE),\p3=(Center),
  %             \n1={max(\x3-\x1,\x2-\x3)} in
  %             (\x3-\n1,\y1) rectangle (\x3+\n1,\y2+.25em);
  % 
  %           % \path[use as bounding box] 
  %           %   let \p1=(RealSW),\p2=(RealNE),\p3=(Center),
  %           %   \n1={max(\x3-\x1,\x2-\x3)} in
  %           %   (\x3-\n1,\y1) rectangle (\x3+\n1,\y2+.25em);
  %           % Lo «.25em» discende dalla traslazione imposta all'etichetta della scala verticale mediante «yScaleLabelShift»
  %         } % Scatola delimitante simmetrica
  %       %\endgroup
  % 
  %       %\begingroup Versione valida senza esternalizzazione
  %         % \newcommand\redefineTikZbounds[2]{
  %         %   % Salva le coordinate della cima e del fondo della figura corrente (cosí da mantenere i titoli e le didascalie)
  %         %   \coordinate (Cima) at (current bounding box.north);
  %         %   \coordinate (Fondo) at (current bounding box.south);
  %     
  %         %   \pgfresetboundingbox % Resetta la scatola dei limiti
  %     
  %         %   \useasboundingbox % Definisce la nuova scatola come un rettangolo
  %         %     (#1.west |- Cima)  % con ascissa sinistra data da «Col1 c1r1.west», 
  %         %     rectangle      % ascissa destra data da «Col2 c1r1.east»
  %         %     (#2.east |- Fondo); % ordinate date dalla cima e fondo originali
  %         %     % L'operatore «|-» in «(A |- B)» prende come ascissa quella del punto A mentre come coordinata quella del punto B
  %         % }
  %       %\endgroup
  %     %\endgroup
  % 
  %     %\begingroup TikZBox
  %       % Parzialmente ispirato da https://gist.github.com/Caramdir/899107
  % 
  %       \makeatletter
  %       \newsavebox\tikzBufferBox
  %       \xdef\offsetX{0pt}
  %       \xdef\offsetY{0pt}
  %       
  %       % \tikzbox[#1]{#2}{#3}[#4]{#5}
  %       % #1 -> Ancora rispetto cui posizionare l'intera scatola (di base è «south west», ossia l'angolo in basso a sinistra)
  %       % #2 -> Nodo dove posizionare l'origine (per la natura di come LaTeX riceve gli argomenti è buon pratica passare questo circondato da parentesi graffe)
  %       % #3 -> Nome della scatola limitante definita dall'ambiente «scope»
  %       % #4 -> Argomenti opzionali da passare all'ambiente «scope»; di base sono vuoti
  %       % #5 -> Contenuto della scatola circondato dall'ambiente «scope»
  %       % Es. \tikzbox{{(0,0)}}{figuraTikZ}{blocco TikZ}
  % 
  %       % Macro di apertura che controlla se v'è l'argomento facoltativo dell'ancora
  %       \def\tikzbox{\pgfutil@ifnextchar[{\tikzbox@anchor}{\tikzbox@anchor[south west]}}
  %       % La funzione «\pgfutil@ifnextchar» serve per controllare se il prossimo carattere è l'apertura di una paretensi graffa «[»
  %       
  %       % Macro intermedia per cogliere gli argomenti obbligatori e controllare quelli successivi facoltativi
  %       \def\tikzbox@anchor[#1]#2#3{\pgfutil@ifnextchar[{\tikzbox@opt[#1]#2#3}{\tikzbox@opt[#1]#2#3[]}}
  % 
  %       % Macro degli argomenti obbligatori
  %       \long\def\tikzbox@opt[#1]#2#3[#4]#5{%
  %         \node[ % Misurazione della scatola in un nodo invisibile
  %           overlay,%
  %           opacity=0,%
  %           inner sep=0pt,%
  %           outer sep=0pt%
  %         ] {%
  %           \sbox{\tikzBufferBox}{%
  %             \begin{tikzpicture}%
  %               \begin{scope}[local bounding box=measuringTikZBox]%
  %                 #5%
  %               \end{scope}%
  %               \pgfpointanchor{measuringTikZBox}{#1}%
  %               \pgfgetlastxy{\tempX}{\tempY}% 
  %               \xdef\offsetX{\tempX}%
  %               \xdef\offsetY{\tempY}%
  %             \end{tikzpicture}%
  %           }%
  %         };%
  %         %
  %         % Disegno reale traslato
  %         \begin{scope}[%
  %           local bounding box=#3,%
  %           shift={#2},%
  %           xshift=-\offsetX,%
  %           yshift=-\offsetY,%
  %           #4%
  %         ]%
  %           #5%
  %         \end{scope}%
  %       }%
  %       \makeatother
  %     %\endgroup
  % 
  %     %\begingroup Chiavi-Valori
  %       %\begingroup Testo
  %         % La dimensione del testo di base è 11pt
  %         \newcommand\titleSize{11}
  %         \newcommand\subCaptionSize{10}
  %         \newcommand\ParameterSize{9}
  %         \newcommand\labelSize{10}
  %         \newcommand\tickSize{9}
  %         \newcommand\legendSize{9}
  %       %\endgroup
  %     
  %       %\begingroup Infopagine
  %         %\begingroup Chiavi condivise
  %         \newcounter{lpc} % «loopcounter»
  %         \newcommand\infopageSharedKeys{
  %           %\begingroup Testo
  %             \renewcommand\titleSize{10}
  %             \renewcommand\subCaptionSize8
  %             \renewcommand\ParameterSize7
  %             \renewcommand\labelSize9
  %             \renewcommand\tickSize6
  %             \renewcommand\legendSize7
  %           %\endgroup
  %     
  %           %\begingroup Dimensioni
  %             \pgfmathsetmacro\plotWidth7
  %             \pgfmathsetmacro\plotHeight{.618*\plotWidth}
  %             \pgfmathsetmacro\halfPlotHeight{\plotHeight/2}
  %             %
  %             \pgfmathparse{0.825*\plotWidth}
  %             \pgfmathsetmacro\plotWidth\pgfmathresult
  %             \pgfmathsetmacro\halfWidth{\plotWidth/2}
  %             \pgfmathsetmacro\plotLastWidth{1.25*\plotWidth}
  %     
  %             \pgfmathsetmacro\xPlotSep{1}
  %             \pgfmathsetmacro\yPlotSep{1.25}
  %             \pgfmathsetmacro\yGroupSep{2.75}
  %           %\endgroup
  %     
  %           %\begingroup Legende
  %             \newcommand\placeLegend[5]{%
  %               /pgfplots/plotLegend,
  %               at={([shift={(##4,##5)}]axis description cs:##1,##2)},
  %               anchor=##3,
  %             }
  %             %
  %             \pgfplotsset{%
  %               plotLegend/.style={%
  %                 font=\textsize{\legendSize},
  %                 /pgfplots/legend cell align=left,
  %                 % align=left,
  %                 % draw=white!15!black,
  %                 draw=none, % Bordo invisibile
  %                 fill=white,
  %                 fill opacity=0.5,
  %                 text opacity=1,
  %                 % text width=2cm,
  %               },
  %               plotLegendNW/.style={legend style={\placeLegend{0}{1}{north west}{0em}{0em}}},
  %               plotLegendNE/.style={legend style={\placeLegend{1}{1}{north east}{0em}{0em}}},
  %               plotLegendSW/.style={legend style={\placeLegend{0}{0}{south west}{0em}{0em}}},
  %               plotLegendSE/.style={legend style={\placeLegend{1}{0}{south east}{0em}{0em}}},
  %               %
  %               % plotLegendW/.style={legend style={plotLegend,at={(0,.5)},anchor=west}},
  %               %
  %               % plotLegendNW/.style={legend style={plotLegend,legend pos=north west}},
  %               % plotLegendNE/.style={legend style={plotLegend,legend pos=north east}},
  %               % plotLegendSW/.style={legend style={plotLegend,legend pos=south west}},
  %               % plotLegendSE/.style={legend style={plotLegend,legend pos=south east}},
  %               %
  %               % plotLegendNWcloser/.style={legend style={plotLegend,at={(0,1)},anchor=north west}},
  %               % plotLegendNEcloser/.style={legend style={plotLegend,at={(1,1)},anchor=north east}},
  %               % plotLegendSWcloser/.style={legend style={plotLegend,at={(0,0)},anchor=south west}},
  %               % plotLegendSEcloser/.style={legend style={plotLegend,at={(0,1)},anchor=south east}},
  %             }
  %             %
  %             \pgfmathsetmacro\lineLegendLength{1}
  %             \pgfmathsetmacro\areaLegendHeigth{.25}
  %             \pgfplotsset{
  %               legend image code/.code={
  %                 \draw [mark repeat=2,mark phase=2,##1]
  %                   plot coordinates {
  %                     (0em,0em)
  %                     (.5*\lineLegendLength em,0em)
  %                     (\lineLegendLength em,0em)
  %                   };
  %               },
  %               /pgfplots/area legend/.style={
  %                 legend image code/.code={
  %                   \draw [##1,fill,draw=none] 
  %                     (0em,-\areaLegendHeigth em) 
  %                     rectangle 
  %                     (\lineLegendLength em,\areaLegendHeigth em);
  %                 },
  %               },
  %             }
  %             %
  %             \newcommand\ExactHistogramLegend{Exact histogram}
  %             \newcommand\ApproxHistogramLegend{Approx. histogram}
  %             \newcommand\RealHistogramLegend{Real histogram}
  %             %
  %             \newcommand\ExactBLNFitLegend{Exact BLN fit}
  %             \newcommand\ApproxBLNFitLegend{Approx BLN fit}
  %             \newcommand\RealBLNFitLegend{Real BLN fit}
  %             %
  %             \newcommand\ExactEmpiricCCDFLegend{Exact empiric CCDF}
  %             \newcommand\ApproxEmpiricCCDFLegend{Approx. empiric CCDF}
  %             \newcommand\RealEmpiricCCDFLegend{Real empiric CCDF}
  %             %
  %             \newcommand\ExactParetoFitLegend{Exact Pareto fit}
  %             \newcommand\ApproxParetoFitLegend{Approx. Pareto fit}
  %             \newcommand\RealParetoFitLegend{Real Pareto fit}
  %             %
  %             \newcommand\ExactScatterLegend{Exact scatter}
  %             \newcommand\ApproxScatterLegend{Approx. scatter}
  %             \newcommand\RealScatterLegend{Real scatter}
  %             %
  %             \newcommand\ClasskILegend{Class k=283}
  %             \newcommand\ClasskIILegend{Class k=83}
  %             \newcommand\ClasskIIILegend{Class k=43}
  %             \newcommand\ClasskIVLegend{Class k=10}
  %           %\endgroup
  %     
  %           %\begingroup Etichette vuote
  %             \pgfplotsset{
  %               yLabelEmpty/.style={ylabel={}},
  %               yTickLabelsEmpty/.style={yticklabels={}},
  %               yTickScaleLabelEmpty/.style={ytick scale label code/.code={}},
  %               xTickScaleLabelEmpty/.style={xtick scale label code/.code={}},
  %               emptyLeged/.style={
  %                 legend style={draw=none,fill=none,nodes={draw=none}},
  %                 legend image post style={opacity=0},
  %               },
  %             }
  %           %\endgroup
  %     
  %           %\begingroup Estetica
  %             \newcommand\smallminus{\scalebox{0.5}[1]{$-$}\hspace{-.11em}}
  %             \pgfplotsset{
  %               /pgfplots/log number format basis/.code 2 args={%
  %                 \pgfmathsetmacro\exponent{##2}%
  %                 \pgfmathparse{sign(\exponent)}%
  %                 \ifnum\pgfmathresult<0%
  %                   \pgfmathparse{abs(\exponent)}%
  %                   $##1^{\smash{\smallminus\pgfmathprintnumber{\pgfmathresult}}}$%
  %                 \else%
  %                   $##1^{\smash{\pgfmathprintnumber{\exponent}}}$%
  %                 \fi%
  %               },
  %               ytick scale label code/.code={%
  %                 \pgfmathsetmacro\exponent{##1}%
  %                 \pgfmathparse{sign(\exponent)}%
  %                 \ifnum\pgfmathresult<0
  %                   \pgfmathparse{abs(\exponent)}%
  %                   $\cdot 10^{\smash{\smallminus\pgfmathprintnumber{\pgfmathresult}}}$%
  %                 \else
  %                   $\cdot 10^{\smash{\pgfmathprintnumber{\exponent}}}$%
  %                 \fi
  %               },
  %             }
  %           %\endgroup
  %     
  %           %\begingroup Colori
  %             \definecolor{blueish}{rgb}{0.00000,0.44700,0.74100}%
  %             \definecolor{reddish}{rgb}{0.83529,0.36863,0.00000}%
  %           %\endgroup
  %     
  %           %\begingroup Commenti
  %             %\begingroup Impostazioni
  %             \tikzset{
  %               title/.style={
  %                 align=center,
  %                 anchor=south,
  %                 font=\bfseries\textsize{\titleSize},
  %               },
  %               fittedParameters/.style={
  %                 align=center,
  %                 font=\textsize{\ParameterSize},
  %               },
  %               varyingParameter/.style={
  %                 align=center,
  %                 % rotate=90,
  %                 font=\textsize{\subCaptionSize},
  %               },
  %               subCaption/.style={
  %                 % text width=\plotWidth cm,
  %                 align=center,
  %                 font=\textsize{\subCaptionSize},
  %               },
  %             }
  %             \newcommand\codetocolor[1]{\ifcase##1blueish\or reddish\else black\fi}
  %             %\endgroup
  %     
  %             %\begingroup Parametri Sardegna
  %             \newcommand\betaReal{\num{1.273}}
  %             \newcommand\xiReal{\num{0.929}}
  %             \newcommand\mIReal{\num{3.227}}
  %             \newcommand\sIReal{\num{0.408}}
  %             \newcommand\mIIReal{\num{4.099}}
  %             \newcommand\sIIReal{\num{0.556}}
  %             %\endgroup
  %     
  %             %\begingroup Parametri Pareto
  %             \newcommand\codetoParetovalue[1]{
  %               \ifcase##1\betaExact\or\betaAprx\else\betaReal\fi
  %             }
  %             \newcommand\writeParetoParameters[7][center]{
  %               \node[fittedParameters,anchor=##1] at ([shift={(##5,##6)}]##4) {%
  %                 \(\begin{aligned}
  %                   \color{\codetocolor{##2}}\beta
  %                   &\color{\codetocolor{##2}}=\codetoParetovalue{##2}\\[##7]
  %                   \color{\codetocolor{##3}}\beta
  %                   &\color{\codetocolor{##3}}=\codetoParetovalue{##3}
  %                 \end{aligned}\)
  %               };
  %             }
  %             %\endgroup
  % 
  %             %\begingroup % Vecchia versione in cui si scrivevano esplicitamente tutte le 9 combinazioni possibili
  %             % \newcommand\writeParameters[6]{%
  %             %   \node[parameter] at ([shift={(##4,##5)}]##1) {%
  %             %     \ifcase##2% Caso 0
  %             %       \ifcase##3% Caso 0
  %             %         \(\begin{aligned}
  %             %           \color{blueish}\beta&\color{blueish}=\betaExact\\[##6]
  %             %           \color{blueish}\beta&\color{blueish}=\betaExact
  %             %         \end{aligned}\)
  %             %       \or% Caso 1
  %             %         \(\begin{aligned}
  %             %           \color{blueish}\beta&\color{blueish}=\betaExact\\[##6]
  %             %           \color{reddish}\beta&\color{reddish}=\betaAprx% Caso 1
  %             %         \end{aligned}\)
  %             %       \else% Caso 2
  %             %         \(\begin{aligned}
  %             %           \color{blueish}\beta&\color{blueish}=\betaExact\\[##6]
  %             %           \color{black}\beta&\color{black}=\betaReal% Caso 2
  %             %         \end{aligned}\)
  %             %       \fi
  %             %     \or% Caso 1
  %             %       \ifcase##3% Caso 0
  %             %         \(\begin{aligned}
  %             %           \color{reddish}\beta&\color{reddish}=\betaAprx\\[##6]
  %             %           \color{blueish}\beta&\color{blueish}=\betaExact
  %             %         \end{aligned}\)
  %             %       \or% Caso 1
  %             %         \(\begin{aligned}
  %             %           \color{reddish}\beta&\color{reddish}=\betaAprx\\[##6]
  %             %           \color{reddish}\beta&\color{reddish}=\betaAprx
  %             %         \end{aligned}\)
  %             %       \else% Caso 2
  %             %         \(\begin{aligned}
  %             %           \color{reddish}\beta&\color{reddish}=\betaAprx\\[##6]
  %             %           \color{black}\beta&\color{black}=\betaReal
  %             %         \end{aligned}\)
  %             %       \fi
  %             %     \else% Caso 2
  %             %       \ifcase##3% Caso 0
  %             %         \(\begin{aligned}
  %             %           \color{black}\beta&\color{black}=\betaReal\\[##6]
  %             %           \color{blueish}\beta&\color{blueish}=\betaExact
  %             %         \end{aligned}\)
  %             %       \or% Caso 1
  %             %         \(\begin{aligned}
  %             %           \color{black}\beta&\color{black}=\betaReal\\[##6]
  %             %           \color{reddish}\beta&\color{reddish}=\betaAprx
  %             %         \end{aligned}\)
  %             %       \else% Caso 2
  %             %         \(\begin{aligned}
  %             %           \color{black}\beta&\color{black}=\betaReal\\[##6]
  %             %           \color{black}\beta&\color{black}=\betaReal
  %             %         \end{aligned}\)
  %             %       \fi
  %             %     \fi
  %             %   };
  %             % }
  %             %\endgroup
  %     
  %             %\begingroup Parametri bilognormale
  %             \newcommand\codetoparameter[1]{
  %               \ifcase##1\xi  % Caso 0
  %               \or\mu_1     % Caso 1
  %               \or\sigma_1   % Caso 2
  %               \or\mu_2     % Caso 3
  %               \else\sigma_2\fi % Caso 4
  %             }
  %             \newcommand\codetoBiLognormalvalue[2]{
  %               \ifcase##1 % Caso 0
  %                 \ifcase##2\xiExact % Caso 0
  %                 \or\mIExact    % Caso 1
  %                 \or\sIExact    % Caso 2
  %                 \or\mIIExact    % Caso 3
  %                 \else\sIIExact\fi % Caso 4
  %               \or % Caso 1
  %                 \ifcase##2\xiAprx % Caso 0
  %                 \or\mIAprx     % Caso 1
  %                 \or\sIAprx     % Caso 2
  %                 \or\mIIAprx    % Caso 3
  %                 \else\sIIAprx\fi  % Caso 4
  %               \else % Caso 2
  %                 \ifcase##2\xiReal % Caso 0
  %                 \or\mIReal     % Caso 1
  %                 \or\sIReal     % Caso 2
  %                 \or\mIIReal    % Caso 3
  %                 \else\sIIReal\fi  % Caso 4
  %               \fi
  %             }
  %             \newcommand\matrixline[2]{\codetoparameter{##2}&=\codetoBiLognormalvalue{##1}{##2}}
  %             \def\lineloop##1##2{%
  %               \matrixline{##1}{\number\value{lpc}}\stepcounter{lpc}%
  %               \ifnum\number\value{lpc}<5\\[##2]\lineloop{##1}{##2}\fi%
  %             } % https://tex.stackexchange.com/questions/491358/tikz-foreach-to-work-with-the-aligned-environment
  %             \newcommand\writeBiLognormalParameters[6][center]{
  %               \node[fittedParameters,anchor=##1] at ([shift={(##4,##5)}]##3) {%
  %                 \setcounter{lpc}{0}
  %                 \(\color{\codetocolor{##2}}
  %                 \begin{aligned}
  %                   \lineloop{##2}{##6}
  %                 \end{aligned}\)
  %               };
  %             }
  %             %\endgroup
  %     
  %             %\begingroup Altre etichette
  %             \newcommand\writeSizeEvolutionDescriptions[1][5]{%
  %               \foreach \col/\clr/\txt in {%
  %                 1/blueish/Esatto,
  %                 2/reddish/Approssimato
  %               }{
  %                 \node[
  %                   fittedParameters,
  %                   anchor=south,
  %                   yshift=.1em,
  %                 ] at (Row##1 c\col r1.south) {\color{\clr}\txt};
  %               }
  %             }
  %     
  %             \newcommand\writeVaryingParameter[5]{%
  %               \foreach \r in {1,...,5}{
  %                 \node[
  %                   varyingParameter,
  %                   anchor=##3,
  %                 ] at ([shift={(##4,##5)}]Row\r\space c##1r1.##2)
  %                 {\(\varyingParameterSymbol=\varyingParameterList\r\)};
  %               }
  %             }
  %             %\endgroup
  %           %\endgroup
  %     
  %           %\begingroup Sottodidascalie
  %             \newcommand\writeSubCaption[5]{%
  %               \node[subCaption,anchor=##3]
  %               at ([shift={(##4,##5)}]##1.##2)
  %               {\newsubcaption[]{}};
  %             }
  %             \newcommand\writeSubCaptions[5][3]{%
  %               \foreach \r in {1,...,5}{
  %                 \foreach \c in {1,...,##1}{
  %                   \writeSubCaption{Row\r\space c\c r1}{##2}{##3}{##4}{##5}
  %                 }
  %               }
  %             }
  %           %\endgroup
  %         }
  %         %\endgroup
  %     
  %         %\begingroup Chiavi del primo tipo d'infopagina
  %         \newcommand\infopageKeysI{
  %           \infopageSharedKeys
  %     
  %           %\begingroup Gruppi
  %             \ifpageisodd
  %               \pgfplotsset{/pgfplots/group/yDescPos/.style={y descriptions at=edge right}}
  %             \else
  %               \pgfplotsset{/pgfplots/group/yDescPos/.style={y descriptions at=edge left}}
  %             \fi
  %             \pgfplotsset{
  %               /pgfplots/group/plotGroup/.style={
  %                 horizontal sep=\xPlotSep em,
  %                 vertical sep=\yPlotSep em,
  %                 yDescPos,
  %                 group size=3 by 1,
  %               },
  %               /pgfplots/group/plotLastGroup/.style={
  %                 plotGroup,
  %                 x descriptions at=edge bottom,
  %               },
  %             }
  %           %\endgroup
  %     
  %           %\begingroup Stili
  %             %\begingroup Selezione «recto»/«verso»
  %             \ifpageisodd
  %               \pgfplotsset{%
  %                 yLineAxis/.style={axis y line*=right},
  %                 /tikz/yLabelShift/.style={%
  %                   yshift=-1.25em,
  %                   at={(axis description cs:1,0.5)},
  %                   anchor=north,
  %                   % rotate=180,
  %                 },
  %                 /tikz/yScaleLabelShift/.style={%
  %                   yshift=.25em,
  %                   at={(1,1)},
  %                   anchor=south east,
  %                 },
  %                 /tikz/yTickLabelShift/.style={xshift=-0.15em},
  %               }
  %             \else
  %               \pgfplotsset{%
  %                 yLineAxis/.style={axis y line*=left},
  %                 /tikz/yLabelShift/.style={%
  %                   yshift=1.25em,
  %                   at={(axis description cs:0,0.5)},
  %                   anchor=south,
  %                   % rotate=180,
  %                 },
  %                 /tikz/yScaleLabelShift/.style={%
  %                   yshift=.25em,
  %                   at={(0,1)},
  %                   anchor=south west,
  %                 },
  %                 /tikz/yTickLabelShift/.style={xshift=+0.15em},
  %               }
  %             \fi
  %             %\endgroup
  %             %
  %             %\begingroup Grafici
  %             \pgfplotsset{
  %               plotRow/.style={
  %                 scale only axis,
  %                 width=\plotWidth cm,height=\plotHeight cm,
  %                 xlabel style={font=\color{white!15!black}\textsize{\labelSize}},
  %                 ylabel style={font=\color{white!15!black}\textsize{\labelSize},yLabelShift},
  %                 xticklabel style={font=\textsize{\tickSize}},
  %                 yticklabel style={font=\textsize{\tickSize},yTickLabelShift},
  %                 y tick scale label style={yScaleLabelShift},
  %                 % axis x line*=bottom,yLineAxis,
  %                 enlarge x limits=0.1,enlarge y limits=0.1,
  %                 xmajorgrids,ymajorgrids,
  %                 xminorticks=true,yminorticks=true,
  %                 scale ticks above exponent=2,
  %                 axis background/.style={fill=white},
  %                 grid style={dashed}
  %               },
  %               % xlabel123/.style={xlabel={\(h\)}},
  %               xlabel3/.style={xlabel={\(h\)}},
  %               ylabel23/.style={ylabel={\(\FRC(h)\)}},
  %               row1Keys/.style={
  %                 % ylabel={\(\underhat{\bar l}_N(h)\)},
  %                 ylabel={\(\bar l_N(h)\)},
  %                 % enlarge x limits={lower,value=.05},
  %                 enlarge x limits=.05,
  %                 enlarge y limits={upper,value=.05},
  %                 xmode=linear,
  %                 log origin=infty,
  %               },
  %               row2KeysWithLabels/.style={
  %                 ylabel23,
  %                 xmode=linear,ymode=log
  %               },
  %               row2Keys/.style={
  %                 ylabel23,
  %                 xmode=linear,ymode=log,
  %                 xticklabels={},
  %               },
  %               row3Keys/.style={
  %                 xlabel3,ylabel23,
  %                 xmode=linear,ymode=log
  %               },
  %               row4Keys/.style={
  %                 xlabel={\(\xi\)},ylabel={\(s(\xi)\)},
  %                 xmode=log,ymode=log
  %               },
  %               row5KeysLin/.style={
  %                 width=\plotLastWidth cm,
  %                 xlabel={\(t\)},ylabel={\(s(t)\)},
  %                 enlarge x limits=0,
  %                 enlarge y limits=0.05,
  %                 x tick scale label style={
  %                   yshift=-1em,
  %                   at={(.95,0)},
  %                   anchor=north,
  %                 },
  %               },
  %               row5Keys/.style={
  %                 row5KeysLin,
  %                 ymode=log,
  %               },
  %             }
  %             %\endgroup
  %           %\endgroup
  %         }
  %         %\endgroup
  %     
  %         %\begingroup Chiavi del secondo tipo d'infopagina
  %         \newcommand\infopageKeysII{
  %           \infopageSharedKeys
  %     
  %           %\begingroup Gruppi
  %             \ifpageisodd
  %               \pgfplotsset{/pgfplots/group/yDescPos/.style={y descriptions at=edge right}}
  %             \else
  %               \pgfplotsset{/pgfplots/group/yDescPos/.style={y descriptions at=edge left}}
  %             \fi
  %             \pgfplotsset{
  %               /pgfplots/group/plotGroup/.style={
  %                 horizontal sep=\xPlotSep em,
  %                 vertical sep=\yPlotSep em,
  %                 y descriptions at=all,
  %                 group size=3 by 1,
  %               },
  %               /pgfplots/group/plotLastGroup/.style={
  %                 plotGroup,
  %                 x descriptions at=edge bottom,
  %               },
  %             }
  %           %\endgroup
  %     
  %           %\begingroup Stili
  %             %\begingroup Selezione «recto»/«verso»
  %             \ifpageisodd
  %               \pgfplotsset{
  %                 yLabel/.style={
  %                   ylabel style={
  %                     font=\color{white!15!black}\textsize{\labelSize},
  %                     yshift=-1.25em,
  %                     at={(axis description cs:1,0.5)},
  %                     anchor=north,
  %                     % rotate=180,
  %                   },
  %                 },
  %                 yLabelInv/.style={yLabelEmpty},
  %                 yScaleLabelShift/.style={
  %                   y tick scale label style={
  %                     yshift=.25em,
  %                     % xshift=1em,
  %                     at={(1,1)},
  %                     anchor=south east,
  %                   },
  %                 },
  %                 yScaleLabelShiftInv/.style={
  %                   y tick scale label style={
  %                     yshift=.25em,
  %                     % xshift=-1em,
  %                     at={(0,1)},
  %                     anchor=south west,
  %                   },
  %                 },
  %                 % yLineAxis/.style={axis y line*=right},
  %               }
  %             \else
  %               \pgfplotsset{
  %                 yLabel/.style={yLabelEmpty},
  %                 yLabelInv/.style={
  %                   ylabel style={
  %                     font=\color{white!15!black}\textsize{\labelSize},
  %                     yshift=1.25em,
  %                     at={(axis description cs:0,0.5)},
  %                     anchor=south,
  %                     % rotate=180,
  %                   },
  %                 },
  %                 yScaleLabelShift/.style={
  %                   y tick scale label style={
  %                     yshift=.25em,
  %                     % xshift=1em,
  %                     at={(1,1)},
  %                     anchor=south east,
  %                   },
  %                 },
  %                 yScaleLabelShiftInv/.style={
  %                   y tick scale label style={
  %                     yshift=.25em,
  %                     % xshift=-1em,
  %                     at={(0,1)},
  %                     anchor=south west,
  %                   },
  %                 },
  %                 % yLineAxis/.style={axis y line*=left},
  %               }
  %             \fi
  %             %\endgroup
  %             %
  %             %\begingroup Grafici
  %             \pgfplotsset{
  %               plotRow/.style={
  %                 scale only axis,
  %                 width=\plotWidth cm,height=\plotHeight cm,
  %                 xlabel style={font=\color{white!15!black}\textsize{\labelSize}},
  %                 % ylabel style={font=\color{white!15!black}\textsize{\labelSize},yLabelShift},
  %                 xticklabel style={font=\textsize{\tickSize}},
  %                 % axis x line*=bottom,yLineAxis,
  %                 enlarge x limits=0.1,enlarge y limits=0.1,
  %                 xmajorgrids,ymajorgrids,
  %                 xminorticks=true,yminorticks=true,
  %                 scale ticks above exponent=2,
  %                 axis background/.style={fill=white},
  %                 grid style={dashed},
  %               },
  %               plotLeftCol/.style={
  %                 yLabelInv,
  %                 yticklabel style={font=\textsize{\tickSize},xshift=0.15em},
  %                 yticklabel pos=left,
  %                 yScaleLabelShiftInv,
  %               },
  %               plotMiddleCol/.style={
  %                 yLabelEmpty,
  %                 yTickLabelsEmpty,
  %                 yTickScaleLabelEmpty,
  %               },
  %               plotRightCol/.style={
  %                 yLabel,
  %                 yticklabel style={font=\textsize{\tickSize},xshift=-0.15em},
  %                 yticklabel pos=right,
  %                 yScaleLabelShift,
  %               },
  %               xlabel3/.style={xlabel={\(h\)}},
  %               ylabel23/.style={ylabel={\(\FRC(h)\)}},
  %               row1Keys/.style={
  %                 ylabel={\(\bar l_N(h)\)},
  %                 enlarge x limits=.05,
  %                 enlarge y limits={upper,value=.05},
  %                 xmode=linear,
  %                 log origin=infty,
  %               },
  %               row2Keys/.style={
  %                 ylabel23,
  %                 xmode=linear,ymode=log,
  %                 % xticklabels={},
  %               },
  %               row3Keys/.style={
  %                 xlabel3,ylabel23,
  %                 xmode=linear,ymode=log,
  %               },
  %               row4KeysLin/.style={
  %                 xlabel={\(\xi\)},ylabel={\(s(\xi)\)},
  %                 xmode=log,
  %               },
  %               row4Keys/.style={row4KeysLin,ymode=log},
  %               row5KeysLin/.style={
  %                 width=\plotLastWidth cm,
  %                 xlabel={\(t\)},ylabel={\(s(t)\)},
  %                 enlarge x limits=0,
  %                 enlarge y limits=0.05,
  %                 x tick scale label style={
  %                   yshift=-1em,
  %                   at={(.95,0)},
  %                   anchor=north,
  %                 },
  %               },
  %               row5Keys/.style={
  %                 row5KeysLin,
  %                 ymode=log,
  %               },
  %             }
  %             %\endgroup
  %           %\endgroup
  %         }
  %         %\endgroup
  %       %\endgroup
  %     %\endgroup
  % 
  %     %\begingroup Diagrammi dei nodi
  %     \newcommand\nodeDiagramSettings{
  %       %\begingroup Dimensioni
  %         \pgfmathsetmacro\xPlotSep{10}
  %         \pgfmathsetmacro\yPlotSep{4}
  %         \pgfmathsetmacro\plotScale{.8}
  % 
  %         \pgfmathsetmacro\titleSize{11}
  %         \pgfmathsetmacro\commentSize{10}
  % 
  %         \pgfmathsetmacro\nodeSize{20pt} % Diametro del nodo
  %         \pgfmathsetmacro\starSize{2.25*\nodeSize}
  %         \pgfmathsetmacro\numberSize{.53*\nodeSize}
  %       %\endgroup
  % 
  %       %\begingroup Colori
  %         % \definecolor{cRed}{HTML}{D68490}
  %         % \definecolor{cPurple}{HTML}{7E6EAC}
  %         % \definecolor{cTeal}{HTML}{87BDB1}
  %         % \definecolor{cOrange}{HTML}{F2D0A4}
  %         % \definecolor{cEdgeLight}{HTML}{EEE8D5} 
  %         % \definecolor{cEdgeDark}{HTML}{998888}
  %         % \definecolor{cStar}{HTML}{FDF6E3}
  %         %
  %         \definecolor{color1}{HTML}{800000} % Deep Red
  %         \definecolor{color2}{HTML}{4682B4} % Steel Blue
  %         \definecolor{color3}{HTML}{6B8E23} % Olive Green
  %         \definecolor{color4}{HTML}{708090} % Slate Gray
  %         \definecolor{color5}{HTML}{BC8F8F} % Rosy Brown
  %         \definecolor{color6}{HTML}{008B8B} % Dark Cyan
  %         \definecolor{color7}{HTML}{B8860B} % Dark Goldenrod
  %       %\endgroup
  % 
  %       %\begingroup Testo
  %         \tikzset{
  %           titleStyle/.style={
  %             font=\textsize{\titleSize},
  %             align=center,
  %             anchor=south,
  %           },
  %           commentStyle/.style={font=\textsize{\commentSize}},
  %           arrowCommentStyle/.style={
  %             midway,
  %             commentStyle,
  %             left,
  %             sloped,
  %             below
  %           },
  %           numberStyle/.style={font=\textsize{\numberSize}}
  %         }
  %       %\endgroup
  % 
  %       %\begingroup Frecce
  %         \tikzset{
  %           transformingArrow/.style={
  %             % =Latex,
  %             ->,
  %             >=Stealth,
  %             dashed,
  %             line width=.8pt,
  %             black!90,
  %             shorten >=5pt,
  %             shorten <=10pt
  %           },
  %           indirectArrow/.style={
  %             draw=black!95,
  %             line width=.5pt,
  %           },
  %           nodeArrowExact/.style={
  %             indirectArrow,
  %             % ->,
  %             % -{Stealth[scale=1.2]},
  %             -{Triangle[scale=.75]},
  %             % draw=cEdgeLight,
  %             % draw=red!50,
  %             % shorten >=2pt,
  %             % shorten <=2pt
  %           },
  %           nodeArrowApprox/.style={
  %             indirectArrow,
  %             -,
  %             dashed,
  %           },
  %           transparentNodeArrow/.style={
  %             opacity=.4,
  %           },
  %           arrowAxis/.style={
  %             ->,
  %             line width=.3pt,
  %             >=Latex,
  %           },
  %         }
  %       %\endgroup
  % 
  %       %\begingroup Nodi
  %         \tikzset{
  %           nodeStyle/.style={
  %             circle,
  %             minimum size=\nodeSize,
  %             % draw=black!80,
  %             draw=black,
  %             line width=.5pt,
  %             inner sep=0pt,
  %           },
  %           starStyle/.style={
  %             star,
  %             star points=5,
  %             star point ratio=2,
  %             % fill=cStar,
  %             fill=black!50,
  %             fill opacity=.25,
  %             minimum size=\starSize,
  %             draw=none
  %           },
  %         }
  %       %\endgroup
  %     }
  %     %\endgroup
  % 
  %     %\begingroup GNU-Plot
  %       % \IfFileExists{\sourceauxpath.aux}{
  %       %   \externaldocument{\sourceauxpath}
  %       % }{\relax}
  %       
  %       \pgfkeys{/pgf/plot/gnuplot call={cd /d "\sourceoutdir" && gnuplot}}
  %     %\endgroup
  %   %\endgroup
  % %\endgroup
  % 
  % %\begingroup For the paper
  %   \def\tsc#1{\csdef{#1}{\textsc{\lowercase{#1}}\xspace}}
  %   \tsc{WGM}
  %   \tsc{QE}
  %   
  %   \newcommand\etal{et al.}
  %   \newcommand\KTMAS{\hyperlink{acrKineticTheoryMultiAgentSystems}{KTMAS}}
  %   
  %   % \newcommand\MathSpacing[3]{
  %   %   \thinmuskip = #1mu\relax
  %   %   \medmuskip = #2mu\relax
  %   %   \thickmuskip = #3mu\relax
  %   % }
  %   % \MathSpacing123
  %   
  %   \newcommand\RPlus{\mathbb R_+}
  %   \newcommand\NPlus{\mathbb N_+}
  %   \DeclareMathAlphabet\pazocal{OMS}{zplm}{m}{n}
  %   
  %   \DeclareMathOperator\Lognormal{Lognormal}
  %   \DeclareMathOperator\BiLognormal{BiLognormal}
  %   \DeclareMathOperator\Bernoulli{Bernoulli}
  %   \DeclareMathOperator\Exponential{Exp}
  %   \DeclareMathOperator\Student{Student}
  %   \DeclareMathOperator\Variance{var}
  %   \DeclareMathOperator\IC{IC}
  %   \DeclareMathOperator\Histogram{\mathcal B}
  %   \DeclareMathOperator\CCDF{\mathcal R}
  %     
  %   \newcommand\ThreeQuarters{{3\hspace{-0.15em}/\hspace{-0.15em}4}}
  %   \newcommand\OneQuarter{{1\hspace{-0.15em}/\hspace{-0.15em}4}}
  %   
  %   \newcommand\de[1][]{\mathrm{d}#1}
  %   \newcommand\derS[3][d]{\frac{\mathrm{#1}#2}{\mathrm{#1}#3}}
  %   \newcommand\derP[2]{\frac{\partial#1}{\partial#2}}
  %   \newcommand\underhat[1]{\underaccent{\hat}{#1}}
  %   \renewcommand\undertilde[1]{\underaccent{\tilde}{#1}}
  %   
  %   \makeatletter
  %   \newcommand\eqtag{\@ifstar{\eqtag@star}{\eqtag@nostar}}
  %   
  %   \newcounter{eqtag}
  %   \newcommand\eqtag@nostar[2]{%
  %     \begin{equation*}%
  %       \refstepcounter{eqtag}%
  %       % \protected@edef\@currentlabeltype{eqtag}%
  %       \protected@edef\@currentlabel{#1}%
  %       #2%
  %       \tag*{#1}%
  %     \end{equation*}%
  %   }
  %   \newcommand\eqtag@star[2]{%
  %     \begin{equation*}%
  %       \refstepcounter{eqtag}%
  %       % \protected@edef\@currentlabeltype{eqtag}%
  %       \protected@edef\@currentlabel{#1}%
  %       \let\label\ltx@label%
  %       #2%
  %       % \notag%
  %     \end{equation*}%
  %   }
  %   \makeatother
  %   
  %   \newcommand\crefsubfigInt[4][-]{Figures~\ref{#2}#3#1\ref{#2}#4}
  %   \newcommand\crefsubfigSep[4]{Figures~\ref{#1}#2\foreach \l in#3{,~\ref{#1}\l}~and~\ref{#1}#4}
  %   
  %   \crefname{section}{§}{§§}
  %   \crefrangeformat{section}{§§~#3#1#4--#5#2#6}
  %   \crefmultiformat{section}{§§~#2#1#3}{ and~#2#1#3}{, #2#1#3}{ and~#2#1#3}
  %   \crefrangemultiformat{section}{§§~#3#1#4--#5#2#6}{ and~#3#1#4--#5#2#6}{, #3#1#4--#5#2#6}{ and~#3#1#4--#5#2#6}
  %   
  %   \crefname{equation}{}{}
  %   \crefrangeformat{equation}{(#3#1#4--#5#2#6)}
  %   \crefmultiformat{equation}{(#2#1#3}{,~#2#1#3)}{,~#2#1#3}{,~#2#1#3)}
  %   \crefrangemultiformat{equation}{(#3#1#4--#5#2#6}{,~#3#1#4--#5#2#6)}{,~#3#1#4--#5#2#6}{,~#3#1#4--#5#2#6)}
  %   
  %   \crefname{figure}{Figure}{Figures}
  %   \crefrangeformat{figure}{Figures~#3#1#4--#5#2#6}
  %   \crefmultiformat{figure}{Figures~#2#1#3}{ and~#2#1#3}{, #2#1#3}{ and~#2#1#3}
  %   \crefrangemultiformat{figure}{Figures~#3#1#4--#5#2#6}{ and~#3#1#4--#5#2#6}{, #3#1#4--#5#2#6}{ and~#3#1#4--#5#2#6}
  %   
  %   \crefname{table}{Table}{Tables}
  %   \crefrangeformat{table}{Tables~#3#1#4--#5#2#6}
  %   \crefmultiformat{table}{Tables~#2#1#3}{ and~#2#1#3}{, #2#1#3}{ and~#2#1#3}
  %   \crefrangemultiformat{table}{Tables~#3#1#4--#5#2#6}{ and~#3#1#4--#5#2#6}{, #3#1#4--#5#2#6}{ and~#3#1#4--#5#2#6}
  %   
  %   \crefname{algorithm}{Algorithm}{Algorithms}
  %   \crefrangeformat{algorithm}{Algorithms~#3#1#4--#5#2#6}
  %   \crefmultiformat{algorithm}{Algorithms~#2#1#3}{ and~#2#1#3}{, #2#1#3}{ and~#2#1#3}
  %   \crefrangemultiformat{algorithm}{Algorithms~#3#1#4--#5#2#6}{ and~#3#1#4--#5#2#6}{, #3#1#4--#5#2#6}{ and~#3#1#4--#5#2#6}
  %   
  %   %\newtheorem{theorem}{Theorem}
  %   %\newtheorem{lemma}[theorem]{Lemma}
  %   \newdefinition{dfn}{Definition}
  %   \newdefinition{rmk}{Remark}
  %   \newdefinition{hyp}{Hypotesis}
  %   %\newproof{pf}{Proof}
  %   %\newproof{pot}{Proof of Theorem \ref{thm}}
  %   
  %   \crefname{dfn}{Definition}{Definitions}
  %   \crefrangeformat{dfn}{Definitions~#3#1#4--#5#2#6}
  %   \crefmultiformat{dfn}{Definitions~#2#1#3}{ and~#2#1#3}{, #2#1#3}{ and~#2#1#3}
  %   \crefrangemultiformat{dfn}{Definitions~#3#1#4--#5#2#6}{ and~#3#1#4--#5#2#6}{, #3#1#4--#5#2#6}{ and~#3#1#4--#5#2#6}
  %   \crefname{hyp}{Hypothesis}{Hypotheses}
  %   
  %   \crefrangeformat{hyp}{Hypotheses~#3#1#4--#5#2#6}
  %   \crefmultiformat{hyp}{Hypotheses~#2#1#3}{ and~#2#1#3}{, #2#1#3}{ and~#2#1#3}
  %   \crefrangemultiformat{hyp}{Hypotheses~#3#1#4--#5#2#6}{ and~#3#1#4--#5#2#6}{, #3#1#4--#5#2#6}{ and~#3#1#4--#5#2#6}
  %   
  %   \crefname{rmk}{Remark}{Remarks}
  %   \crefrangeformat{rmk}{Remarks~#3#1#4--#5#2#6}
  %   \crefmultiformat{rmk}{Remarks~#2#1#3}{ and~#2#1#3}{, #2#1#3}{ and~#2#1#3}
  %   \crefrangemultiformat{rmk}{Remarks~#3#1#4--#5#2#6}{ and~#3#1#4--#5#2#6}{, #3#1#4--#5#2#6}{ and~#3#1#4--#5#2#6}
  %   
  %   % \newcommand\tikzfigure[5][0em]{%
  %   %   \begin{figure}[%
  %   %     pos=#2,%
  %   %     width=#4,%
  %   %     align=\centering%
  %   %   ]%
  %   %     \centering%
  %   %     \vspace*{#1}%
  %   %     \IfSubStr{#2}{p}{%
  %   %       \leavevmode\stepcounter{figposcount}%
  %   %       \zsaveposx{figpos-\thefigposcount}%
  %   %       \makebox[0pt][l]{%
  %   %         \hspace*{-\dimexpr\zposx{figpos-\thefigposcount}sp\relax}%
  %   %         \hspace*{\dimexpr0.5\paperwidth\relax}%
  %   %         \makebox[0pt][c]{%
  %   %           \scalebox{0.86}{\includegraphics{figs/#5.pdf}}%
  %   %         }%
  %   %       }
  %   %     }{%
  %   %       \includegraphics{figs/#5.pdf}%
  %   %     }%
  %   %     \caption{%
  %   %       \hyphenpenalty=0%
  %   %       \exhyphenpenalty=0%
  %   %       #3%
  %   %     }%
  %   %     \protect\label{#5}%
  %   %   \end{figure}%
  %   % }
  % %\endgroup
%\endgroup

  % 
% !TEX root = ../../Esperimenti/.tex/MEF.tex

% TikZFigureExternalized
\newcommand\compilerpath{../Altro/tikz2pdf.tex}
\newcommand\pdfdirectory{../../Figure/.fnl/Articolo/TikZpdf/}
\newcommand\auxpath{../Altro/mainCopy}
\newcommand\outdirectory\pdfdirectory

\makeatletter

% \tikzexternalizedfigure[#1]{#2}[#3]{#4}{#5}{#6}{#7}
% #1 -> Tipo di compilatore (di base è «pdflatex»)
% #2 -> Posizione della figura flottante
% #3 -> Tralazione verticale (di base è nulla)
% #4 -> Percorso assoluto o relativo della filza .tex da compilare esternamente
% #5 -> Didascalia
% #6 -> Lunghezza della didascalia
% #7 -> Nome della filza «.pdf» della figura tikz compilata

% https://tex.stackexchange.com/questions/57788/understanding-ifnextchar
\def\tikzexternalizedfigure{%
  \@ifnextchar[{\tikz@ext@figure@compiler}{\tikz@ext@figure@compiler[pdflatex]}
}
\def\tikz@ext@figure@compiler[#1]#2{%
  \@ifnextchar[{\tikz@ext@figure@translation[#1]#2}{\tikz@ext@figure@translation[#1]#2[0em]}
}
\def\tikz@ext@figure@translation[#1]#2[#3]#4#5#6#7{%
  % \newcommand\outdirectory{\pdfdirectory#6.}
  \begin{figure}[%
   pos=#2,%
   width=#6,%
   align=\centering%
  ]%
    \centering%
    % Recupera la pagina dall'etichetta (default a 1 se non esiste)
    \checkoddpage%
    \Ifthispageodd{\def\pagetype{odd}}{\def\pagetype{even}}%
    %
    % Alternativa con «\getpagerefnumber»
    % \edef\currentfigpage{\getpagerefnumber{#7}}
    % \ifx\currentfigpage\relax\def\currentfigpage{1} \fi
    % \pgfmathparse{int(mod(\currentfigpage,2) == 1 ? 1 : 0)}
    % \ifnum\pgfmathresult=1\def\pagetype{odd}
    % \else\def\pagetype{even}\fi
    %
    \def\shlcmd{%
      \immediate\ShellEscape{%
        #1 %
        --shell-escape %
        -interaction=batchmode %
        -output-directory="\pdfdirectory" %
        -jobname="#7" %
        "\string\def\string\isodd{\pagetype}%
        \string\def\string\sourcefilepath{#4}%
        \string\def\string\sourceoutdir{\outdirectory}%
        \string\def\string\sourceauxpath{\auxpath}%
        \string\input{\compilerpath}"%
      }%
    }%
    \IfSubStr{#2}{p}{% 
      \leavevmode\stepcounter{figposcount}\zsaveposx{figpos-\thefigposcount}%
      \makebox[0pt][l]{%
        \hspace*{-\dimexpr\zposx{figpos-\thefigposcount}sp\relax}%
        % \hspace*{\dimexpr0.5\paperwidth\relax}%
        \def\hsep{0cm}\Ifthispageodd{\def\offset{+\hsep}}{\def\offset{-\hsep}}% per far toccare il bordo sinistro/destro dei grafici al margine interno/esterno si trasla di -/+1.9 cm
        \hspace*{\dimexpr0.5\paperwidth+\offset\relax}%
        \makebox[0pt][c]{%
          \begin{minipage}\paperwidth%
            \centering%
            \vspace*{#3}%
            \addtocounter{figure}{1}%
            \setcounter{subfigure}{0}%
            \IfFileExists{\pdfdirectory#7.pdf}{%
              \ifbool{forcetikzrecompile}{\shlcmd}{\relax}%
            }{\shlcmd}%
            \scalebox{0.86}{%
              \includegraphics{\pdfdirectory#7.pdf}%
            }%
            \addtocounter{figure}{-1}%
            \newcaption{%
              \hyphenpenalty=0%
              \exhyphenpenalty=0%
              #5%
            }{#6}%
            \protect\label{#7}%
          \end{minipage}%
        }%
      }%
    }{%
      \makebox[\textwidth][c]{%
        \begin{minipage}\linewidth
          \centering%
          \addtocounter{figure}{1}\setcounter{subfigure}{0}%
          \IfFileExists{\pdfdirectory#7.pdf}{%
            \ifbool{forcetikzrecompile}{\shlcmd}{\relax}
          }{\shlcmd}%
          \includegraphics{\pdfdirectory#7.pdf}%
          \addtocounter{figure}{-1}%
          \newcaption{%
            \hyphenpenalty=0%
            \exhyphenpenalty=0%
            #5%
          }{#6}%
          \protect\label{#7}%
        \end{minipage}%
      }%
    }%
  \end{figure}%
}
\makeatother

%\begingroup Vecchia versione con un solo argomento facoltativo
  % \newcommand\tikzexternalizedfigure[6][0em]{%
  %   % \newcommand\outdirectory{\pdfdirectory#6.}
  %   \begin{figure}[#2]%
  %     \centering%
  %     % Recupera la pagina dall'etichetta (default a 1 se non esiste)
  %     \checkoddpage
  %     \Ifthispageodd{\def\pagetype{odd}}{\def\pagetype{even}}
  %     %
  %     % Alternativa con «\getpagerefnumber»
  %     % \edef\currentfigpage{\getpagerefnumber{#6}}
  %     % \ifx\currentfigpage\relax\def\currentfigpage{1} \fi
  %     % \pgfmathparse{int(mod(\currentfigpage,2) == 1 ? 1 : 0)}
  %     % \ifnum\pgfmathresult=1\def\pagetype{odd}
  %     % \else\def\pagetype{even}\fi
  %     %
  %     \edef\shlcmd{
  %       \immediate\ShellEscape{%
  %         lualatex %
  %         --shell-escape %
  %         -interaction=batchmode %
  %         -output-directory="\pdfdirectory" %
  %         -jobname="#6" %
  %         "\string\def\string\isodd{\pagetype}%
  %         \string\def\string\sourcefilepath{#3}%
  %         \string\def\string\sourceoutdir{\outdirectory}%
  %         \string\def\string\sourceauxpath{\auxpath}%
  %         \string\input{\compilerpath}"%
  %       }
  %     }
  %     \IfSubStr{#2}{p}{% 
  %       \leavevmode\stepcounter{figposcount}\zsaveposx{figpos-\thefigposcount}%
  %       \makebox[0pt][l]{%
  %         \hspace*{-\dimexpr\zposx{figpos-\thefigposcount}sp\relax}%
  %         \hspace*{0.5\paperwidth}%
  %         \makebox[0pt][c]{%
  %           \begin{minipage}{\paperwidth}%
  %             \vspace*{#1}\centering
  %             \addtocounter{figure}{1}
  %             \setcounter{subfigure}{0}%
  %             \IfFileExists{\pdfdirectory#6.pdf}{%
  %               \ifbool{forcetikzrecompile}{\shlcmd}{\relax}
  %             }{\shlcmd}%
  %             \includegraphics{\pdfdirectory#6.pdf}%
  %             \addtocounter{figure}{-1}
  %             \didascalia{#4}{#5}
  %             \protect\label{#6}
  %           \end{minipage}
  %         }%
  %       }%
  %     }{%
  %       \makebox[\textwidth][c]{
  %         \begin{minipage}{\linewidth}
  %           \centering
  %           \addtocounter{figure}{1}\setcounter{subfigure}{0}
  %           \IfFileExists{\pdfdirectory#6.pdf}{%
  %             \ifbool{forcetikzrecompile}{\shlcmd}{\relax}
  %           }{\shlcmd}%
  %           \includegraphics{\pdfdirectory#6.pdf}%
  %           \addtocounter{figure}{-1}
  %           \didascalia{#4}{#5}
  %           \protect\label{#6}
  %         \end{minipage}%
  %       }%
  %     }%
  %   \end{figure}%
  % }
%\endgroup


  % \setbool{forcetikzrecompile}{true}

  % \IfFileExists{main.aux}{%
  %   \immediate\ShellEscape{copy /Y "main.aux" "../Altro/mainCopy.aux"}
  % }{\relax}

  % \begin{document}
  %   \let\WriteBookmarks\relax
  %   \def\floatpagepagefraction{1}
  %   \def\textpagefraction{.001}

    %\begingroup Articolo
      % Main title of the paper
  %     \title[mode=title]{Modelling population distribution among cities on spatial networks using the kinetic theory of multi-agent systems}%: the Sardinia case}

  %     \shorttitle{Modelling population distribution on spatial networks}
  %     \shortauthors{V. Taralli, A. Tosin}

      % \tnotemark[1] % Title footnote mark
      % \tnotetext[1]{This paper is a rework of Taralli's Master's Degree in Mathematical Engineering, Academic Year 2025-2026} % Title footnote

      % Authors
      % Options: Use if required
      % eg: \author[1,3]{Author Name}[type=editor,
      %    style=chinese,
      %    auid=000,
      %    bioid=1,
      %    prefix=Sir,
      %    orcid=0000-0000-0000-0000,
      %    facebook=<facebook id>,
      %    twitter=<twitter id>,
      %    linkedin=<linkedin id>,
      %    gplus=<gplus id>]

  %     \author{Valerio Taralli}[%
      % \author[1]{Valerio Taralli}[%
      % linkedin=https://www.linkedin.com/in/valerio-taralli/,%
  %     ]
      % \cormark[1]% Corresponding author indication
      % \fnmark[1]% Footnote of the first author
  %     \ead{valerio.taralli2001@gmail.com}% Email id of the first author
      % \ead[url]{}% URL of the first author
  %     \credit{Conceptualization, Methodology, Software, Formal analysis, Investigation, Writing - Original Draft}% Credit authorship
      % eg: \credit{Conceptualization of this study, Methodology, Software}

      % % Address/affiliation
      % \affiliation[1]{organization={},
      %       addressline={}, 
      %       city={},
      % %     citysep={}, % Uncomment if no comma needed between city and postcode
      %       postcode={}, 
      %       state={},
      %       country={}}

  %     \author[1]{Andrea Tosin}[%
  %       orcid=0000-0003-4955-1818,%
  %     ]
  %     \cormark[1]% Footnote of the second author
  %     \ead{andrea.tosin@polito.it}% Email id of the second author
      % \ead[url]{https://orcid.org/0000-0003-4955-1818}% URL of the second author
  %     \credit{Supervision}% Credit authorship

      % Address/affiliation
  %     \affiliation[1]{%
  %       organization={Department of Mathematical Sciences “G. L. Lagrange”},%
  %       addressline={Politecnico di Torino},%
  %       city={Turin},%
        % citysep={}, % Uncomment if no comma needed between city and postcode%
        % postcode={10129},%
  %       state={Piedomont},%
  %       country={Italy}%
  %     }

      % Corresponding author text
  %     \cortext[1]{Corresponding author: A. Tosin; Email: andrea.tosin@polito.it}
      % \fntext[1]{A. Tosin; Email: \href{mailto:andrea.tosin@polito.it}{andrea.tosin@polito.it}}% Footnote text

      % For a title note without a number/mark
      %\nonumnote{}

  %     \begin{abstract}
  %       This work investigates the population distribution through the kinetic theory of multi-agent systems on spatial networks, interpreting cities both as agents and vertices, and representing interurban links as edges. Various emigration rules have been proposed to regulate the interactions between agents, yielding two primary results: the bimodal lognormal fit of the total population distribution among cities and the Pareto fit of the tail. Overall, the model provides results consistent with the actual distribution derived from ISTAT data of Sardinia; furthermore, the bimodal lognormal fit generally follows the distribution tail more accurately than the Pareto one, although the two coincide in the first half.
  %       \nocite{*}
  %     \end{abstract}

      % Research highlights
  %     \begin{highlights}
  %       \item Application of the kinetic theory of multi-agent systems to population distribution on spatial networks.
  %       \item Interpretation of cities as interacting agents and interurban links as edges.
  %       \item Validation of emigration rules yielding a bimodal lognormal fit for total population distribution and a Pareto fit for the tail.
  %     \end{highlights}

      % Keywords
  %     \begin{keywords}%
        % Kinetic theory of 
  %       Multi-agent systems\sep%
  %       Boltzmann-type equations\sep%
  %       Spatial networks\sep%
  %       Population distribution\sep%
  %       Pareto distribution%
  %     \end{keywords}
      % Each keyword is seperated by \sep

  %     \maketitle

      % Main text
  %     \section{Introduction}\label{secIntroduction}

  %       The study of population distribution among cities is a phenomenon that has been sporadically analyzed in the past. The first was Auerbach \cite{Auerbach1913} who, at the beginning of the 20th century, noted a characteristic later formalized by Zipf\footnote{%
  %         Actually, as Zipf himself admits \cite[p. 546]{Zipf2016}, he was not the first to discover such a law; in any case, he popularized it to such an extent that today it is known by his name.%
  %       } \cite{Zipf2016} almost fifty years later, albeit initially in a linguistic context; indeed, Zipf's [empirical] law states
        %
  %       \begin{equation}
  %         \label{eqZipfsLaw}
  %         f\propto\frac1 r\quad\text{where}
  %         \left\{\begin{aligned}
  %           f&\text{ is the frequency of the word,}\\
  %           r&\text{ is its rank in the classification;}
  %         \end{aligned}\right.
  %       \end{equation}
        %
  %       to give a concrete example, if the tenth most frequent word appears \(2000\) times, then there exists a constant \(C\) such that \(2000\approx C/10\). The same law described in \eqref{eqZipfsLaw} can be found in many other contexts, including the distribution of population among cities if, instead of \(f\), one writes the population \(s\) \cite[Fig. 9-2 p. 375]{Zipf2016}; however, this is true only when \(s\gg1\), namely in the tail of the distribution, contrary to Zipf's original linguistic context where \eqref{eqZipfsLaw} always holds.

  %       More generally, Zipf's law can be described by a Pareto distribution: let \(S\in\RPlus\) be the random variable related to city size and let \(f_S(s)\) be its probability density function\footnote{%
  %         For \(f_S(s)\) to exist, one must strictly also assume that \(S\) is absolutely continuous, but the details are not necessary here.%
  %       }, which allows interpreting \(f_S(s)\de s\) as the number of cities with \(s\in[s,s+\de s]\);
  %       then \(S\) is said to follow a Pareto distribution in the tail if it satisfies
        %
  %       \begin{equation}
  %         \label{eqComplementaryCumulativePareto}
  %         \pazocal R(s)\equiv\int_{s}^{+\infty}f_S(z)\de z\approx\frac c{s^\beta}
  %         \quad\text{if }s\gg1,
  %       \end{equation}
        %
  %       for a certain constant \(c\in\RPlus\) and with \(\pazocal R(s)\) equal to the complementary cumulative distribution function (CCDF) of \(S\): it represents the number of cities with a population greater than or equal to \(s\) and, if the Pareto index \(\beta\in\mathbb R\) were unitary, it is analogous to the rank \(r\) in \eqref{eqZipfsLaw}.

  %       More recently \cite{Eeckhout2004,Gualandi2019HBLD}, it has been discovered that outside the tail, \(f_S(s)\) is well fitted by the lognormal distribution with probability density
        %
  %       \begin{equation}
  %         \label{eqLognormalDistribution}
  %         f_X(x)=\frac1{\ln(10)}\frac1{x\sigma\sqrt{2\pi}}\exp\left( -\frac{(\log x-\mu)^2}{2\sigma^2}\right),
  %       \end{equation}
        %
  %       where \(X\sim\Lognormal(\mu,\sigma)\); as the name suggests, the peculiarity of this distribution is that \(Y=\log X\sim\mathcal N(\mu,\sigma)\) holds, where \(\mathcal N(\mu,\sigma)\) indicates the Gaussian distribution with mean \(\mu\) and standard deviation \(\sigma\).

  %       Furthermore, Gualandi \etal{} \cite{Gualandi2019HBLD,Gualandi2019SDC} realized that the entire distribution is well fitted by a bimodal lognormal (here also denoted as bilognormal), namely a convex combination of \eqref{eqLognormalDistribution}
        %
  %       \begin{equation}
  %         \label{eqBimodalLognormalDistribution}
  %         f_S(s)\approx f_X(s)=\xi f_{X_1}(s)+(1-\xi)f_{X_2}(s),
  %       \end{equation}
        %
  %       where \(\xi\in(0,1)\), \(X_1,X_2\sim\Lognormal(\mu,\sigma)\) and \(X\sim\BiLognormal(\xi,\mu_1,\sigma_1,\mu_2,\sigma_2)\).

  %       \tikzexternalizedfigure{p}
  %         {../../Figure/.fnl/Elaborato/Capitolo1/EsempioFunzioneLognormale.tex}
  %         {Function $X \sim \text{Lognormal}(\mu, \sigma)$ for varying parameters, but with equal expected value $\mathbb{E}[X] = 10^{\mu + \sigma^2 \ln(10) / 2} = 10$, in linear scale (a) and logarithmic scale (b); instead in (c) the same distributions are reported in logarithmic space $y = \log(x)$. The dashed lines represent the expected values of each distribution.}
  %         {.65\linewidth}{figLognormalFunctionExample}

  %       Therefore, the objective of this paper is as follows: we want to simulate through the \hypertarget{acrKineticTheoryMultiAgentSystems}{Kinetic Theory of Multi-Agent Systems} (\KTMAS) the interactions between cities on a spatial graph, attempting to reproduce a population distribution with characteristics analogous to those actually measured (mainly bilognormal body and Pareto tail) in the Sardinia region of Italy. Once this goal is achieved, the migration phenomenon is interpreted through the proposed emigration laws, which are the basis of the interactions between cities.

  %       Moreover, it is emphasized that in this work, cities are considered as agents characterized by their population, which interact through the underlying structure of a [spatial] graph; this description, except for the network, is not unlike the classical one of gas molecules characterized by their velocity and position. In fact, in the literature, at least to the authors' knowledge, there are no articles dealing simultaneously with the \KTMAS{}, population distribution among cities, and graphs, despite the natural representation of these on a network; more in detail:

  %       \begin{itemize}
  %         \item some either apply the \KTMAS{} to cities without considering a natural underlying topology \cite{Gualandi2019SDC}\footnote{%
  %           In particular \cite{Gualandi2019SDC}, despite being a very interesting article that reaches equally important conclusions, excessively abstracts the interactions between cities and fails to define with sufficient clarity what is meant by \textit{neighbourhoods}.%
  %         } or, conversely, consider the graph structure without applying it \cite{Berliant2018};
  %         \item others view nodes more as places where agents live and interact, thus creating a model described by a system of large, coupled Boltzmann equations \cite{Loy2021};
  %         \item finally, the article closest to the current paper applies its theory in the context of social networks \cite{Nurisso2024}.
  %       \end{itemize}

  %       This gap is probably attributable to the fact that developments in graph theory applied to the \KTMAS{} have been quite recent and primarily focused on the node-place perspective rather than node-agent. Therefore, the most innovative aspect of this work is showing that the second perspective is as valid as the first, despite being less common; otherwise, this paper should be considered as a formal applicative exploration with few analytical in-depth studies.

  %       Overall, the manuscript is divided as follows: the second section will briefly and superficially address some notions of the \KTMAS{} mediated by graphs; then the third will showcase the numerical results for various emigration rules for the Sardinia region; finally, some conclusions are drawn with potential future developments.

  %     \section{Notes of the \protect\KTMAS{} on graphs}\label{secNotesKTMAS}

  %       In this section some hypotheses about the type of graph used and the two main Boltzmann type equations on networks are briefly discussed. For the general theory of the \KTMAS{} see \cite{Loy2025EKTMAS,PareschiIMS2013}.

  %       Before continuing, two recurrent definitions are the following ones: \(\RPlus\equiv\{x\in\mathbb R \colon x\ge0\}\subset\mathbb R\) is the subset of the real non-negative numbers, while \(\NPlus\equiv\{n\in\mathbb N \colon n\geq1\}\subset\mathbb N\) is the subset of the natural positive numbers.
        % https://en.wikipedia.org/wiki/Positive_real_numbers

  %       \subsection{Networks and cities}\label{secNetworksAndCities}

  %         First, it's important to clearly define what a graph is:

  %         \begin{dfn}[Graph]
  %           \label{defGraph}
  %           A graph is a set \(\mathcal G=(\mathcal I,\mathcal E)\) formed by the pair of the set of node indices \(\mathcal I\subseteq\NPlus\) and the set of edges \(\mathcal E\), namely pairs of indices of \(\mathcal I\): two nodes \(i,j\in\mathcal I\) are connected if and only if \((i,j)\in\mathcal E\subseteq\NPlus\times\NPlus\).
  %         \end{dfn}

  %         On the other hand, other related concepts are henceforth taken for granted though they can be found in any reference about graph theory like \cite{Barthelemy2011,vanSteen2010}.

  %         The next problem is how to exactly interpret a link between two nodes, i.e. cities, given the multitude of possible urban connections. In fact, there are many ways to represent cities by means of graphs: the first are at the \textit{intraurban} level, either considering streets as edges and their intersections as nodes \cite[§ 3.1.3.1 p. 17]{Barthelemy2011}, or representing transport networks (tram, bus, and metro) \cite[§ 3.2.1 p. 22]{Barthelemy2011}; others are more properly at the \textit{interurban} level, taking individually into consideration various transport networks (railway \cite[§ 3.1.3.2 p. 17]{Barthelemy2011}, naval \cite[§ 3.1.4 p. 19]{Barthelemy2011}, air \cite[§ 3.1.2 p. 13]{Barthelemy2011}, etc.). However, from what was stated in \cref{secIntroduction}, it is clear that in order to predict the population distribution among cities, the following observations hold true:

  %         \begin{enumerate}
  %           \item every intraurban representation must be discarded because its movements are limited to a single city;
  %           \item on the other hand, all the interurban ones must be considered simultaneously and not individually, since each person can choose different modes of transport to move around.
  %         \end{enumerate}

  %         This is why the correct network used in the simulation is the one related to commuting movements [\citealp{Barthelemy2011}, § 3.1.3.3 p. 18; \citealp{DeMontis2007}] between cities, which holistically shows all their possible connections regardless of the chosen transportation mode; furthermore, it shows realistic connections associated with daily rather than extraordinary movements (holidays, medical visits, etc.). Once the representation is fixed, it is necessary to understand the type of graph we are dealing with. Yet, even here, the answer is immediate after two further notes:

  %         \begin{enumerate}
  %           \item a city can interact with another without the latter interacting with the former;
  %           \item assuming a city can interact with another, then the opposite is always possible.
  %         \end{enumerate}

  %         Therefore the sense of direction must be considered and the potential interactions are symmetric, that is the graph in question is directed but with an undirected structure:

  %         \begin{hyp}
  %           \label{hypGraphType}
  %           The graph \(\mathcal G\) is directed and symmetric, hence the adjacency matrix \(\mathbf A\) is symmetric.
  %         \end{hyp}

  %         Furthermore, the following fundamental hypothesis holds in this discussion:

  %         \begin{hyp}
  %           \label{hypStaticGraph}
  %           The graph \(\mathcal G\) is static: \(\mathcal I\) and \(\mathcal E\) are constant over time.
  %         \end{hyp}

  %         In other words, we are only interested in predicting how the population distributes itself with respect to a topology established \textit{a priori}; naturally, this is a strong simplification since in reality the topology, especially the interurban connections in \(\mathcal E\), has clearly co-evolved along with the development of the cities themselves. Finally, it is observed that this type of graph is spatial \cite{Barthelemy2011,DeMontis2007}, meaning its nodes occupy a point in Euclidean space, even though, at the moment, this is a somewhat formal characteristic; however, it helps to keep it in mind when interpreting certain emigration rules in \cref{secNumericalResults}. Moreover, contrary to what one might think at first glance, this graph is not scale-free \cite{Barabasi1999} precisely because of its spatial nature \cite{Broido2019,Barthelemy2011}; this does not rule out the existence of nodes that are more central than others, but only implies that the maximum degree of a node is upper-bounded due to physical constraints.

  %       \subsection{Exact Boltzmann-type equation}\label{secExactBoltzmannType}

  %         The \KTMAS{} is centred around agents interacting to one another in so changing their microstate \(\mathbf x_t\in\mathbb R^n\); however, usually the precise knowledge of a single agent's microstate is neither useful nor necessary, hence why a stochastic approach is taken by denoting \(\mathbf X_t\in\mathbb R^n\) as the random variable of the microstate at time \(t\); finally, the sequence in time of microstates \(\{X_t\}_{t\in\RPlus}\) is called a stochastic process. To derive the equation for the [probability] density function \(f(\mathbf x,t)\) four assumptions are made \cite{Loy2025EKTMAS}:

  %         \begin{hyp}
  %           \label{hypBoltzmannType}
  %           Agents are characterized by the following hypotheses:

  %           \begin{enumerate}[label=B\arabic*]
  %             \item\label{hypBoltzmannTypeHomogeneousDensity} the statistical distribution \(f\) of microstates is spatially uniform;
  %             \item\label{hypBoltzmannTypeBinaryInteractions} non-binary interactions between agents are negligible;
  %             \item\label{hypBoltzmannTypeIndistinguishableAgents} the agents are indistinguishable, and each represents the ensemble;
  %             \item\label{hypBoltzmannTypeMolecularChaos} two interacting agents are statistically independent.
  %           \end{enumerate}
  %         \end{hyp}

  %         Thus a inconsistency arises from \ref{hypBoltzmannTypeIndistinguishableAgents} if agents are confused with graph nodes: the former are indistinguishable, but the underlying network makes the latter distinct due to the indices \(\mathcal I\); to solve this apparent issue it's sufficient to consider the index as a random variable \(I\in\mathcal I\), resulting in the stochastic process \(\{\mathbf X_t\}\equiv\{(I,S_t)\}_{t\in\RPlus}\) where \(S_t\in\RPlus\) is the random variable related to the population of the representative agent at time \(t\). Of course by definition \(S_t\) is the only component of \(\mathbf X_t\) to vary over time due to \cref{hypStaticGraph}.

  %         \begin{rmk}[On the numerical nature of the population]
  %           \label{rmkPopulationNumericalNature}
  %           It is natural that fractions of people do not exist and therefore strictly \(S_t\in\mathbb N\), but the most realistic choice is not always the easiest modelling-wise. At any rate, if necessary, it's possible to gauge the actual population by rounding it down to the nearest natural number.
  %         \end{rmk}

  %         Reasoning in this way, one can overall statistically describe the microscopic state \(\mathbf X_t\equiv(I,S_t)\) of the representative agent with the joint density
  %           \(f\colon\mathcal I\times\RPlus\times\RPlus\to\RPlus,\)
  %         which is discrete in \(i\in\mathcal I\) and continuous in \(s\in\RPlus\), equal to
          %
  %         \begin{equation}
  %           \label{eqUrbanJointDensity}
  %           f(i,s,t)\equiv\frac1N\sum_{j\in\mathcal I}f_j(s,t)\otimes\delta(i-j),
  %         \end{equation}
          %
  %         where \(N\equiv\vert\mathcal I\vert\) is the total number of agents in the graph, while \(\delta(\cdot)\) denotes the Dirac delta centered at the origin; on the other hand, \(f_i:\RPlus\times\RPlus\rightarrow\RPlus\) is the density of \(S_t\) relative to the \(i\)-th agent (\cref{figGraphicRepresentationKTMASNetworks}).

  %         \tikzexternalizedfigure{thb}
  %           {../../Figure/.fnl/Elaborato/Capitolo3/RappresentazioneGraficaTCSMAReti.tex}
  %           {Graphical representation of the \protect\KTMAS{} on graphs.}
  %           {\linewidth}{figGraphicRepresentationKTMASNetworks}

  %         In this way, \(f\) is effectively indistinguishable with respect to all agents and integrates over its domain to one:
          %
  %         \begin{equation*}
  %           \begin{aligned}
  %             \mathbb P_t(\mathcal I\times\RPlus)
  %             &=\int_{\mathcal I}\int_{\RPlus}f(i,s,t)\de s\de i
  %             =\frac1N\sum_{i\in\mathcal I}\int_{\mathcal I}\int_{\RPlus}
  %             f_i(s,t)\de s\otimes\delta(i-j)\de i
  %             \\&=\frac1N\sum_{i\in\mathcal I}\int_{\mathcal I}\delta(i-j)\de i
  %             =\frac1N\sum_{i\in\mathcal I}1=\frac NN=1,
  %             \quad\forall t\geq0,
  %           \end{aligned}
  %         \end{equation*}
          %
  %         allowing to preserve the natural distinguishability of the agents as well, induced by the underlying topology through the individual densities \(f_i\).

  %         When two agents interact they alter their respective microstate following the following general scheme
          %
  %         \eqtag*{\([\text{RI}]_S\)}{%
  %           \label{eqUrbanInteractionRules}
  %           [\text{RI}]_S
  %           \left\{\begin{aligned}
  %             S'_t&=\psi(S_t,I,S^*_t,I_*),\\
  %             S^{*\prime}_t&=\psi_*(S_t,I,S^*_t,I_*).
  %           \end{aligned}\right.
  %           \quad\text{where the functions}\;\;
  %           \begin{aligned}
  %             \psi&\colon\RPlus\times\NPlus\times\RPlus\times\NPlus\to\RPlus,\\
  %             \psi_*&\colon\RPlus\times\NPlus\times\RPlus\times\NPlus\to\RPlus
  %           \end{aligned}
  %         }%
  %         describe the post-interaction microstate; however, it's not certain that two randomly picked agents can interact: this can be taken into account with the Bernoulli
          %
  %         \eqtag{\([\text{Ber}]^{\mathbf A}_S\)}{%
  %           \label{eqUrbanInteractionBernoulli}
  %           \Theta\sim\Bernoulli(A(I,I_*)\Delta t),
  %           \text{ with }
  %           \newcommand\hspz{\hspace{1em}}
  %           A(I,I_*)\equiv a_{I,I_*}=
  %           \left\{\begin{aligned}
  %             1&\hspz\text{if }(I,I_*)\in\mathcal E,\\
  %             0&\hspz\text{otherwise},%(I,I_*)\notin\mathcal E,
  %           \end{aligned}\right.
  %         }%
  %         where \(A\colon\mathcal I\times\mathcal I\to\{0,1\}\) is a function that associates each pair of indices \((I,I_*)\) with the corresponding element of the adjacency matrix; defined in this way, \ref{eqUrbanInteractionBernoulli} inhibits \textit{a priori} the interaction if the two agents are not connected.

  %         \begin{rmk}
  %           By definition of the Bernoulli distribution, it is necessary to impose \(A(I,I^*)\Delta t\leq1\), but given that \(A(I,I^*)\in\{0,1\}\) this translates into the natural condition that \(\Delta t\leq1\).
  %         \end{rmk}

  %         Combining \ref{eqUrbanInteractionBernoulli} and \ref{eqUrbanInteractionRules}, the following algorithm is obtained
          %
  %         \eqtag*{\([\text{AR}]_S\)}{
  %           \label{eqUrbanInteractionAlgorithm}
  %           [\text{AR}]_S\left\{
  %           \begin{aligned}
  %             S_{t+\Delta t}&=(1-\Theta)S_t+\Theta S'_t,\\
  %             S^*_{t+\Delta t}&=(1-\Theta)S^*_t+\Theta S^{*\prime}_t,
  %           \end{aligned}\right.
  %         }%
  %         consistently with the fact that in an urban context, if a city-node interacts with another, both must vary their state.

  %         \begin{rmk}
  %           \label{rmkActionReactionAlgorithm}
  %           The acronym \ref{eqUrbanInteractionAlgorithm} stands for Action-Reaction and it assumes that every interaction (action), i.e., \(\Theta=1\), necessarily modifies the states of both involved agents (reaction); in certain contexts \cite[(2.5) p. 5]{Nurisso2024} this rule can be altered such that only the interacting agent is modified, leading to Action-Action interactions.
  %         \end{rmk}

  %         Let now \(\Phi\colon\mathcal I\times\RPlus\to\mathbb R\) be an arbitrary observable quantity, then by averaging it while evaluating \ref{eqUrbanInteractionAlgorithm} it can be shown \cite{Nurisso2024} to lead to the homogeneous Boltzmann-type equation
          %
  %         \eqtag{\([\text{EtB}]^{\mathbf A}_S\)}{
  %           \label{eqUrbanAsymmetricHomogeneousWeakFormBoltzmannType}
  %           \derS{}t\int_{\mathcal I}\int_{\RPlus}
  %           \Phi f\de s\de i=
  %           \int_{\mathcal I^2}\int_{\RPlus^2}A(i,i_*)
  %           \frac{\Phi'+\Phi_*'-\Phi-\Phi_*}2
  %           ff_*\de s\de s_*\de i\de i_*,
  %         }%
  %         where an abbreviated notation has been used in which the arguments are omitted and the asterisks as well as the primes change their type: for instance \(\Phi_*'\equiv\Phi(i_*,s_*')\) while \(f_*\equiv f(i_*,s_*,t)\), and so forth.

  %       \subsection{Approximated Boltzmann-type equation}\label{secApproximatedBoltzmannType}

  %         Recently an approximation of \ref{eqUrbanAsymmetricHomogeneousWeakFormBoltzmannType} has been proposed \cite{Nurisso2024} based on the adjacency matrix \(\mathbf A\):

  %         \begin{dfn}[Degree matrix \(\mathbf B\)]
  %           \label{defDegreeMatrix}
  %           Let \(\mathbf B\in\mathbb R^{N\times N}\) be the rank-one matrix
  %           \[
  %             \mathbf B\equiv\frac{\mathbf k^+(\mathbf k^-)^\top}{D_N}
  %             \quad\text{with}\;
  %             \left\{\begin{aligned}
  %               \mathbf k^+&\equiv(k^+_1,k^+_2,\ldots,k^+_N)^\top=\mathbf A\mathbf 1\in\mathbb R^N,\\
  %               \mathbf k^-&\equiv(k^-_1,k^-_2,\ldots,k^-_N)^\top=\mathbf A^\top\mathbf 1\in\mathbb R^N,
  %             \end{aligned}\right.
  %           \]
            %
  %           being the out-degrees and in-degrees respectively, and the constant \(D_N\equiv\mathbf1^\top\mathbf A\mathbf1\). Then \(\mathbf B\) approximates adjacency matrix \(\mathbf A\) since by definition they share the same in- and out-degrees:
            %
  %           \begin{equation}
  %             \label{eqApproximationCriterionA}
  %             \left.\begin{aligned}
  %               \mathbf B\mathbf1
  %               &=\frac1{D_N}\mathbf k^+(\mathbf k^-)^\top\mathbf1
  %               =\mathbf k^+\frac{D_N}{D_N}=\mathbf k^+=\mathbf A\mathbf1\\
  %               \mathbf B^\top\mathbf1
  %               &=\frac1{D_N}\mathbf k^-(\mathbf k^+)^\top\mathbf1
  %               =\mathbf k^-\frac{D_N}{D_N}=\mathbf k^-=\mathbf A^\top\mathbf1
  %             \end{aligned}\right\}
  %             \ \ \implies\ \ \mathbf B\approx\mathbf A.
  %           \end{equation}
            %
  %           Furthermore, since \(\mathbf A\) is symmetric (\cref{hypGraphType}), \(\mathbf k^+=\mathbf k^-=\mathbf k\) holds and therefore \(\mathbf B=\mathbf k\mathbf k^\top/D_N.\)
  %         \end{dfn}

  %         By using \(\mathbf B\) \ref{eqUrbanAsymmetricHomogeneousWeakFormBoltzmannType} becomes
          %
  %         \eqtag{\([\text{EtB}]^{\mathbf B}_S\)}{
  %           \label{eqApproximatedUrbanAsymmetricHomogeneousWeakFormBoltzmannType}
  %           \derS{}t\int_{\mathcal I}\int_{\RPlus}
  %           \Phi g\de v \de i=
  %           \int_{\mathcal I^2}\int_{\RPlus^2}B(i,i_*)
  %           \frac{\Phi'+\Phi_*'-\Phi-\Phi_*}2
  %           gg_*\de s\de s_*\de i\de i_*,
  %         }%
  %         in which the density \(g\) is used to distinguish it from \(f\) though it has its same definition found in \cref{eqUrbanJointDensity}; \ref{eqApproximatedUrbanAsymmetricHomogeneousWeakFormBoltzmannType} can be shown \cite[Theorem 6.1.]{Nurisso2024} to converge in the limit \(N\to+\infty\) to
          %
  %         \eqtag{\([\text{EtB}]^k_S\)}{
  %           \label{eqRelaxedGraphBoltzmannType}
  %           \derS{}t\int_0^1\int_{\RPlus}\varphi\underhat d\ \de s\de\underhat k=
  %           \int_0^1\int_0^1\int_{\RPlus^2}\frac{\underhat k\underhat k_*}{\bar D}
  %           \frac{\varphi'+\varphi'_*-\varphi-\varphi_*}2
  %           \underhat d\underhat d^*\ \de s\de s_*\de\underhat k\de\underhat k_*.
  %         }%
  %         where

  %         \begin{itemize}
  %           \item \(\Phi\) has been reduced to an arbitrary observable \(\varphi\colon\RPlus\to\mathbb R\) function only of the population \(s\);
  %           \item \(\underhat d(s,\underhat k,t)\) is the limit of the distribution of normalized degrees \(
  %             \underhat d_N\colon\RPlus\times\underhat{\mathcal K}\times\RPlus\to\RPlus,
  %           \)
  %           where
  %           \[
  %             \underhat{\mathcal K}\equiv
  %             \left\{k/N\ |\ k=0,1,\ldots,N\right\}
  %             \subseteq[0,1];
  %           \]
  %           \item the constant \(\bar D\) is defined as \(\bar D\equiv\int_0^1\int_{\RPlus}\underhat k\underhat d\ \de s\de\underhat k\).
  %         \end{itemize}

  %         The implications of \ref{eqRelaxedGraphBoltzmannType} are profound: approximating \(\mathbf A\) with \(\mathbf B\) coincides with losing the distinguishability induced by the graph as the microstate loses the index \(i\) in place of the degree \(k\), thereby relaxing the topology (\cref{figGraphicalRepresentationApproximation}) which does not vanish but remains solely as a degree distribution; such a reduction is akin to blurring the connection details among the vertices: one transitions from a specific graph to a class of graphs sharing the same degree distribution. Therefore \ref{eqRelaxedGraphBoltzmannType} is closer the original context of the \KTMAS{} making it a good candidate to approximate \ref{eqUrbanAsymmetricHomogeneousWeakFormBoltzmannType}.

  %         \tikzexternalizedfigure{htb}
  %           {../../Figure/.fnl/Articolo/Capitolo2/RappresentazioneGraficaApprossimazione.tex}
  %           {Graphical representation of the \protect\KTMAS{} approximation on graphs; for visual simplicity an undirected graph has been drawn, instead of the more correct directed one (\cref{hypGraphType}), for the approximated and relaxed topology.}
  %           {.85\linewidth}{figGraphicalRepresentationApproximation}

  %         \begin{rmk}
  %           \ref{eqRelaxedGraphBoltzmannType} justifies why one cannot consider a generic matrix \(\mathbf C\) approximating \(\mathbf A\) in any way: it is not guaranteed to relax the topology and lead to an equation analogous to it like \(\mathbf B\) did.
  %         \end{rmk}

  %       \subsection{Emigration rules}\label{secEmigrationRules}

  %         First, by using realizations of all considered random variables, writing them in lowercase and omitting time dependence—equations, \ref{eqUrbanInteractionRules} can be written as follows\footnote{%
  %           They were partially inspired by \cite[(2.1) p. 223]{Gualandi2019SDC}.%
  %         }:
          %
  %         \eqtag*{\([\text{REI}]\)}{
  %           \label{eqExplicitUrbanInteractionRules}
  %           [\text{REI}]
  %           \left\{\begin{aligned}
  %             \psi(s,i,s_*,i_*)&\equiv s\bigl(1-E(s,i,s_*,i_*)\bigr)\\
  %             \psi_*(s,i,s_*,i_*)&\equiv s_*+sI(s,i,s_*,i_*)
  %           \end{aligned}\right.
  %         }%
  %         where \(E,I\colon\RPlus\times\mathcal I\times\RPlus\times\mathcal I\to\RPlus\), are the emigration and immigration rates, respectively, whose arguments are implicitly assumed hereafter for brevity like in \ref{eqUrbanAsymmetricHomogeneousWeakFormBoltzmannType}.

  %         \begin{rmk}
  %           One could also more realistically define \( \psi(s,i,s_*,i_*\gamma)\equiv s\bigl(1-E(s,i,s_*,i_*)+\gamma\bigr)\) with \(\gamma\) being the random variable of the variation rate, i.e. the sum of the death and birth rates, representing the natural stochastic fluctuations of the population; \ref{eqExplicitUrbanInteractionRules} just assumes that \(\gamma=0\): the birth and death rates are always equal.
  %         \end{rmk}

  %         We further characterize \(E\) and \(I\) with two hypotheses:

  %         \begin{hyp}[Positivity constraints]
  %           \label{hypEmigrationRateLimits}
  %           To ensure a physically meaningful (i.e., non-negative) post-interaction state \(s'\), it is assumed that \(E\) is upper-bounded by the constant \(\lambda\equiv\sup E\in(0,1)\) and lower-bounded by the origin:
            %
  %           \begin{equation*}
  %             0\leq E\leq\lambda<1,\ \ \forall s,\ \forall s_*,\ \forall i,\ \forall i_*.
  %           \end{equation*}
            %
  %           Conversely, the post-interaction state \(s_*'\) is not restricted since it is by definition always positive regardless of the immigration rate \(I\in\RPlus\).
  %         \end{hyp}

  %         \begin{hyp}[Conservation of average population]
  %           \label{hypConservationAveragePopulation}
  %           It is postulated that the \ref{eqExplicitUrbanInteractionRules} conserve the total population on average: 
  %           \(
  %             s+s_*=\langle s+s_*\rangle=
  %             \langle s'+s_*'\rangle
  %             =s-sE+s_*+sI,
  %           \)
  %           which yields \(I\equiv E\), a reasonable result since emigration and immigration are relative phenomena (inverting \(s\) and \(s_*\) would produce the opposite case).
  %         \end{hyp}

  %         \begin{rmk}
  %           \label{rmkEqualityInteractionEmigration}
  %           The \ref{eqExplicitUrbanInteractionRules} are not yet complete since the emigration rule \(E\) has not been specified, the choice of which can be seen as the interaction rule itself; this is done later on in \cref{secNumericalResults} when the results are actually discussed.
  %         \end{rmk}

  %     \section{Implementation and simulations}\label{secImplementationSimulation}

  %       \subsection{Monte Carlo method}\label{secMonteCarloMethod}

  %         As is clear from \cref{secNotesKTMAS}, the main objective is to obtain the density \(f\) in \ref{eqUrbanAsymmetricHomogeneousWeakFormBoltzmannType} which is, however, quite difficult to handle due to, among others, its integro-differential nature. Because of this, and as usual in the \KTMAS{}, a simpler procedure is adopted: knowing the algorithm \ref{eqUrbanInteractionAlgorithm}, which governs the binary interactions between agents, all multiple collisions can therefore be directly simulated using a Monte Carlo method of the Nanbu-Babovsky type, described in detail in \cref{algMonteCarlo} in which \(T>0\) is the final simulation time, while \(P\) is the initial total population.

  %         \begin{rmk}
  %           In line \ref{algMonteCarloDistribuzioneNesima}, \(\bar z_N\) is an arbitrary marginal density w.r.t. the population: 
  %           \(
  %             \bar z_N(s,t)
  %             \equiv\int_{\mathcal I}z(i,s,t)\ \de i,
              % =\sum_{i\in\mathcal I}z_i(s,t)/N;
  %           \) 
  %           where \(z\) is similarly defined as \cref{eqUrbanJointDensity} and should be understood as \(f\) with \(\mathbf A\) and as \(g\) with \(\mathbf B\).
  %         \end{rmk}

  %         \begin{rmk}
  %           More generally, in line \ref{algMonteCarloDistribuzioneIniziale}, \(\pazocal S_0\) can be sampled from an initial density \(\bar z_0\); however, since none is known for the population distribution, it was decided to define \(\pazocal S^0\) as a uniform vector with respect to an initial total population \(P\).
  %         \end{rmk}

  %         \begin{algorithm}[tb]%h
  %           \caption{Nanbu-Babovsky type algorithm \ref{eqUrbanInteractionAlgorithm}}
  %           \label{algMonteCarlo}
  %           \KwData{%
  %             \justifying
  %             \(N\in\NPlus\), (number of agents), %
  %             \(\Delta t\leq 1\), (time step), %
  %             \(\sigma\), (fluctuation), %
  %             \(T>0\), (final time), %
  %             \(P\), (initial total population), %
  %             \(\mathbf A\) (adjacency matrix) %
  %             and \(\mathbf B\) (approximated adjacency matrix)%
  %             ;\vspace{.5em}%
  %           }%
  %           \(\pazocal S^0\gets(s_1^0,s_2^0,\dots,s_N^0)\equiv(P/N)\mathbf 1\in\RPlus^N\)\;\label{algMonteCarloDistribuzioneIniziale}
  %           \For{\(n=0,1,2,\dots,\left\lfloor T/\Delta t\right\rfloor-1\)}{
              %
  %             \(p\gets\) independent permutation of \(\{1,2,\ldots,N\}\)\;
              %
  %             \For(\tcp*[h]{number of pairs}){\(c=1,2,\dots,\left\lfloor N/2\right\rfloor\)}{
                %
  %               \(i\gets p(c)\); \(r\gets p(\left\lfloor N/2\right\rfloor+c)\)\;
                %
                % Si campiona \(\Theta\sim\DstrBernoulli(A(i,j)\Delta t)\)\;
  %               \uIf{exact}{
  %                 \(\Theta\gets\Bernoulli(A(i,r)\Delta t)\)\;
  %               }
  %               \Else(\tcp*[h]{it is approximated}){
  %                 \(\Theta\gets\Bernoulli(B(i,r)\Delta t)\)\;
  %               }
                %
  %               \uIf{\(\Theta=1\)}{
  %                 \(E\gets E(s_i^n,i,s_r^n,r)\)\;
                  % \(\gamma\gets\DstrGamma((1-E)^2/\sigma^2,\sigma^2/(1-E))\)\;
                  %
  %                 \(s_i^{n+1}\gets s_i^n(1-E)\)\tcp*[r]{interacting city}
  %                 \(s_r^{n+1}\gets s_r^n+s_i^nE\)\tcp*[r]{receiving city}
  %               }
  %               \Else{
  %                 \(s_i^{n+1}\gets s_i^n\); \(s_r^{n+1}\gets s_r^n\)\;
  %               }
                %
                % \(s_i^{n+1}\gets s_i^n(1-\Theta)+\Theta\psi(s_i^n,i,s_j^n,j,\gamma)\)\;
                % \(s_j^{n+1}\gets s_j^n(1-\Theta)+\Theta\psi_*(s_i^n,i,s_j^n,j)\)\;
  %             }
              %
  %             \(\pazocal S^{n+1}\gets(s_1^{n+1},s_2^{n+1},\dots,s_N^{n+1})\)\;
  %             \(\bar z_N(s,(n+1)\Delta{t})\gets\) histogram of \(\pazocal S^{n+1}\)\;\label{algMonteCarloDistribuzioneNesima}
  %           }
  %         \end{algorithm}

  %       \subsection{Numerical results}\label{secNumericalResults}

  %         \subsubsection{Foreword}

  %           Firstly, all data used was taken from the \href{https://en.wikipedia.org/wiki/Italian_National_Institute_of_Statistics}{Italian National Institute of Statistics} (\textit{ISTAT}) in the year 199I and limited to the Sardinia region, choice primarily inspired by \cite{DeMontis2007}; in detail, \cite{ISTATPendolarismo1991} provided the data to build the graph and its adjacency matrix of the commuting matrix, while \cite{ISTATCensimento1991} the actual population distribution to which the results were also compared. Therefore, given the importance of Sardinia, the main data of the region are listed in \cref{tabSardiniaDataParameters} along with the fitted parameters of the Pareto and bilognormal distributions.

  %           \begin{table}[width=.7\linewidth,cols=4,pos=t]
  %             \centering
              % \captionsetup{width=6.5cm}
  %             \caption{Data \cite{ISTATCensimento1991,ISTATPendolarismo1991} and parameters of Sardinia in 1991.}
  %             \begin{tabular*}{\tblwidth}{@{} CCCCCCCC @{}}
  %               \toprule
  %               \(N\)&\(P\)&\(\beta\)&\(\xi\)&
  %               \(\mu_1\)&\(\sigma_1\)&\(\mu_2\)&\(\sigma_2\)\\
  %               \midrule
  %               \(\num{375}\)&\(\approx\num{1648248}\)&\(\num{1.273}\)&\(\num{0.929}\)&
  %               \(\num{3.227}\)&\(\num{0.408}\)&\(\num{4.099}\)&\(\num{0.556}\)\\
  %               \bottomrule
  %             \end{tabular*}
  %             \label{tabSardiniaDataParameters}
  %           \end{table}

  %           \begin{rmk}
  %             With minor, medium, and major we refer to centres with small, medium, and large degrees respectively.
  %           \end{rmk}

  %           Regarding the plots, their complex definitions are not necessary to understand the results shown hereafter; hence, they are omitted but still present in Appendix \ref{apxPlotDefinitions} for the interested reader. Nonetheless, only three features have to be kept in mind: due to the stochastic nature of \cref{algMonteCarlo}, to have significant results \(R=100\) simulations have to be considered and later averaged in some way; moreover the time step is always \(\Delta t=0.01\) with \(N_t\) number of time steps; finally, for brevity sake, only the best fitting configuration empirically found by the authors is analyzed.
            % https://english.stackexchange.com/questions/270288/timestep-time-step-time-step-which-variant-to-use

  %         \subsubsection{Size rule}\label{secSER}

  %           \tikzexternalizedfigure{p}[-.5em]
  %             {../../Figure/.fnl/Articolo/Capitolo3/20RETConfigurazioneRiferimento/Stabile/Infopagina.tex}
  %             {Study of the stable reference configuration of \ref{eqSizeEmigrationRule} with parameters from \cref{tabSERParameters}; for the explanation of the graphs, see Appendix \ref{apxPlotDefinitions}.}
  %             {.6\linewidth}{figSERStableReferenceConfiguration}

  %           \tikzexternalizedfigure{p}[-1em]
  %             {../../Figure/.fnl/Articolo/Capitolo3/20RETConfigurazioneRiferimento/Instabile/Infopagina.tex}
  %             {Study of the unstable reference configuration of \ref{eqSizeEmigrationRule} with parameters from \cref{tabSERParameters}; for the explanation of the graphs, see Appendix \ref{apxPlotDefinitions}.}
  %             {.62\linewidth}{figSERUnstableReferenceConfiguration}

  %           A first possibility\footnote{%
  %             It is a Hill function of order \(\alpha\), partially inspired by \cite[(2.2), p. 223, § 2, and (4.5), p. 228, § 4]{Gualandi2019SDC}, widely applied in the Biomedical field.%
  %           } consists in%
            %
  %           \eqtag{\([\text{ER}]_S\)}{
  %             \label{eqSizeEmigrationRule}
  %             E(s,s_*)\equiv\lambda\frac{(s_*/s)^\alpha}{1+(s_*/s)^\alpha},
  %           }%
  %           where \(\lambda\in(0,1)\) (\cref{hypEmigrationRateLimits}) and \(\alpha\in\mathbb R^+\). The logic is that the emigration rate towards the receiving city \(s_*\) is greater the larger its relative population, given by the ratio \(s_*/s\), compared to the interacting city \(s\).

  %           In total, two reference configurations described by the parameters in \cref{tabSERParameters} were studied, while the results are illustrated in \cref{figSERStableReferenceConfiguration,figSERUnstableReferenceConfiguration}; these describe two behaviours:

  %           \begin{enumerate}
  %             \item a stable one but not distant from the initial uniform distribution, and
  %             \item an unstable one in which \textit{on average} the medium-sized centres depopulate.
  %           \end{enumerate}

  %           Mind that «On average» is emphasized as in single realizations it is much more likely that the larger centres depopulate, whereas over multiple realizations at least one with a large size manifests, thus balancing the average; this is also the reason why only 10 iterations were considered instead of the usual 100 in \cref{tabSERParameters}, given that the average tends to mask the depopulation, which is nevertheless inevitable but simply slower. In addition to \crefsubfigInt{figSERUnstableReferenceConfiguration}{j}{l}, the histograms in \crefsubfigInt{figSERUnstableReferenceConfiguration}{a}{c} also confirm this tendency in the form of a tail that elongates the distribution into negative values, corresponding to the exponents of the cities that depopulate.

  %           More generally, it has been verified, albeit only empirically, that for every \(\lambda\) there exists an \(\bar\alpha\in\RPlus\) such that \(\alpha<\bar\alpha\) results in a stable configuration while \(\alpha\geq\bar\alpha\) yields an unstable one.

            % \begin{table}[width=.4\linewidth,cols=4,pos=htb]
            %   \centering
            %   \caption{Parameters of \cref{figSERStableReferenceConfiguration,figSERUnstableReferenceConfiguration}.}
            %   \begin{tabular*}{\tblwidth}{@{} CCCCC @{}}
            %     \hline
            %     &\(\lambda\)&\(\alpha\)&\(N_t\)&\(R\)\\
            %     \hline
            %     Stable&\(\num{.1}\)&\(\num{.5}\)&\(\num{1e7}\)&\(\num{100}\)\\
            %     Unstable&\(\num{.1}\)&\(\num{1.5}\)&\(\num{1e6}\)&\(\num{10}\)\\
            %     \hline
            %   \end{tabular*}
            %   \label{tabSERParameters}
            % \end{table}

  %         \subsubsection{Size-degree rule}\label{secSDER}

  %           \tikzexternalizedfigure{p}
  %             {../../Figure/.fnl/Articolo/Capitolo3/20RETGConfigurazioneRiferimento/Infopagina.tex}
  %             {Study of the reference configuration of \ref{eqSizeDegreeEmigrationRule} with parameters from \cref{tabSDERParameters}; for the explanation of the graphs, see Appendix \ref{apxPlotDefinitions}.}
  %             {.58\linewidth}{figSDERReferenceConfiguration}

  %           In order to correct \ref{eqSizeEmigrationRule}, the relative degree ratio can also be introduced:
            %
  %           \eqtag{\([\text{ER}]_{SK}\)}{
  %             \label{eqSizeDegreeEmigrationRule}
  %             E_{SK}(s,i,s_*,i_*)\equiv\lambda\frac{[(s_*/s)(k_{i_*}/k_i)]^\alpha}{1+[(s_*/s)(k_{i_*}/k_i)]^\alpha}.
  %           }%
  %           Thus, the logic becomes that the emigration rate towards the receiving city \((s_*,k_*)\) is higher the more populous and connected it is compared to the interacting city \((s,k)\)\

  %           \newsavebox{\stabularbox}
  %           \begin{table}[pos=t]
  %             \caption{Parameters of all the reference configurations studied.}
  %             \label{tabParmetersReferenceConfigurations}
              %
  %             \newcommand\subtabular[4]{
  %               \sbox{\stabularbox}{\begin{tabular}{@{} #1 @{}}#2\end{tabular}}%
  %               \begin{subtable}[b]{\wd\stabularbox}
  %                 \centering
  %                 \begin{tabular}{@{} #1 @{}}#2\end{tabular}
  %                 \caption{#3}
  %                 \label{#4}
  %               \end{subtable}
  %             }\newcommand\tableindent{\hskip 0pt plus .4fill}%
  %             \tableindent
  %             \subtabular{ccccc}{%
  %               \toprule
  %               &\(\lambda\)&\(\alpha\)&\(N_t\)&\(R\)\\
  %               \midrule
  %               Stable&\(\num{.1}\)&\(\num{.5}\)&\(\num{1e7}\)&\(\num{100}\)\\
  %               Unstable&\(\num{.1}\)&\(\num{1.5}\)&\(\num{1e6}\)&\(\num{10}\)\\
  %               \bottomrule
  %             }{\cref{figSERStableReferenceConfiguration,figSERUnstableReferenceConfiguration}.}
  %             {tabSERParameters}
  %             \hfill
  %             \subtabular{ccc}{%
  %               \toprule
  %               \(\lambda\)&\(\alpha\)&\(N_t\)\\
  %               \midrule
  %               \(\num{.15}\)&\(\num{.63}\)&\(\num{3e6}\)\\
  %               \bottomrule
  %             }{\cref{figSDERReferenceConfiguration}.}
  %             {tabSDERParameters}
  %             \hfill
  %             \subtabular{ccc}{%
  %               \toprule
  %               \(\lambda\)&\(\alpha\)&\(N_t\)\\
  %               \midrule
  %               \(\num{.18}\)&\(\num{.44}\)&\(\num{3e6}\)\\
  %               \bottomrule
  %             }{\cref{figSSERReferenceConfiguration}.}
  %             {tabSSERParameters}
  %             \tableindent
  %             \mbox{}
  %           \end{table}

            % \begin{table}[width=.2\linewidth,cols=4,pos=htb]
            %   \centering
            %   \caption{Parameters of \cref{figSDERReferenceConfiguration}.}
            %   \begin{tabular*}{\tblwidth}{@{} CCC @{}}
            %     \hline
            %     \(\lambda\)&\(\alpha\)&\(N_t\)\\
            %     \hline
            %     \(\num{.15}\)&\(\num{.63}\)&\(\num{3e6}\)\\
            %     \hline
            %   \end{tabular*}
            %   \label{tabSDERParameters}
            % \end{table}

  %           The studied reference configuration is described by the parameters in \cref{tabSDERParameters}, while the results are illustrated in \cref{figSDERReferenceConfiguration}. Overall, they are better than those obtained using \ref{eqSizeEmigrationRule}: 

  %           \begin{itemize}
  %             \item \crefsubfigSep{figSDERReferenceConfiguration}{a}{{d}}{j} illustrate that the approximation corresponds almost perfectly to the exact simulation;
  %             \item in \crefsubfigInt[ and ]{figSDERReferenceConfiguration}{b}{c} the simulated distributions are fairly similar to the real ones, whereas in \crefsubfigInt[ and ]{figSDERReferenceConfiguration}{e}{f} the Pareto indices are reasonably close;
  %             \item with the exception of the real case in which they coincide, in \crefsubfigInt{figSDERReferenceConfiguration}{g}{i} the bilognormal follows the tail better towards the end, although it coincides with the Pareto in the first half;
  %             \item in \crefsubfigInt[ and ]{figSDERReferenceConfiguration}{k}{l} the scatter plots present a similar positive correlation, but much tighter compared to the real one: it is almost a one-to-one correlation, which is, of course, unrealistic.
  %           \end{itemize}

  %         \subsubsection{Size-strength rule}\label{secSSER}

  %           \tikzexternalizedfigure{p}
  %             {../../Figure/.fnl/Articolo/Capitolo3/20RETFConfigurazioneRiferimento/Infopagina.tex}
  %             {Study of the reference configuration of \ref{eqSizeStrengthEmigrationRule} with parameters from \cref{tabSSERParameters}; for the explanation of the graphs see Appendix \ref{apxPlotDefinitions}.}
  %             {.577\linewidth}{figSSERReferenceConfiguration}

  %           One may think to substitute the strengths, i.e. the total amount of commuters, in place of the degrees in \ref{eqSizeDegreeEmigrationRule}:
            %
  %           \eqtag{\([\text{ER}]_{SW}\)}{
  %             \label{eqSizeStrengthEmigrationRule}
  %             E_{SW}(s,i,s_*,i_*)\equiv\lambda\frac{[(s_*/s)(w_{i_*}/w_i)]^\alpha}{1+[(s_*/s)(w_{i_*}/w_i)]^\alpha},
  %           }%
  %           meaning, in this way, that there is an emigration rate towards the receiving city \((s_*,w_*)\) that is greater the more populous and heavily trafficked it is compared to the interacting city \((s,w)\).

            % \begin{table}[width=.2\linewidth,cols=4,pos=htb]
            %   \centering
            %   \caption{Parameters of \cref{figSSERReferenceConfiguration}.}
            %   \begin{tabular*}{\tblwidth}{@{} CCC @{}}
            %     \hline
            %     \(\lambda\)&\(\alpha\)&\(N_t\)\\
            %     \hline
            %     \(\num{.18}\)&\(\num{.44}\)&\(\num{3e6}\)\\
            %     \hline
            %   \end{tabular*}
            %   \label{tabSSERParameters}
            % \end{table}

  %           The analyzed reference configuration is described by the parameters in \cref{tabSSERParameters}, while the relative results are shown in \cref{figSSERReferenceConfiguration}. They are in fact the best ones: 

  %           \begin{itemize}
  %             \item \crefsubfigSep{figSSERReferenceConfiguration}{a}{{d}}{j} confirm that the approximation corresponds almost perfectly to the exact simulation;
  %             \item in \crefsubfigInt[ and ]{figSSERReferenceConfiguration}{b}{c} the simulated distributions are almost identical to the real one, and the relative Pareto indices in \crefsubfigInt[ and ]{figSSERReferenceConfiguration}{e}{f} also correspond;
  %             \item \crefsubfigInt[ and ]{figSSERReferenceConfiguration}{h}{i} do not deviate significantly from the previous configuration;
  %             \item in \crefsubfigInt[ and ]{figSSERReferenceConfiguration}{k}{l} the scatter plots, besides having the same positive correlation with respect to the real one, show greater spread, correctly capturing the size differentiation for a given degree.
  %           \end{itemize}

  %           The only negative aspect is found in \crefsubfigInt[ and ]{figSSERReferenceConfiguration}{k}{l}: the emigration rule fails to correctly predict on average the sizes of the lower portion of the real scatter plot; in other words, the prediction overestimates the size of the minor centres and slightly underestimates that of the major centres. Yet, on second thought, this is a characteristic also shared by the other case in \crefsubfigInt[ and ]{figSDERReferenceConfiguration}{k}{l}, as can be noted by observing the slope of the line described by the scatter plot. We believe its cause to be found in the absence of intra-urban commuters in the strength definition which only considers inter-urban ones; this would increase only the strength of the medium and major centres, but only marginally that of the minor ones, correcting in theory the slope.

  %       \subsection{Interpretations}\label{secInterpretazioni}

  %         First of all, it is necessary to elucidate a characteristic shared by all previous emigration rules: from an interpretative point of view, they are myopic, meaning they do not see the details but only the common factors that might govern migrations. Therefore, given the arbitrariness of potential interpretations, it has been decided in this subsection not to express details, but only to highlight which factors most likely influence migration in light of the comparisons with real data in \cref{figSERStableReferenceConfiguration,figSERUnstableReferenceConfiguration,figSDERReferenceConfiguration,figSSERReferenceConfiguration}.

  %         \subsubsection{Size rule - \ref{eqSizeEmigrationRule}}

  %           From the results of \ref{eqSizeEmigrationRule}, it is clear that considering only the relative population is not sufficient to describe the size distribution among cities. The reason for its degeneration can be attributed to a binary behaviour: if \(s\le s_*\) then \(s_*\) will always be an attractive pole for \(s\) [and vice versa]; therefore, as soon as a city loses even a little population compared to all its neighbouring nodes, its fate is sealed to depopulate. With this logic, the depopulation \textit{on average} of the medium centres can be explained: the major centres deplete the populations of the medium centres, rendering them relatively repellent poles right from the beginning of the simulation; in the short term, the medium centres depopulate to such an extent that they become repellent even with respect to the minor centres; the minor centres never interface directly with the major ones, therefore they do not initially depopulate like the medium ones. The result of this dynamics is shown in \crefsubfigInt[ and ]{figSERUnstableReferenceConfiguration}{k}{l}, at least in the unstable case.

            % \begin{enumerate}
            %   \item the major centres deplete the populations of the medium centres, rendering them relatively repellent poles right from the beginning of the simulation;
            %   \item in the short term, the medium centres depopulate to such an extent that they become repellent even with respect to the minor centres;
            %   \item the minor centres never interface directly with the major ones, therefore they do not initially depopulate like the medium ones.
            % \end{enumerate}

  %           Conversely, in the stable case, it is unclear what this behaviour can be attributed to: certainly \(\alpha\) is not small enough to make the Hill function in \ref{eqSizeEmigrationRule} constant (\cref{figDegreeAHillFunction}) and equal to \(\lambda/2\), especially considering that \ref{eqSizeStrengthEmigrationRule} in the reference configuration from \cref{tabSSERParameters} also presents a lower value of \(\alpha\); consequently, the only conclusion that can be drawn is that the problem lies in the emigration rule itself.

  %           \tikzexternalizedfigure{t}
  %             {../../Figure/.fnl/Elaborato/Capitolo4/FunzioneHillGradoa.tex}
  %             {Hill function as the degree \(\alpha\) varies.}
  %             {.775\linewidth}{figDegreeAHillFunction}

  %           Overall, therefore, it can be deduced that migration most likely does not depend \textit{solely} on relative population.

  %         \subsubsection{Size-degree rule - \ref{eqSizeDegreeEmigrationRule}}

  %           These results are decidedly more realistic; however, contrary to relative population, the introduction of the relative degree via the ratio \(k_{i_*}/k_i\) is not intuitive. In other words, one might wonder why it is more reasonable to postulate that a city with a low degree, i.e., less connected, also has a low population. The fear is that the model is being forced toward results that, although corresponding to real ones, are not realistic in their deepest logic; nevertheless, to motivate this hypothesis it is sufficient to highlight two naturally observable constraints.

  %           \begin{enumerate}[label=V\arabic*]
  %             \item The first is of a physical nature induced by the spatial graph: by their nature, edges occupy spaces between cities and connect to them by taking a portion of their boundary; therefore, the only way a city can increase its degrees is to expand to have a larger boundary, but this necessarily implies an increase in population because the latter must maintain it.
  %             \item\label{rmkSDRSocialConstraint} The second is of a social nature induced by commuter flows: greater connectivity implies a more complex transport network, which can therefore be justified if there are enough commuters using it; these, however, can be attracted to a city if it has a sufficiently high population that naturally motivates such flows.
  %           \end{enumerate}

  %           The second constraint is indeed confirmed by \cref{figStrengthVsDegree20}, which reports a positive correlation between degree and strength: the greater the former, the greater the latter will be, i.e., the commuter flow, from which \ref{rmkSDRSocialConstraint} follows.

  %           Therefore, the physical constraint justifies this type of interaction regardless, because the graph is necessarily spatial; similarly, the social constraint motivates the interaction if and only if there is a positive correlation between degree and strength\footnotemark{}. All this implies the relation %\(k\gg1\implies s\gg1\)
  %           \(s\propto k\), also confirmed by the same real dispersion in \crefsubfigInt[ and ]{figSDERReferenceConfiguration}kl, which makes the introduction of \(k_{i_*}/k_i\) realistic from a modelling perspective.%
  %           \footnotetext{This characteristic is also confirmed for each Italian region, as well as for Italy itself.}%

  %           \tikzexternalizedfigure{t}
  %             {../../Figure/.fnl/Elaborato/Capitolo2/AnalisiRete20/SVsD.tex}
  %             {Strength vs degree of the Sardinia network.}
  %             {.75\linewidth}{figStrengthVsDegree20}

  %           It is also noteworthy to dwell on, albeit without a conclusive answer, why the emigration rule converges at all: one might wonder, in fact, why the flow 
  %           \(
              % \text{minor centers}\ \ \to\ \ \text{medium centers}\ \ \to\ \ \text{major centers}
  %             \text{minor centers}\to\text{medium centers}\to\text{major centers}
  %           \) 
  %           does not manifest itself.
  %           The only reason for the lack of degeneration in the model can be traced back to the presence of a backflow naturally induced by \ref{eqSizeDegreeEmigrationRule}, which balances the flow toward relatively more populated and connected centers; in other words, not only does the introduction of degrees in \ref{eqSizeEmigrationRule} make medium centers attractive compared to minor ones in the short term, but it also stabilizes long-term interactions, i.e., stationary ones. Yet the exact dynamics establishing a stabilizing reverse flow are unclear, since the ratio \(r\) in the Hill function (\cref{figDegreeAHillFunction}) should be smaller in \ref{eqSizeDegreeEmigrationRule} than in \ref{eqSizeEmigrationRule}.

  %           In conclusion, from these results, \ref{eqSizeDegreeEmigrationRule} implies that migration is probably governed by \textit{at least} two factors: the relative population and connectivity between two cities.

  %         \subsubsection{Size-strength rule - \ref{eqSizeStrengthEmigrationRule}}

  %           The introduction of strength is motivated by the same constraints described to motivate the introduction of degree in \ref{eqSizeEmigrationRule}. Reasoning this way, it is now possible to distinguish cities that have the same degree but not necessarily the same traffic; this natural distinction allows correcting the almost bijective relationship between degree and population in \crefsubfigInt[ and ]{figSDERReferenceConfiguration}{k}{l}, transforming it into a more realistic positive correlation described by \crefsubfigInt[ and ]{figSSERReferenceConfiguration}{k}{l}.

  %           Overall, given that these results are the best among those seen, \ref{eqSizeStrengthEmigrationRule} implies that \textit{most likely} the migration phenomenon is governed by two factors: the relative population and traffic between two cities.

  %         \subsubsection{\(\lambda\) and \(\alpha\) Interpretation}

  %           In all analyzed emigration rules there are two main parameters: \(\lambda\) and \(\alpha\); to try to understand their meaning, one can proceed through two parametric studies: both take the reference configuration of \ref{eqSizeStrengthEmigrationRule} in \cref{tabSSERParameters} and perturb only one of the two parameters leaving the other unchanged; \cref{tabRETFParametricStudies} summarizes what has been stated where \(N_p=5\) indicates the number of values uniformly sampled in the two intervals.

  %           \begin{table}[width=.6\linewidth,cols=5,pos=thb]
  %             \centering
  %             \caption{Parametric studies in the supplementary materials.}
  %             \begin{tabular*}{\tblwidth}{@{} CCCCC @{}}
  %               \toprule
  %               &\(N_p\)&\(\lambda\)&\(\alpha\)&\(N_t\)\\
  %               \midrule
  %               Parametric study \(\lambda\)&5&\([\num{.01},\num{.3}]\)&\(\num{.44}\)&\(\num{1e7}\)\\
  %               Parametric study \(\alpha\)&5&\(\num{.18}\)&\([\num{.3},\num{.6}]\)&\(\num{3e6}\)\\
  %               \bottomrule
  %             \end{tabular*}
  %             \label{tabRETFParametricStudies}
  %           \end{table}

  %           Before proceeding, it is necessary to keep in mind the difference between \(\mu\) and the lognormal average:

  %           \begin{rmk}
  %             \label{rmkLognormalExpectedValueConservation}
  %             By mass conservation (\cref{hypConservationAveragePopulation}), the sample mean%\footnote{%
                %For Sardinia, the value is \(\approx\num{4395.33}\) from \cref{tabSardiniaDataParameters}.%
              %}
  %             \(P/N\) is constant regardless of the chosen parametric pair \((\lambda,\alpha)\). However, this does \textit{not} imply that the parameter \(\mu\) in \(X\sim\Lognormal(\mu,\sigma)\) is also constant; indeed, what is preserved is its expected value of the form \( \mathbb E[X]=10^{\left(\mu+\sigma^2\ln(10)/2\right)}=P/N, \) from which it is immediately understood that if the variance \(\sigma^2\) were to increase, \(\mu\) would decrease and vice versa.

  %             As can be seen in \cref{figLognormalFunctionExample} this behaviour is hard to perceive both in linear and logarithmic scales, while in logarithmic space everything changes since \(\mu\) and \(\sigma\) correspond precisely to the mean and variance of the equivalent Gaussian. A similar reasoning also apply to the bilognormal fit \cref{eqBimodalLognormalDistribution}, as it is a convex combination of two lognormals.
  %           \end{rmk}

  %           Regarding the first parametric study, varying \(\lambda\), overall the following holds:

  %           \begin{itemize}
  %             \item the variance increases, broadening the distribution for increasing \(\lambda\), albeit slightly;
  %             \item the Pareto index de facto remains almost unaltered, varying between the extremes by only \(\approx\num{.05}\) on average;
  %             \item the relationship between Pareto and bilognormal remains unchanged: it confirms in any case that the second follows the end of the tail better than the first, while they coincide at the beginning;
  %             \item the scatter is unaltered as \(\lambda\) varies: the average behaviour is the same;
  %             \item all considered cases converge with shorter convergence time for increasing \(\lambda\).
  %           \end{itemize}

  %           It is noted that all these variations are nonetheless slight, considering that \(\lambda\) increases by an order of magnitude; this means that \(\lambda\) has little influence on the final results, except for the convergence time; However, for \(\lambda>\num{.3}\), an unstable configuration is still reached. Regarding the second parametric study, varying \(\alpha\), overall the following holds:

  %           \begin{itemize}
  %             \item as \(\lambda\) grows, the variance increases significantly, in so decreasing the mean and broadening the distribution;
  %             \item the Pareto index varies significantly, decreasing by nearly an order of magnitude;
  %             \item the relationship in the tail between Pareto and bilognormal continues to remain unchanged;
  %             \item the dispersion initially overestimates the minor-medium centres and underestimates the major ones, to then invert this relationship with increasing \(\alpha\);
  %             \item all considered cases converge, with a longer convergence time for increasing \(\alpha\).
  %           \end{itemize}

  %           Contrary to the \(\lambda\) parametric study, the variation of the results is significant; in particular, the Pareto index \(\beta\) seems to largely depend on the value of \(\alpha\); similarly, for values of \(\alpha>\num{.6}\), an unstable configuration is obtained. Altogether, in light of these studies, \(\lambda\) and \(\alpha\) can be interpreted as follows:

  %           \begin{itemize}
  %             \item \(\lambda\) is the attraction: the greater its value, the more the relatively larger centers succeed in attracting, in a single interaction, high fractions of the population of the relatively smaller centers;
  %             \item \(\alpha\) is the repulsion: the greater its value, the more the relatively larger centers are repelled by the relatively smaller ones; it is clearly inversely related to the magnitude of the stabilizing backflow.
  %           \end{itemize}

  %           Another point of view is the following: \(\lambda\) mainly influences the portion of the Hill function (\cref{figDegreeAHillFunction}) beyond unity (\(r>1\)), and thus the attraction, while \(\alpha\) conditions more the portion below unity (\(r\le1\)), and therefore the repulsion.

  %     \section{Conclusions}\label{secConclusions}

  %       All results illustrated in the previous section corroborate the node-agent perspective in an urban context to describe both the city size distribution and their Pareto index. The unit-rank approximation from \cref{defDegreeMatrix} produces results corresponding almost perfectly to those obtained exactly, confirming \cite{Nurisso2024}. All interaction rules faithfully reproduce the Pareto index of the actual distribution tail, while the distribution itself was overall well reproduced with almost every considered rule. In any case, the bimodal lognormal distribution fits the tail better than the Pareto distribution, confirming \cite{Gualandi2019SDC,Gualandi2019HBLD}; however, both fits present advantages and disadvantages: from an engineering perspective, the Pareto distribution is easier to handle having only one parameter, although it is simultaneously worse mathematically; conversely, the bimodal lognormal distribution is mathematically more correct but consists of \(5\) parameters, making any potential comparison more difficult.

  %       Of course, the work presented here is naturally non-exhaustive and leads to multiple future developments:

  %       \begin{enumerate}[label=\arabic*.]
  %         \item experiment with interaction rules containing other graph theory coefficients, such as betweenness centrality or the clustering coefficient, and by adding intra-urban commuters in \cref{eqSizeStrengthEmigrationRule};
  %         \item derive the Fokker-Planck equation of \ref{eqUrbanAsymmetricHomogeneousWeakFormBoltzmannType} to analytically study the stationary density, although the nonlinear and symmetric nature of the emigration rules proposed here makes this objective rather difficult to achieve;
  %         \item abandon the simplifying hypothesis of the static network to describe a dynamic one evolving in parallel with the population, with a possible comparison against real data [if available];
  %         \item try to reproduce the results with a European or international networks to confirm the validity of the current theory in other contexts or at larger scales.
  %       \end{enumerate}

  %     \appendix
      % \section*{Appendix}

  %     \section{Result reproducibility}

  %       \subsection{Technical specifications}

  %         The code is available in this repository \cite{Taralli2026} and was written entirely in Python 3.13.7, while the simulations were executed on a machine with a Windows operating system and an Intel Core i5-7500 processor (quad-core, 3.40 GHz).

  %       \subsection{Confidence interval}

  %         Let \(R\in\NPlus\) be the number of performed simulations and \(\mathbf S\in\RPlus^R\) be the random vector of a city's population related to each simulation; then by virtue of \cref{algMonteCarlo}, the components \((S_1,S_2,\ldots,S_R)\) of \(\mathbf S\) are independent and identically distributed according to the density \(\bar z_N\), a property that allows defining\nobreak
          %
  %         \begin{equation*}
  %           \bar S\equiv\frac1R\sum_{r=1}^RS_r
  %           \hspace{1em}\text{and}\hspace{1em}
  %           V\equiv\frac1{R-1}\sum_{r=1}^R(S_r-\bar S)^2,
  %         \end{equation*}
          %
  %         as the sample mean and sample variance respectively. Then, given \(\mu\) as the \textit{true}\footnote{%
  %           That is to say, \(\mu\) is \textit{not} a random variable but the exact mean parameter of \(\bar z_N\).%
  %         } mean of the density \(\bar z_N\), Student's \(T\)-distribution with parameter \(R-1\) is written as%
          %
  %         \begin{equation*}
  %           T=\frac{\bar S-\mu}{\sqrt{V/R}}\sim\Student(R-1).
  %         \end{equation*}
          %
  %         From this, one can define the [symmetric] confidence interval with confidence level \(\alpha\) by setting
          %
  %         \begin{equation*}
  %           \label{eqConfidenceIntervalCondition}
  %           \mathbb P\left(-t_{R-1}^{\alpha/2}\leq T\leq t_{R-1}^{\alpha/2}\right)=1-\alpha,
  %         \text{ where \(t_{R-1}^{\alpha/2}\in\mathbb R\) is such that }
  %           \mathbb P\left(T<t_{R-1}^{\alpha/2}\right)=1-\frac\alpha2.
  %         \end{equation*}
          %
  %         Expanding \(T\) and isolating the mean \(\mu\), we obtain
          %
  %         \begin{equation*}
  %           \mathbb P\left(
  %             \bar S-t_{R-1}^{\alpha/2}\sqrt{\frac VR}
  %             \leq\mu\leq
  %             \bar S+t_{R-1}^{\alpha/2}\sqrt{\frac VR}
  %           \right)=1-\alpha,
  %         \text{ whence }
  %           \IC^\alpha_R(\mathbf S)\equiv\left[
  %             \bar S-t_{R-1}^{\alpha/2}\sqrt{\frac VR},
  %             \bar S+t_{R-1}^{\alpha/2}\sqrt{\frac VR}
  %           \right]
  %         \end{equation*}
          %
  %         is the interval estimator defining the desired confidence interval, which can also be expressed more compactly as
          %
  %         \begin{equation}
  %           \label{eqConfidenceInterval}
  %           \IC^\alpha_R(\mathbf S)\equiv\bar S\pm t_{R-1}^{\alpha/2}\sqrt{V/R}.
  %         \end{equation}
          %
  %         The confidence interval is essentially the statistical error committed: it assesses how likely it is that the interval estimator \cref{eqConfidenceInterval} contains the parameter \(\mu\); in other words, \(\IC^\alpha_R\) measures the uncertainty of the mean estimation: given a sample \(\mathbf s\), the wider the confidence interval, the more uncertain the sample mean \(\bar s\) is as an estimation of the actual mean \(\mu\); conversely, the narrower it is, the more precise the estimation \(\bar s\) is, in the sense that \(\mu\) lies within a small neighbourhood of the sample mean. From this perspective, it is conceptually analogous to the precision of a measuring instrument.
          %
  %         \begin{hyp}[Confidence level]
  %           \label{hypConfidenceLevel}
  %           In all results \(\alpha=0.05\) is chosen to have an interval with a \(0.95\) confidence level.  
  %         \end{hyp}
          %
  %         \begin{rmk}
  %           Naturally, the interval estimator \cref{eqConfidenceInterval} just derived holds as much for the random vector \(\mathbf S\) as it does for its realizations\footnote{%
  %             In this case, \cref{eqConfidenceInterval} is referred to as an interval estimate.%
  %           } \(\mathbf s\), which are precisely the result of the \(R\) simulations.
  %         \end{rmk}%

  %       \subsection{Data structure}

  %         The data have the structure of a fourth-order tensor \(\mathbf s\in\RPlus^{N_f\times2\times N\times R}\) whose indices have the following meanings:
          %
  %         \begin{equation*}
  %           s^{n,a}_{i,r}
  %           \left\{\begin{aligned}
  %             n\ &=\text{ time instant},&\ 
  %             a\ &=\text{ simulation type},\\
  %             i\ &=\text{ city index},&\ 
  %             r\ &=\text{ simulation number}.
  %           \end{aligned}\right.
  %         \end{equation*}
          % 
  %         The time instant is defined via three parameters in \(\NPlus\):

  %         \begin{itemize}
  %           \item \(N_t\equiv\left\lfloor T/\Delta t\right\rfloor\) is the total number of simulated times;
  %           \item \(N_s\ll N_t\) is the number of snapshots from the simulated times;
  %           \item \(N_f<N_s\) is the number of reduced times from the snapshots.
  %         \end{itemize}

  %         Both the snapshots and the reductions are sampled at equally spaced intervals with steps \(\Delta s=N_t/N_s\) and \(\Delta f=N_s/N_f\) where it is assumed, for simplicity, that \(N_s\) and \(N_f\) are divisors of \(N_t\) and \(N_s\) respectively\footnote{%
  %           Namely \(N_t/N_s-\left\lfloor N_t/N_s\right\rfloor=0\) and \(N_s/N_f-\left\lfloor N_s/N_f\right\rfloor=0\).%
  %         },
  %         and the same is assumed between \(\Delta t\) and \(T\). By means of \cref{algMonteCarlo} a total of \(N_t\) times are simulated with step \(\Delta t\), of which \(N_s\) are saved in the snapshot tensor \(\undertilde{\mathbf s}\in\RPlus^{N_s\times2\times N\times R}\):
          %
  %         \begin{equation*}
  %           \undertilde s^{n,a}_{\cdot,r}\equiv\pazocal S^{n\Delta s,a}_{\cdot,r},
  %           \quad\forall n,\ \forall a,\ \forall r,
  %         \end{equation*}
          %
  %         where \(\pazocal S^{n\Delta s,a}_{\cdot,r}\) is to be understood as the vector of predicted sizes at the \(n\Delta s\)-th step, in the \(r\)-th exact simulation if \(a=1\), or approximated simulation if \(a=2\); subsequently, \(\undertilde{\mathbf s}\) is convolved with respect to the time index:
          %
  %         \begin{equation*}
  %           s^{n,a}_{i,r}\equiv\frac1{\Delta f}\sum_{j=1}^{\Delta f}
  %           \undertilde s^{(n-1)\Delta f+j,a}_{i,r},
  %           \quad\forall n,\ \forall a,\ \forall i,\ \forall r,
  %         \end{equation*}
          %
  %         that is to say, every \(\Delta f\) elements of the \(N_s\) snapshots are averaged. This mollification is necessary to make \(\mathbf s\) less noisy than \(\undertilde{\mathbf s}\) and therefore more readable once plotted.

  %         Overall, for the actual times, three forms are available depending on how the index \(n\) is interpreted:
          %
  %         \begin{equation*}
  %           t_n=n\Delta t,\;\;\forall n\in\{1,2,\ldots,N_t\},\hfill
  %           t^s_n=t_n\Delta s,\;\;\forall n\in\{1,2,\ldots,N_s\},\hfill
  %           t^f_n=t^s_n\Delta f,\;\;\forall n\in\{1,2,\ldots,N_f\},
  %         \end{equation*}
          %
  %         respectively for the discretized, sampling, and reduced times; regardless, \(t_{N_t}=t^s_{N_s}=t^f_{N_f}=T\) holds in any case.

  %         \begin{rmk}
  %           The index \(n\) in \(\mathbf s\) does not include \(0\) because it refers to the initial distribution, which is always equal to \(\pazocal S^0\) from \cref{algMonteCarlo}; if one wishes to extend the tensor, however, it can be written as \(\mathbf s^{0,a}_{\cdot,r}\equiv\pazocal S^0\), \(\forall a\), \(\forall r\).
  %         \end{rmk}

  %         \begin{hyp}
  %           All results consider \(\Delta t=0.01\), \(N_s=1000\), \(N_f=50\), while \(N_t\) is chosen to have \(t_n\) slightly higher than the convergence time, if the results converge, but it is certainly a multiple of \(10\) to ensure that \(N_s\) is one of its divisors.
  %         \end{hyp}

  %         \begin{rmk}
  %           For most plots, only the distribution at the final time \(T\) is considered \textit{without} time convolution, so that, for lightness of notation, the third-order tensor \(\mathbf s^T\equiv\undertilde{\mathbf s}^{N_s}\) can be improperly denoted by \(\mathbf s^T\in\RPlus^{2\times N\times R}\).
  %         \end{rmk}

  %       \subsection{Plot definitions}\label{apxPlotDefinitions}

  %         \subsubsection{Histograms}

  %           The first plots consider the [normalized] histograms of the distribution at the final time \(T\). We begin by considering the maximum and minimum elements of the tensor \(\mathbf s^T\),
            %
  %           \begin{equation*}
  %             s_{\max}\equiv\max_{a,i,r}\ s^{T,a}_{i,r}
  %             \quad\text{and}\quad
  %             s_{\min}\equiv\min_{a,i,r}\ s^{T,a}_{i,r},
  %           \end{equation*}
            %
  %           with which a common equally spaced grid can be defined to construct the histogram classes. Let \(N_c\) be the number of bins/classes of the grid; then, defining \(\Histogram\colon\mathbb R^N\to\mathbb R^{N_c}\) as the function that returns the [normalized] values of the histogram classes, the tensor \(\mathbf b\in\RPlus^{2\times N_c\times R}\) is written as \(\mathbf b^a_{\cdot,r}\equiv\Histogram(\mathbf s^{T,a}_{\cdot,r})\), \(\forall a\text{ and }\forall r\). Once the confidence intervals
            %
  %           \begin{equation*}
  %             \IC_R^\alpha(\mathbf b^a_c)=
  %             \bar b^a_c\pm t^{\alpha/2}_{R-1}\sqrt{v^a_c/R},
  %             \ \ \forall a\text{ and }\forall c,
  %           \end{equation*}
            %
  %           are evaluated, we can represent the mean histogram of the exact and approximated simulations with \(\bar{\mathbf b}^a\), and the stochastic error on the classes with \(\pm t^{\alpha/2}_{R-1}\sqrt{\mathbf v^a/R}\).

  %         \subsubsection{Bimodal Lognormal}

  %           The logic is similar to that of the histograms: let \(\mathcal L(\cdot;\mathbf s^{T,a}_{\cdot,r})\colon\RPlus\to\RPlus\), be the density of a bilognormal distribution \cref{eqBimodalLognormalDistribution} fitted from the vector \(\mathbf s^{T,a}_{\cdot,r}\), and let \(\mathcal L(\cdot;\mathbf s^{T,a})\colon\RPlus\to\RPlus^R\) be the vector-valued function such that
            %
  %           \begin{equation*}
  %             \mathcal L(x;\mathbf s^{T,a})\equiv
  %             \begin{bmatrix}
  %               \mathcal L(x;\mathbf s^{T,a}_{\cdot,1})&
  %               \mathcal L(x;\mathbf s^{T,a}_{\cdot,2})&
  %               \cdots&
  %               \mathcal L(x;\mathbf s^{T,a}_{\cdot,R})
  %             \end{bmatrix}^\top,
  %           \end{equation*}
            %
  %           which essentially collects all the fits pointwise into a single vector. Then the confidence intervals take the form
            %
  %           \begin{equation*}
  %             \IC_R^{\alpha}(\mathcal L(x;\mathbf s^{T,a}))=
  %             \bar{\mathcal L}^a(x)\pm t^{\alpha/2}_{R-1}\sqrt{v^a(x)/R},
  %             \ \ \forall x\in[s_{\min},s_{\max}]\text{ and }\forall a,
  %           \end{equation*}
            %
  %           where, as before, \(\bar{\mathcal L}^a(x)\) is the mean function while \(\pm t^{\alpha/2}_{R-1}\sqrt{v^a(x)/R}\) is its stochastic uncertainty.

  %           \begin{rmk}
  %             Since \(x\) is continuous, the set of confidence intervals forms a confidence band for the functions.
  %           \end{rmk}

  %           \begin{rmk}
  %             \label{ossBiLognormalParameterFitting}
  %             The confidence interval is applied both to the entire fit of the bilognormal and to its five parameters \(\xi\), \(\mu_1\), \(\sigma_1\), \(\mu_2\), and \(\sigma_2\), of which only the form of the first is reported for brevity: letting \(\boldsymbol\xi\in\mathbb R^{2\times R}\) be the matrix of all fractions obtained in the simulations, then it holds that
  %             \(
  %               \IC_R^\alpha\left(\boldsymbol\xi\right)
  %               =\bar\xi^a\pm t^{\alpha/2}_{R-1}\sqrt{v^a/R}.
  %             \)
  %           \end{rmk}

  %         \subsubsection{Empirical Complementary Cumulative Distribution Function}\label{secEmpiricalCCDF}

  %           The empirical CCDF is defined in the tail of the population distribution, which however needs to be specified: in this case, the last quartile is chosen as the threshold, i.e., that value \(s_\ThreeQuarters\) such that \(\mathbb P(S\leq s_\ThreeQuarters)=3/4\), given any random variable \(S\); in practice, this corresponds to finding that element \(s_\ThreeQuarters\) of the tensor \(\mathbf s^{T,a}_{\cdot,r}\) such that
            %
  %           \begin{equation*}
  %             \left\vert\left\{
  %               s^{T,a}_{i,r}\geq s_\ThreeQuarters\ |\ i\in\{1,2,\ldots,N\}
  %             \right\}\right\vert=\left\lceil\frac N4\right\rceil,
  %             \ \ \forall a\text{ and }\forall r,
  %           \end{equation*}
            %
  %           which in words means that \textit{at least} \(1/4\) of all the values of \(\mathbf s^{T,a}_{\cdot,r}\) are greater than or equal to \(s_\ThreeQuarters\). Then the interval \([s_\ThreeQuarters,s_{\max}]\) characterizes the tail.

  %           Accordingly, the empirical CCDF, \(\CCDF_{\mathcal E}(\cdot;\mathbf s^{T,a}_{\cdot,r})\colon\RPlus\to\RPlus\), which approximates \(\mathbb P(S>x)\), is written as
            %
  %           \begin{equation}
  %             \label{eqComplementaryCumulativeDistributionFunction}
  %             \CCDF_{\mathcal E}(x;\mathbf s^{T,a}_{\cdot,r})\equiv
  %             1-\frac1N\sum_{i=1}^N\chi_{[0,x]}
  %             \left(s^{T,a}_{i,r}\right),
  %             \ \ \forall a\text{ and }\forall r,
  %             \quad\text{where }
  %             \newcommand\hsep{\hspace{0.5em}}
  %             \chi_{[0,x]}(w)\equiv
  %             \left\{\begin{aligned}
  %               1&\hsep\text{if }w\in[0,x],\\
  %               0&\hsep\text{otherwise}.
  %             \end{aligned}\right.
  %           \end{equation}
            %
  %           is the indicator function. However, since \cref{eqComplementaryCumulativeDistributionFunction} must be evaluated for \(x\in[s_\ThreeQuarters,s_{\max}]\), and specifically on the tensor \(\mathbf p\in\mathbb R^{2\times N_\OneQuarter\times R}\) of \(\mathbf s\) restricted to the tail \(p^a_{i,r}\equiv s^{T,a}_{N-N_\OneQuarter+i,r}\), \(\forall a\), \(\forall i\in\{1,2,\ldots,N_\OneQuarter\}\) and \(\forall r\) with \(N_\OneQuarter\equiv\left\lceil N/4\right\rceil\), the following rescaling holds:
            %
  %           \begin{equation*}
  %             \begin{aligned}
  %               \CCDF_{\mathcal E}(x;\mathbf s^{T,a}_{\cdot,r})&=1-\frac{N-N_\OneQuarter}N-\frac1N\sum_{i=1}^{N_\OneQuarter}
  %               \chi_{[0,x]}\left(s^{T,a}_{N-N_\OneQuarter+i,r}\right)\\&
  %               =\frac{N_\OneQuarter}N\left[
  %                 1-\frac1{N_\OneQuarter}\smash{\sum_{i=1}^{N_\OneQuarter}}
  %                 \chi_{[0,x]}\left(p^{T,a}_{i,r}\right)
  %               \right]%\\&
  %               =\frac{N_\OneQuarter}N\CCDF_{\mathcal E}(x;\mathbf p^a_{\cdot,r}),
  %               \ \ \forall a\text{ and }\forall r.
  %             \end{aligned}
  %           \end{equation*}
            %
  %           That is to say, \(\CCDF_{\mathcal E}(x;\mathbf p^a_{\cdot,r})\) is equal to \(\CCDF_{\mathcal E}(x;\mathbf s^{T,a}_{\cdot,r})\) up to a scaling factor \(N_\OneQuarter/N\). At any rate, the Pareto index is invariant with respect to rescalings:

  %           \begin{rmk}
  %             \label{ossLogarithmParetoCCDF}
  %             In logarithmic scale, \cref{eqComplementaryCumulativePareto} is written as
  %             \(
  %               \log(\pazocal R(s))\approx-\beta\log(s)+\log(c),
  %             \)
  %             which is equivalent to a straight line with slope \(\beta\) and intercept \(\log(c)\); it follows that \( \log(\pazocal R'(s))=\log(\pazocal R(s))-\log(c)\approx-\beta\log(s)\), with \(\pazocal R'(s)\equiv\pazocal R(s)/c\), for which the Pareto index \(\beta\) remains unchanged.
  %           \end{rmk}

  %           From this observation, one can therefore choose \(\CCDF_{\mathcal E}(x;\mathbf s^{T,a}_{\cdot,r})\) in logarithmic scale for greater readability, precisely because what is of greatest interest is not the values themselves but rather their trend encapsulated in the parameter \(\beta\). The use of this scale, however, introduces a minor issue: the last element \(p^a_{N_\OneQuarter,r}\) has a zero ordinate and is thus not visible; to obviate the problem, the empirical CCDF can be translated in linear scale:
            %
  %           \begin{equation*}
  %             \CCDF_{\mathcal E}(x;\mathbf p^a_{\cdot,r})+\frac1{2N_\OneQuarter},
  %             \ \ \forall a\text{ and }\forall r,
  %           \end{equation*}
            %
  %           which corresponds in logarithmic scale to a non-uniform translation of the values, which are therefore distorted, albeit slightly since \(1/2N_\OneQuarter\) is a small quantity; for this reason, this modification is valid \textit{only graphically} and not for the computation of the Pareto index, which is always retrieved from the tensor \(\mathbf p\).

  %           Consequently, as a whole, the empirical CCDF is depicted in logarithmic scale by the scatter plot of the pairs
            %
  %           \begin{equation*}
  %             \left(
  %               \frac N{N_\OneQuarter}\CCDF_{\mathcal E}(\bar p^a_i;\mathbf s^{T,a}_{\cdot,r})
  %               +\frac1{2N_\OneQuarter},\bar p^a_i
  %             \right)
  %             =\left(
  %               \CCDF_{\mathcal E}(\bar p^a_i;\mathbf p^a_{\cdot,r})
  %               +\frac1{2N_\OneQuarter},\bar p^a_i
  %             \right),
  %             \ \ \forall a\text{ and }\forall i,
  %           \end{equation*}
            %
  %           where \(\bar p^a_i\) comes from the confidence intervals
            %
  %           \begingroup
  %             \setlength\belowdisplayskip{0pt}
  %             \setlength\belowdisplayshortskip{0pt}
  %             \begin{equation*}
  %               \IC_R^{\alpha}(\mathbf p^a_i)=
  %               \bar p^a_i\pm t^{\alpha/2}_{R-1}\sqrt{v^a_i/R},
  %               \ \ \forall a\text{ and }\forall i.
  %             \end{equation*}
  %           \endgroup
            %
  %           \begin{rmk}
  %             Note that, contrary to the previous cases, the confidence interval is horizontal and not vertical; indeed, the image of the empirical CCDF is invariant with respect to the simulations, which only modify the populations of the largest hubs, i.e., the abscissa of the empirical CCDF.
  %           \end{rmk}

  %         \subsubsection{Pareto CCDF and Relative Index}

  %           Regarding the fitting of the Pareto distribution, the reasoning is analogous to that seen for the bilognormal: let \(\CCDF_{\mathcal P}(\cdot;\mathbf p^a_{\cdot,r})\colon\RPlus\to\RPlus\), \(\forall a\text{ and }\forall r\), be the CCDF of the Pareto distribution density, analogous to \cref{eqComplementaryCumulativePareto}, fitted from the vector \(\mathbf p^a_{\cdot,r}\), and let \(\CCDF_{\mathcal P}(\cdot;\mathbf p^a)\colon\RPlus\to\RPlus\) be the vector-valued function such that
            %
  %           \begin{equation*}
  %             \CCDF_{\mathcal P}(x;\mathbf p^a)\equiv
  %             \begin{bmatrix}
  %               \CCDF_{\mathcal P}(x;\mathbf p^a_{\cdot,1})&
  %               \CCDF_{\mathcal P}(x;\mathbf p^a_{\cdot,2})&
  %               \cdots&
  %               \CCDF_{\mathcal P}(x;\mathbf p^a_{\cdot,R})
  %             \end{bmatrix}^\top;
  %           \end{equation*}
            %
  %           then the confidence intervals take the form
            %
  %           \begin{equation*}
  %             \IC_R^{\alpha}\left(\frac N{N_\OneQuarter}\CCDF_{\mathcal P}(x;\mathbf p^a)\right)
  %             =\bar{\CCDF}_{\mathcal P}^a(x)\pm t^{\alpha/2}_{R-1}\sqrt{\frac{v^a(x)}R},
  %             \ \ \forall x\in[s_\ThreeQuarters,s_{\max}]\text{ and }\forall a,
  %           \end{equation*}
            %
  %           where \(N/N_\OneQuarter\) is the rescaling factor described above and is necessary to compare the fit with the empirical CCDF.

  %           \begin{rmk}
  %             As in \cref{ossBiLognormalParameterFitting}, the confidence interval is applied both to the Pareto CCDF and to its parameter \(\beta\): letting \(\boldsymbol\beta\in\mathbb R^{2\times R}\) be the matrix of all Pareto indices obtained in the simulations, it then holds that
  %             \[
  %               \IC_R^\alpha\left(\boldsymbol\beta\right)
  %               =\bar\beta^a\pm t^{\alpha/2}_{R-1}\sqrt{v^a/R},
  %               \ \ \forall a.
  %             \]
  %           \end{rmk}

  %         \subsubsection{Bimodal Lognormal CCDF}

  %           The bilognormal is necessary not only for the entire distribution but also for the tail: let
            %
  %           \begin{equation*}
  %             \CCDF_{\mathcal L}(\cdot;\mathbf p^a_{\cdot,r}),\ 
  %             \CCDF_{\mathcal B}(\cdot;\mathbf p^a_{\cdot,r})
  %             \colon\RPlus\to\RPlus,
  %             \ \ \forall a\text{ and }\forall r,
  %           \end{equation*}
            %
  %           be the CCDFs of the lognormal and bilognormal distribution densities respectively, fitted to the vector \(\mathbf p^a_{\cdot,r}\); then by linearity of the CCDF, induced by the linearity of the integral, from \cref{eqBimodalLognormalDistribution} it holds that
            %
  %           \begin{equation*}
  %             \CCDF_{\mathcal B}(x;\mathbf p^a_{\cdot,r})
  %             =\xi\CCDF_{\mathcal L}(x;\mathbf p^a_{\cdot,r})
  %             +(1-\xi)\CCDF_{\mathcal L}(x;\mathbf p^a_{\cdot,r}),
  %             \ \ \forall a\text{ and }\forall r.
  %           \end{equation*}
            %
  %           Therefore, as before, we can define \(\CCDF_{\mathcal B}(\cdot;\mathbf p^a)\colon\RPlus\to\RPlus\) as the vector-valued function such that
            %
  %           \begin{equation*}
  %             \CCDF_{\mathcal B}(x;\mathbf p^a)\equiv
  %             \begin{bmatrix}
  %               \CCDF_{\mathcal B}(x;\mathbf p^a_{\cdot,1})&
  %               \CCDF_{\mathcal B}(x;\mathbf p^a_{\cdot,2})&
  %               \cdots&
  %               \CCDF_{\mathcal B}(x;\mathbf p^a_{\cdot,R})
  %             \end{bmatrix}^\top,
  %           \end{equation*}
            %
  %           which leads to the confidence intervals
            %
  %           \begin{equation*}
  %             \IC_R^{\alpha}\left(\frac N{N_\OneQuarter}\CCDF_{\mathcal B}(x;\mathbf p^a)\right)
  %             =\bar{\CCDF}_{\mathcal B}^a(x)\pm t^{\alpha/2}_{R-1}\sqrt{\frac{v^a(x)}R},
  %             \ \ \forall x\in[s_\ThreeQuarters,s_{\max}]\text{ and }\forall a.
  %           \end{equation*}

  %         \subsubsection{Data Transformation}

  %           As shown in \cref{figLognormalFunctionExample}(c), if some data \(\mathbf s\) are distributed according to a lognormal, then their reformulation \(\mathbf h\in\RPlus^{N_f\times2\times N\times R}\) in logarithmic space \(h_{i,r}^{n,a}\equiv\log(s_{i,r}^{n,a})\), \(\forall n\), \(\forall a\), \(\forall i\), \(\forall r\), is normally distributed; the advantage of using \(\mathbf h\) instead of \(\mathbf s\) lies in the easier interpretability of the data, especially when parameters vary, and thus in parametric studies like those shown in the supplementary material.

  %           This transformation is therefore applied to the plots defined so far, which change in the following way: first, the histograms are computed as before but with the reformulated data, so that the class values becomes \(\mathbf b^a_{\cdot,r}\equiv\Histogram(\mathbf h^{T,a}_{\cdot,r})\), \(\forall a\text{ and }\forall r\), and similarly for the confidence intervals. It is equally important to point out that in the space \(h\equiv\log(s)\), the density \(\bar f_N(s)\equiv\bar f_N(s,T)\) also transforms\footnotemark{} while preserving the total probability; this translates mathematically into
  %           \footnotetext{More specifically, as already mentioned at the beginning, it tends to a Gaussian density.}%
  %           \[
  %             \de P=\bar f_N(s)\de s
  %             =\bar l_N(h)\de h
  %             =\bar l_N(h)\frac{\de s}{s\ln(10)}
  %             \implies
  %             \bar l_N(h)=10^h\ln(10)\bar f_N(10^h),
  %           \]
  %           where \(\de P\) is the probability variation relative to the variations \(\de s / \de h\) at generic points \(s/h\), while \(\bar l_N(h)\) is the transformation of \(\bar f_N(s)\) into the new space \(h\).

  %           On the other hand, the empirical CCDFs, both Pareto or bilognormal, remain unchanged and are simply remapped from the interval \([s_{\min},s_{\max}]\) to \([h_{\min},h_{\max}]\), with \(h_{\min}=\log(s_{\min})\) and \(h_{\max}=\log(s_{\max})\). This follows from the fact that the logarithm is a strictly increasing function, thereby preserving the order of the transformed data:
            %
  %           \begin{equation*}
  %             \mathbb P(X<x)=\mathbb P(\log(X)<\log(x)),\ \ \forall x\in\RPlus,
  %           \end{equation*}
            %
  %           given a generic random variable \(X\in\RPlus\), it holds that:

  %           \begin{rmk}
  %             From \cref{ossLogarithmParetoCCDF}, transforming \(s\) into \(h\), the straight line becomes \(\log(\pazocal R'(h))\approx-\beta h\), which is the same relation but in semilogarithmic scale, hence this is the scale used in all previous CCDF plots.
  %           \end{rmk}

  %         \subsubsection{Evolution of the Mean Size}

  %           Expressing the tensor of the mean sizes with respect to the nodes as \(\bar{\mathbf s}\in\RPlus^{N_f\times2\times R}\)
            %
  %           \begin{equation*}
  %             \bar s^{n,a}_r\equiv
  %             \frac1N\sum_{i=1}^Ns^{n,a}_{i,r}
  %             ,\ \ \forall n,\ \forall a\text{ and }\forall r,
  %           \end{equation*}
            %
  %           where the confidence intervals of \(\bar{\mathbf s}\) are
            %
  %           \begin{equation*}
  %             \IC_R^{\alpha}(\bar{\mathbf s}^{n,a})=
  %             \bar{\bar s}^{n,a}\pm t^{\alpha/2}_{R-1}\sqrt{v^{n,a}/R},
  %             \ \ \forall n\text{ and }\forall a,
  %           \end{equation*}
            %
  %           then the trend of the total mean size is depicted by linearly interpolating the pairs \(\bigl( t^f_n,\bar{\bar s}^{n,a}\bigr)\), \(\forall n\) and \(\forall a\),
  %           at contiguous time instants with a confidence band of \(\pm t^{\alpha/2}_{R-1}\sqrt{v^{n,a}/R}\).

  %         \subsubsection{Mean Sizes versus Degrees}

  %           Let us now denote by \(\bar{\mathbf s}\in\RPlus^{2\times R}\) the tensor of the mean sizes with respect to the simulations
            %
  %           \begin{equation*}
  %             \bar s^a_i\equiv
  %             \frac1R\sum_{r=1}^Rs^{T,a}_{i,r}
  %             ,\ \ \forall a\text{ and }\forall i;
  %           \end{equation*}
            %
  %           thus, the pairs \(\left(k_i,\bar s^a_i\right)\), \(\forall a\) and \(\forall i\), defines the scatter plot linking the mean size of a node across all simulations to its degree.

  %           \begin{rmk}
  %             \label{ossLackOfDensityCI}
  %             In this case, confidence intervals are not drawn because the plot is already very dense and adding further vertical/horizontal bars would make it difficult to read.
  %           \end{rmk}

  %         \subsubsection{Evolutions of the Mean Sizes of Degree Classes}

  %           Let us denote by \(\bar s\in\mathbb R^{N_f\times2\times\vert\pazocal K\vert}\) the tensor
            %
  %           \begin{equation*}
  %             \bar s^{n,a}_k\equiv
  %             \frac1{\vert\mathcal I_k\vert}\sum_{i\in\mathcal I_k}
  %             \frac1R\sum_{r=1}^Rs^{n,a}_{i,r},
  %             \ \ \forall n,\ \forall a\text{ and }\forall k,
  %           \end{equation*}
            %
  %           where \(\pazocal K\subset\{1,2,\ldots,N\}\) is the set of distinct degrees of \(\mathbf k\), while \(\mathcal I_k\equiv\{i\in\mathcal I\ \vert\ k_i=k\}\subseteq\mathcal I\) represents the set of indices with degree \(k\); therefore, by linearly interpolating the pairs \(\bigl(t^f_n,\bar s^{n,a}_k\bigr)\), \(\forall n\),\(\ \forall a\) and \(\forall k\), at contiguous time instants, one can represent the time evolution of the mean population of the \(k\)-th class across all simulations. This plot is then used to ensure that the system is no longer in the transient phase and has therefore reached a steady state.

  %           We conclude by noting that confidence intervals are not drawn for the same reason clarified in \cref{ossLackOfDensityCI}.

      % To print the credit authorship contribution details
  %     \printcredits

      %% Loading bibliography style file
      %\bibliographystyle{model1-num-names}
  %     \bibliographystyle{../Corpo/cas-model2-names}

      % Loading bibliography database
  %     \bibliography{../Corpo/cas-refs}
    %\endgroup

    %\begingroup Studi Parametrici
      % \shortauthors{V. Taralli, A. Tosin}

      % %\begingroup λ
      %   % \shorttitle{Supplementary material - Parametric study \(\lambda\)}

      %   % \tikzexternalizedfigure{p}
      %   %   {../../Figure/.fnl/Elaborato/Capitolo4/20RETFStudioParametricol/SDF1/Infopagina.tex}
      %   %   {Distributions of the parametric study \(\lambda\).}
      %   %   {\linewidth}{figStudioParametricolRETFSDF1}

      %   % \tikzexternalizedfigure{p}
      %   %   {../../Figure/.fnl/Elaborato/Capitolo4/20RETFStudioParametricol/SDF2/Infopagina.tex}
      %   %   {Pareto fits of the parametric study \(\lambda\).}
      %   %   {\linewidth}{figStudioParametricolRETFSDF2}

      %   % \tikzexternalizedfigure{p}
      %   %   {../../Figure/.fnl/Elaborato/Capitolo4/20RETFStudioParametricol/SDF3/Infopagina.tex}
      %   %   {Pareto and bimodal lognormal fits of the parametric study \(\lambda\).}
      %   %   {\linewidth}{figStudioParametricolRETFSDF3}

      %   % \tikzexternalizedfigure{p}
      %   %   {../../Figure/.fnl/Elaborato/Capitolo4/20RETFStudioParametricol/SVsD/Infopagina.tex}
      %   %   {Sizes versus degrees of the parametric study \(\lambda\).}
      %   %   {\linewidth}{figStudioParametricolRETFSVsD}

      %   % \tikzexternalizedfigure{p}
      %   %   {../../Figure/.fnl/Elaborato/Capitolo4/20RETFStudioParametricol/SE/Infopagina.tex}
      %   %   {Evolutions of the average sizes of the degree classes of the parametric study \(\lambda\).}
      %   %   {\linewidth}{figStudioParametricolRETFSE}
      % %\endgroup

      % %\begingroup α
      %   % \shorttitle{Supplementary material - Parametric study \(\alpha\)}

      %   % \tikzexternalizedfigure{p}
      %   %   {../../Figure/.fnl/Elaborato/Capitolo4/20RETFStudioParametricoa/SDF1/Infopagina.tex}
      %   %   {Distributions of the parametric study \(\alpha\).}
      %   %   {\linewidth}{figStudioParametricoaRETFSDF1}

      %   % \tikzexternalizedfigure{p}
      %   %   {../../Figure/.fnl/Elaborato/Capitolo4/20RETFStudioParametricoa/SDF2/Infopagina.tex}
      %   %   {Pareto fits of the parametric study \(\alpha\).}
      %   %   {\linewidth}{figStudioParametricoaRETFSDF2}

      %   % \tikzexternalizedfigure{p}
      %   %   {../../Figure/.fnl/Elaborato/Capitolo4/20RETFStudioParametricoa/SDF3/Infopagina.tex}
      %   %   {Pareto and bimodal lognormal fits of the parametric study \(\alpha\).}
      %   %   {\linewidth}{figStudioParametricoaRETFSDF3}

      %   % \tikzexternalizedfigure{p}
      %   %   {../../Figure/.fnl/Elaborato/Capitolo4/20RETFStudioParametricoa/SVsD/Infopagina.tex}
      %   %   {Sizes versus degrees of the parametric study \(\alpha\).}
      %   %   {\linewidth}{figStudioParametricoaRETFSVsD}

      %   % \tikzexternalizedfigure{p}
      %   %   {../../Figure/.fnl/Elaborato/Capitolo4/20RETFStudioParametricoa/SE/Infopagina.tex}
      %   %   {Evolutions of the average sizes of the degree classes of the parametric study \(\alpha\).}
      %   %   {\linewidth}{figStudioParametricoaRETFSE}
      % %\endgroup
    %\endgroup
  % \end{document}
%#endregion Prima bozza
